% Transformations affines du plan — Mathématiques 1re Tech — Cours signé OnBuch+
% Programme officiel MINESEC. Compiler avec : tectonic N15-transformations-affines-plan.tex
\newcommand{\DOCMATIERE}{Mathématiques}
\newcommand{\DOCNIVEAU}{1re Tech}
\newcommand{\DOCMODULE}{Mathématiques – classe de Première technique}
\newcommand{\DOCLECON}{N15}
\newcommand{\DOCTITRE}{Transformations affines du plan}
% OnBuch+ — préambule « cours signés » ENRICHI (compilé avec tectonic / XeLaTeX)
% Système visuel « Pop » d'OnBuch+ : fond crème (jamais blanc), bordures encre
% épaisses, ombres « dures » décalées, titres Archivo Black en capitales.
% Version MATHÉMATIQUES : graphes (pgfplots/tikz), boîtes pédagogiques
% (définition, propriété, méthode, attention, le savais-tu, exercice résolu,
% exercice, TP, bilan), page de garde et signature OnBuch+ ancrée en bas.
%
% Macros attendues dans main.tex AVANT \input{preamble} :
%   \DOCMATIERE  \DOCNIVEAU  \DOCMODULE  \DOCLECON (numéro)  \DOCTITRE
\documentclass[11pt]{article}

\usepackage{fontspec}
\usepackage[french]{babel}
\usepackage[a4paper,margin=2cm,top=3.4cm,bottom=2.8cm,headheight=1.4cm,footskip=1.3cm]{geometry}
\usepackage{amsmath,amssymb,amsfonts}
\usepackage{mathtools} % \xmapsto, \coloneqq... souvent utilisés par les modèles
\usepackage{cancel} % simplifications barrées
% \abs{z} (module, valeur absolue) et \norm{u} (norme), utilisés par les modèles
\providecommand{\abs}{}\let\abs\relax\DeclarePairedDelimiter{\abs}{\lvert}{\rvert}
\providecommand{\norm}{}\let\norm\relax\DeclarePairedDelimiter{\norm}{\lVert}{\rVert}
\usepackage[table]{xcolor} % [table] : fournit aussi \rowcolors (alternance de lignes)
\usepackage{pagecolor}
\usepackage{tikz}
\usetikzlibrary{arrows.meta,positioning,calc,shapes.geometric,shapes.misc,shapes.arrows,decorations.pathmorphing,decorations.pathreplacing,decorations.markings,patterns,shadows,mindmap,trees,fit,backgrounds,matrix,intersections,angles,quotes}
\usepackage{pgfplots}
\usepackage{pgf-pie} % diagrammes circulaires \pie (statistiques)
\pgfplotsset{compat=1.18}
\usepgfplotslibrary{fillbetween,groupplots} % aires entre courbes (calcul intégral)
\usepackage{enumitem}
\usepackage{fancyhdr}
% [most] charge toutes les bibliothèques tcolorbox (listings, minted, raster,
% documentation...) et gonfle inutilement la mémoire TeX sur un document long
% (~300 boîtes) : on ne charge que ce qui est réellement utilisé.
\usepackage[skins,breakable]{tcolorbox}
\usepackage{graphicx}
\usepackage{colortbl}
\usepackage{booktabs}
\usepackage{tabularx}
\usepackage{diagbox} % case d'en-tête diagonale des tableaux à double entrée
% Colonnes extensibles centrée / gauche / droite (C, L, R), souvent utilisées par les modèles
\newcolumntype{C}{>{\centering\arraybackslash}X}
\newcolumntype{L}{>{\raggedright\arraybackslash}X}
\newcolumntype{R}{>{\raggedleft\arraybackslash}X}
\usepackage{array}
\usepackage{multicol}
\usepackage{siunitx}
% parse-numbers=false : accepte aussi \SI{2,5 \times 10^{-3}}{...}
\sisetup{locale=FR,output-decimal-marker={,},group-minimum-digits=4,parse-numbers=false}
\DeclareSIUnit\ohm{\ensuremath{\symup{\Omega}}} % U+2126 absent de la police
\usepackage{qrcode}
% \degree : commande fréquemment utilisée par les modèles pour un angle ou une
% température en dehors de \SI{}{} (non fournie par les paquets chargés).
\providecommand{\degree}{^{\circ}}
% \qed (fin de démonstration), utilisé par les modèles sans amsthm
\providecommand{\qed}{\hfill$\square$}
% \barre : double barre des valeurs interdites dans un tableau de variations (tkz-tab)
\providecommand{\barre}{\ensuremath{\|}}
% environnement proof (amsthm non chargé) utilisé par les modèles
\ifdefined\proof\else\newenvironment{proof}[1][Démonstration]{\par\noindent\textit{#1.}\ }{\qed\par}\fi
\usepackage{lastpage}
\usepackage{hyperref}
\hypersetup{hidelinks}

% ============================================================
% Palette « Pop » OnBuch+ — mêmes hex que la classe P (plus_ui.dart)
% ============================================================
\definecolor{poporange}{HTML}{F59321}
\definecolor{popdark}{HTML}{E07A0C}
\definecolor{popcream}{HTML}{FFF6E8}
\definecolor{popink}{HTML}{1C1714}
\definecolor{poppurple}{HTML}{7A5AE0}
\definecolor{popgreen}{HTML}{2E9E6A}
\definecolor{poppink}{HTML}{E0457B}
\definecolor{popblue}{HTML}{2F7DE1}
\definecolor{popgold}{HTML}{C98A12}
\definecolor{popmuted}{HTML}{8A7F76}
\definecolor{popline}{HTML}{EADFD2}
\definecolor{popgray}{HTML}{8A7F76}
\colorlet{popgrey}{popgray}
% Alias tolérants (le modèle nomme parfois les couleurs différemment)
\colorlet{popwhite}{white}
\colorlet{popblack}{popink}
\colorlet{popred}{poppink}
\colorlet{popyellow}{popgold}
% Teintes claires pour fonds de tableaux / zones
\colorlet{popblueL}{popblue!10!white}
\colorlet{poporangeL}{poporange!14!white}
\colorlet{popgreenL}{popgreen!12!white}
\colorlet{poppurpleL}{poppurple!12!white}
\colorlet{poppinkL}{poppink!12!white}
\colorlet{popgoldL}{popgold!14!white}

\pagecolor{popcream}

% ============================================================
% Typographie
% ============================================================
\defaultfontfeatures{Ligatures=TeX}
\setmainfont{PlusJakartaSans-Medium.ttf}[Path=../../../fonts/,BoldFont=PlusJakartaSans-Bold.ttf]
% Maths en sans-serif (Fira Math) avec chiffres et lettres droites de Plus
% Jakarta, pour que les formules se fondent dans le texte.
\usepackage[math-style=ISO,bold-style=ISO]{unicode-math}
\setmathfont{FiraMath-Regular.otf}[Path=../../../fonts/]
\setmathfont{PlusJakartaSans-Medium.ttf}[Path=../../../fonts/,range={up/num,up/{Latin,latin},\mathup}]
\usepackage{icomma} % virgule décimale sans espace en mode math
% Caractères mathématiques Unicode absents des polices de texte (Plus Jakarta,
% Archivo Black) : rendus par leur équivalent mathématique, en texte comme en
% formule — sinon ils s'affichent en carré vide (ex. « Arithmétique dans ℤ »).
\usepackage{newunicodechar}
\newunicodechar{ℝ}{\ensuremath{\mathbb{R}}}
\newunicodechar{ℤ}{\ensuremath{\mathbb{Z}}}
\newunicodechar{ℂ}{\ensuremath{\mathbb{C}}}
\newunicodechar{ℕ}{\ensuremath{\mathbb{N}}}
\newunicodechar{ℚ}{\ensuremath{\mathbb{Q}}}
\newunicodechar{𝔻}{\ensuremath{\mathbb{D}}}
\newunicodechar{α}{\ensuremath{\alpha}}
\newunicodechar{β}{\ensuremath{\beta}}
\newunicodechar{γ}{\ensuremath{\gamma}}
\newunicodechar{δ}{\ensuremath{\delta}}
\newunicodechar{ε}{\ensuremath{\varepsilon}}
\newunicodechar{θ}{\ensuremath{\theta}}
\newunicodechar{λ}{\ensuremath{\lambda}}
\newunicodechar{μ}{\ensuremath{\mu}}
\newunicodechar{σ}{\ensuremath{\sigma}}
\newunicodechar{τ}{\ensuremath{\tau}}
\newunicodechar{φ}{\ensuremath{\varphi}}
\newunicodechar{ψ}{\ensuremath{\psi}}
\newunicodechar{ω}{\ensuremath{\omega}}
\newunicodechar{Σ}{\ensuremath{\Sigma}}
% majuscules produites par \MakeUppercase dans les titres (α-aminés -> Α)
\newunicodechar{Α}{\ensuremath{\alpha}}
\newunicodechar{Β}{\ensuremath{\beta}}
\newunicodechar{Γ}{\ensuremath{\Gamma}}
\newunicodechar{Δ}{\ensuremath{\Delta}}
\newunicodechar{Ε}{\ensuremath{\varepsilon}}
\newunicodechar{Θ}{\ensuremath{\Theta}}
\newunicodechar{Λ}{\ensuremath{\Lambda}}
\newunicodechar{Μ}{\ensuremath{\mu}}
\newunicodechar{Τ}{\ensuremath{\tau}}
\newunicodechar{Φ}{\ensuremath{\Phi}}
\newunicodechar{Ψ}{\ensuremath{\Psi}}
\newunicodechar{Ω}{\ensuremath{\Omega}}
\newunicodechar{↦}{\ensuremath{\mapsto}}
\newunicodechar{∈}{\ensuremath{\in}}
\newunicodechar{∉}{\ensuremath{\notin}}
\newunicodechar{∧}{\ensuremath{\wedge}}
\newunicodechar{⇔}{\ensuremath{\Leftrightarrow}}
\newunicodechar{⟺}{\ensuremath{\Longleftrightarrow}}
\newunicodechar{⇒}{\ensuremath{\Rightarrow}}
\newunicodechar{⟹}{\ensuremath{\Longrightarrow}}
\newunicodechar{∘}{\ensuremath{\circ}}
\newunicodechar{∪}{\ensuremath{\cup}}
\newunicodechar{∩}{\ensuremath{\cap}}
\newunicodechar{≡}{\ensuremath{\equiv}}
\newunicodechar{⊂}{\ensuremath{\subset}}
\newunicodechar{⊥}{\ensuremath{\perp}}
\newunicodechar{∃}{\ensuremath{\exists}}
\newunicodechar{∀}{\ensuremath{\forall}}
\newunicodechar{∛}{\ensuremath{\sqrt[3]{\,}}}
\newunicodechar{∜}{\ensuremath{\sqrt[4]{\,}}}
\newunicodechar{⃗}{\ensuremath{{}^{\to}}}
\newunicodechar{ⁿ}{\ifmmode^{n}\else\textsuperscript{\ensuremath{n}}\fi}
\newunicodechar{ˣ}{\ifmmode^{x}\else\textsuperscript{\ensuremath{x}}\fi}
\newunicodechar{ᵇ}{\ifmmode^{b}\else\textsuperscript{\ensuremath{b}}\fi}
\newunicodechar{ʳ}{\ifmmode^{r}\else\textsuperscript{\ensuremath{r}}\fi}
\newunicodechar{ᵘ}{\ifmmode^{u}\else\textsuperscript{\ensuremath{u}}\fi}
\newunicodechar{ʸ}{\ifmmode^{y}\else\textsuperscript{\ensuremath{y}}\fi}
\newunicodechar{ᵃ}{\ifmmode^{a}\else\textsuperscript{\ensuremath{a}}\fi}
\newunicodechar{ᵉ}{\ifmmode^{e}\else\textsuperscript{\ensuremath{e}}\fi}
\newunicodechar{ᵏ}{\ifmmode^{k}\else\textsuperscript{\ensuremath{k}}\fi}
\newunicodechar{ᵖ}{\ifmmode^{p}\else\textsuperscript{\ensuremath{p}}\fi}
\newunicodechar{ᴺ}{\ifmmode^{N}\else\textsuperscript{\ensuremath{N}}\fi}
\newunicodechar{ᶜ}{\ifmmode^{c}\else\textsuperscript{\ensuremath{c}}\fi}
\newunicodechar{ʰ}{\ifmmode^{h}\else\textsuperscript{\ensuremath{h}}\fi}
\newunicodechar{ᵠ}{\ifmmode^{q}\else\textsuperscript{\ensuremath{q}}\fi}
\newunicodechar{ₙ}{\ifmmode_{n}\else\textsubscript{\ensuremath{n}}\fi}
\newunicodechar{ₐ}{\ifmmode_{a}\else\textsubscript{\ensuremath{a}}\fi}
\newunicodechar{ᵢ}{\ifmmode_{i}\else\textsubscript{\ensuremath{i}}\fi}
\newunicodechar{ᵣ}{\ifmmode_{r}\else\textsubscript{\ensuremath{r}}\fi}
\newunicodechar{ₖ}{\ifmmode_{k}\else\textsubscript{\ensuremath{k}}\fi}
\newunicodechar{ᵦ}{\ifmmode_{\beta}\else\textsubscript{\ensuremath{\beta}}\fi}
\newunicodechar{ₓ}{\ifmmode_{x}\else\textsubscript{\ensuremath{x}}\fi}
\newfontfamily\popdisplay{ArchivoBlack-Regular.ttf}[Path=../../../fonts/]
\newfontfamily\poplogo{SpaceGrotesk-Bold.ttf}[Path=../../../fonts/]
\newfontfamily\popsemi{PlusJakartaSans-SemiBold.ttf}[Path=../../../fonts/]
\newfontfamily\popxbold{PlusJakartaSans-ExtraBold.ttf}[Path=../../../fonts/]
\newfontfamily\popxboldls{PlusJakartaSans-ExtraBold.ttf}[Path=../../../fonts/,LetterSpace=3.0]

\setlist[enumerate]{leftmargin=1.7em,itemsep=5pt,topsep=4pt}
\setlist[itemize]{leftmargin=1.7em,itemsep=4pt,topsep=4pt,label={\color{poporange}\textbullet}}
\setlength{\parindent}{0pt}
\setlength{\parskip}{5pt}
\renewcommand{\arraystretch}{1.25}

% pgfplots : style Pop par défaut
\pgfplotsset{
  every axis/.append style={axis line style={popink,line width=1pt},tick style={popink},
    grid style={popline,line width=0.5pt},label style={font=\small},tick label style={font=\footnotesize}},
  popcurve/.style={poporange,line width=1.8pt,no markers,smooth},
}
% Même style disponible dans un \draw TikZ simple (hors pgfplots)
\tikzset{popcurve/.style={poporange,line width=1.8pt}}

% ============================================================
% Pill / squiggle
% ============================================================
\newcommand{\poppill}[3]{%
  \tikz[baseline=-0.55ex]\node[draw=popink,line width=1.2pt,rounded corners=2.6pt,%
    fill=#2,inner xsep=6.5pt,inner ysep=3pt]{%
    \popxboldls\fontsize{7}{7}\selectfont\color{#3}\MakeUppercase{#1}};%
}
\newcommand{\popsquiggle}[1]{%
  \tikz[baseline=-0.4ex]{
    \draw[#1,line width=2.6pt,line cap=round]
      plot [smooth] coordinates {(0,0.09) (0.34,0.01) (0.68,0.21) (1.02,0.01) (1.36,0.10)};
  }%
}
% Mot-clé surligné (marqueur orange sous le texte)
\newcommand{\cle}[1]{{\popxbold\color{popdark}#1}}

% ============================================================
% En-tête / pied
% ============================================================
\pagestyle{fancy}
\fancyhf{}
\renewcommand{\headrulewidth}{0pt}
\renewcommand{\footrulewidth}{0pt}
\lhead{\raisebox{-0.15ex}{\poplogo\fontsize{13}{13}\selectfont\color{popink}OnBuch\color{poporange}+}}
\rhead{\poppill{\DOCMATIERE\ \textbullet\ \DOCNIVEAU\ \textbullet\ Leçon \DOCLECON}{white}{popink}}
\lfoot{\popxbold\fontsize{7.6}{7.6}\selectfont\color{popmuted}\MakeUppercase{Cours signé OnBuch+}\ \textbullet\ {\popsemi\DOCTITRE}}
\rfoot{\tikz[baseline=-0.6ex]\node[rounded corners=8.5pt,fill=popink,minimum height=17pt,minimum width=17pt,inner xsep=5pt,inner sep=0]{\popxbold\fontsize{8}{8}\selectfont\color{popcream}\thepage\,/\,\pageref*{LastPage}};}

% ============================================================
% Titres
% ============================================================
\newcommand{\chaptitre}[2]{%
  \begin{tcolorbox}[enhanced,colback=poporange,colframe=popink,boxrule=2.6pt,arc=11pt,
      drop shadow=popink,left=15pt,right=15pt,top=12pt,bottom=11pt,before skip=3pt,after skip=18pt]
    \poppill{#2}{popink}{popcream}\par\vspace{9pt}
    {\raggedright\hyphenpenalty=10000\exhyphenpenalty=10000\popdisplay\fontsize{22}{24}\selectfont\color{popcream}\MakeUppercase{#1}\par}
  \end{tcolorbox}
}
% Section numérotée : pastille orange avec le numéro + titre Archivo
\newcounter{popsec}
\newcommand{\coursec}[1]{%
  \stepcounter{popsec}\setcounter{popsub}{0}%
  \par\vspace{16pt}\noindent
  \begin{minipage}{\linewidth}
  \tikz[baseline=-0.7ex]\node[draw=popink,line width=1.6pt,fill=poporange,rounded corners=4pt,minimum size=22pt,inner sep=0]{\popdisplay\fontsize{12}{12}\selectfont\color{popcream}\thepopsec};%
  \hspace{8pt}{\popdisplay\fontsize{15}{17}\selectfont\color{popink}\MakeUppercase{#1}}\par
  \vspace{4pt}\noindent{\color{poporange}\rule{36pt}{3.6pt}}
  \end{minipage}\par\nopagebreak\vspace{8pt}
}
\newcounter{popsub}
\newcommand{\courssub}[1]{%
  \stepcounter{popsub}%
  \par\vspace{9pt}\noindent{\popxbold\fontsize{11.5}{13}\selectfont\color{popink}{\color{poporange}\thepopsec.\thepopsub}\hspace{6pt}#1}\par\nopagebreak\vspace{4pt}
}

% ============================================================
% Boîtes pédagogiques — même recette : fond blanc, bordure épaisse colorée,
% ombre dure encre, badge plein. Titre optionnel [#1].
% ============================================================
\tcbset{popbase/.style={breakable,enhanced,colback=white,boxrule=2.3pt,arc=9pt,
  shadow={3pt}{-3pt}{0pt}{popink},left=14pt,right=14pt,top=11pt,bottom=11pt,before skip=11pt,after skip=15pt,
  fontupper=\popsemi\fontsize{10.6}{14.5}\selectfont\color{popink},
  fontlower=\popsemi\fontsize{10.6}{14.5}\selectfont\color{popink}}}
% Coche dessinée (le glyphe ✓ n'existe pas dans les polices du document)
\newcommand{\popcheck}[1]{\tikz[baseline=-0.5ex]\draw[#1,line width=1.6pt,line cap=round,line join=round](-0.13,0.0)--(-0.04,-0.09)--(0.13,0.1);}
\newcommand{\popboxhead}[3]{\poppill{#1}{#2}{white}\ifx\relax#3\relax\else\hspace{6pt}{\popxbold\fontsize{10.6}{12}\selectfont\color{popink}#3}\fi\par\vspace{7pt}}

\newtcolorbox{aretenir}[1][]{popbase,colframe=popgreen,before upper={\popboxhead{À retenir}{popgreen}{#1}}}
\newtcolorbox{exemplebox}[1][]{popbase,colframe=poporange,before upper={\popboxhead{Exemple}{poporange}{#1}}}
\newtcolorbox{definition}[1][]{popbase,colframe=popblue,before upper={\popboxhead{Définition}{popblue}{#1}}}
\newtcolorbox{propriete}[1][]{popbase,colframe=poppurple,before upper={\popboxhead{Propriété}{poppurple}{#1}}}
\newtcolorbox{methode}[1][]{popbase,colframe=popgold,before upper={\popboxhead{Méthode}{popgold}{#1}}}
\newtcolorbox{attention}[1][]{popbase,colframe=poppink,before upper={\popboxhead{Attention}{poppink}{#1}}}
\newtcolorbox{experience}[1][]{popbase,colframe=popblue,colback=popblueL,before upper={\popboxhead{Expérience}{popblue}{#1}}}
% « Le savais-tu ? » : carte violette pleine (contexte, culture, Cameroun)
\newtcolorbox{savaistu}[1][]{popbase,colback=poppurple,colframe=popink,
  fontupper=\popsemi\fontsize{10.4}{14.2}\selectfont\color{white},
  before upper={\poppill{Le savais-tu\,?}{white}{poppurple}\ifx\relax#1\relax\else\hspace{6pt}{\popxbold\color{white}#1}\fi\par\vspace{7pt}}}
% Exercice résolu : énoncé en haut, \tcblower puis la solution sur fond crème
\newcounter{popexo}\newcounter{popexr}
\newtcolorbox{exoresolu}[1][]{popbase,colframe=poporange,colbacklower=popcream,
  segmentation style={popink,line width=1.2pt,dashed},
  before upper={\stepcounter{popexr}\popboxhead{Exercice résolu \thepopexr}{poporange}{#1}},
  before lower={\poppill{Solution}{popink}{popcream}\par\vspace{6pt}}}
% Exercice (non résolu en place) — la correction est dans la partie Corrigés
\newtcolorbox{exercice}[1][]{popbase,colframe=popink,
  before upper={\stepcounter{popexo}\popboxhead{Exercice \thepopexo}{popink}{#1}}}
% Corrigé
\newtcolorbox{corrige}[1][]{popbase,colframe=popgreen,colback=popgreenL,
  before upper={\popboxhead{Corrigé}{popgreen}{#1}}}
% Objectifs / prérequis / bilan
\newtcolorbox{objectifs}{popbase,colframe=popink,colback=poporangeL,
  before upper={\popboxhead{Objectifs de la leçon}{popink}{}}}
\newtcolorbox{prerequis}{popbase,colframe=popink,colback=popblueL,
  before upper={\popboxhead{Prérequis}{popblue}{}}}
\newtcolorbox{bilan}{popbase,colframe=popink,colback=white,boxrule=2.8pt,
  before upper={\popboxhead{Fiche bilan}{poporange}{L'essentiel à connaître pour l'examen}}}
% Figure encadrée avec légende
\newtcolorbox{popfigure}[1][]{enhanced,breakable=false,colback=white,colframe=popline,boxrule=1.4pt,arc=8pt,
  left=10pt,right=10pt,top=10pt,bottom=8pt,before skip=10pt,after skip=12pt,halign=center,
  lower separated=false,halign lower=center,
  fontlower=\popsemi\fontsize{9}{11.5}\selectfont\color{popmuted},
  #1}
\newcommand{\legende}[1]{\par\vspace{6pt}{\popsemi\fontsize{9}{11.5}\selectfont\color{popmuted}#1}}

% ============================================================
% Signature OnBuch+
% ============================================================
\newcommand{\poplogoinline}[1]{{\poplogo\fontsize{#1}{#1}\selectfont\color{popink}OnBuch\color{poporange}+}}
\newcommand{\signatureonbuch}{%
  \begin{tcolorbox}[enhanced,colback=popink,colframe=popink,boxrule=2.2pt,arc=10pt,
      drop shadow=poporange,left=14pt,right=14pt,top=10pt,bottom=10pt,before skip=0pt,after skip=0pt]
    \tikz[baseline=-0.7ex]\node[circle,fill=popgreen,draw=popcream,line width=1.3pt,minimum size=22pt,inner sep=0]{\popcheck{white}};%
    \hspace{9pt}\begin{minipage}[c]{0.62\linewidth}
      {\popxbold\fontsize{12}{13}\selectfont\color{popcream}Signé\ \textbullet\ {\poplogo OnBuch\color{poporange}+}}\par\vspace{2pt}
      {\raggedright\popsemi\fontsize{8}{10.5}\selectfont\color{popline}\MakeUppercase{Équipe pédagogique OnBuch+}\\\MakeUppercase{Conforme au programme officiel MINESEC}\par}
    \end{minipage}\hfill
    {\poplogo\fontsize{22}{22}\selectfont\color{popcream}OnBuch\color{poporange}+}
  \end{tcolorbox}%
}

% ============================================================
% Page de garde
% #1 titre · #2 matière · #3 niveau · #4 module · #5 n° leçon · #6 durée ·
% #7 description / objectifs (texte ou itemize)
% ============================================================
\newcommand{\pagegarde}[7]{%
  \thispagestyle{empty}
  \newgeometry{a4paper,margin=1.7cm,top=1.6cm,bottom=1.4cm}%
  \begin{tikzpicture}[remember picture,overlay]
    \fill[poporange!18] ([xshift=-3.2cm,yshift=-3.6cm]current page.north east) circle (4.6cm);
    \fill[poppurple!14] ([xshift=2.4cm,yshift=7.6cm]current page.south west) circle (3.8cm);
  \end{tikzpicture}%
  {\poplogo\fontsize{30}{30}\selectfont\color{popink}OnBuch\color{poporange}+}\hfill
  \raisebox{0.6ex}{\poppill{Cours signé \textbullet\ Premium}{popink}{popcream}}\par
  \vspace{1pt}\popsquiggle{poppurple}\par
  \vspace{8pt}
  \begin{tcolorbox}[enhanced,colback=poporange,colframe=popink,boxrule=2.8pt,arc=13pt,
      drop shadow=popink,left=18pt,right=18pt,top=12pt,bottom=13pt,after skip=10pt]
    \poppill{#2 \textbullet\ #3}{popink}{popcream}\hspace{5pt}\poppill{#4}{white}{popink}\par\vspace{8pt}
    {\popdisplay\fontsize{14}{14}\selectfont\color{popink}LEÇON #5}\par\vspace{4pt}
    {\raggedright\hyphenpenalty=10000\exhyphenpenalty=10000\popdisplay\fontsize{25}{27}\selectfont\color{popcream}\MakeUppercase{#1}\par}
    \vspace{8pt}
    {\popxbold\fontsize{9.5}{9.5}\selectfont\color{popink}\MakeUppercase{Durée conseillée : #6}}
  \end{tcolorbox}
  \begin{tcolorbox}[enhanced,colback=white,colframe=popink,boxrule=2.2pt,arc=10pt,
      drop shadow=popink,left=16pt,right=16pt,top=10pt,bottom=10pt,after skip=8pt]
    \poppill{Dans cette leçon}{poppurple}{white}\par\vspace{5pt}
    \popsemi\fontsize{9.3}{12.6}\selectfont\color{popink}#7
  \end{tcolorbox}
  \noindent\begin{minipage}[t]{0.487\linewidth}
    \begin{tcolorbox}[enhanced,colback=popgreen,colframe=popink,boxrule=2.2pt,arc=10pt,
        drop shadow=popink,left=11pt,right=11pt,top=10pt,bottom=10pt,height=3.1cm,valign=center]
      \tcbox[colback=white,colframe=popink,boxrule=1.3pt,arc=5pt,boxsep=0pt,nobeforeafter,
        left=3pt,right=3pt,top=3pt,bottom=3pt,valign=center]{\qrcode[height=1.9cm]{https://play.google.com/store/apps/details?id=com.luvvix.onbuch&hl=fr}}%
      \hspace{8pt}\begin{minipage}[c]{0.5\linewidth}
        {\popdisplay\fontsize{11}{12.5}\selectfont\color{white}\MakeUppercase{Télécharge OnBuch}}\par\vspace{4pt}
        {\raggedright\popsemi\fontsize{8.4}{11}\selectfont\color{white}Gratuit \textbullet\ Google Play\\Tuteur IA, annales, concours, résultats\par}
      \end{minipage}
    \end{tcolorbox}
  \end{minipage}\hfill
  \begin{minipage}[t]{0.487\linewidth}
    \begin{tcolorbox}[enhanced,colback=poppurple,colframe=popink,boxrule=2.2pt,arc=10pt,
        drop shadow=popink,left=11pt,right=11pt,top=10pt,bottom=10pt,height=3.1cm,valign=center]
      \tikz[baseline=-0.7ex]\node[circle,fill=white,draw=popink,line width=1.3pt,minimum size=1.6cm,inner sep=0]{\popdisplay\fontsize{24}{24}\selectfont\color{poppurple}+};%
      \hspace{8pt}\begin{minipage}[c]{0.6\linewidth}
        {\popdisplay\fontsize{11}{12.5}\selectfont\color{white}\MakeUppercase{Passe à OnBuch+}}\par\vspace{4pt}
        {\raggedright\popsemi\fontsize{8.4}{11}\selectfont\color{white}Tous les cours signés, Tuteur IA illimité, fiches PDF — onglet OnBuch+ de l'app\par}
      \end{minipage}
    \end{tcolorbox}
  \end{minipage}
  \vspace{8pt}
  \popcontenu
  \vfill
  \signatureonbuch
  \restoregeometry
  \newpage
}

% Rangée « Ce cours contient » (page de garde)
\newcommand{\poptile}[3]{% #1 couleur #2 gros texte #3 légende
  \begin{tcolorbox}[enhanced,colback=#1,colframe=popink,boxrule=2pt,arc=9pt,shadow={2.5pt}{-2.5pt}{0pt}{popink},
      left=6pt,right=6pt,top=9pt,bottom=9pt,halign=center,valign=center,height=1.75cm,width=\linewidth,nobeforeafter]
    {\popdisplay\fontsize{12.5}{13}\selectfont\color{white}\MakeUppercase{#2}}\par\vspace{3pt}
    {\centering\popsemi\fontsize{7.8}{9.5}\selectfont\color{white}#3\par}
  \end{tcolorbox}}
\newcommand{\popcontenu}{%
  \par\noindent{\popxboldls\fontsize{7.5}{7.5}\selectfont\color{popmuted}CE COURS SIGNÉ CONTIENT}\par\vspace{7pt}
  \noindent\begin{minipage}[t]{0.235\linewidth}\poptile{popblue}{Cours}{complet, illustré, pas à pas}\end{minipage}\hfill
  \begin{minipage}[t]{0.235\linewidth}\poptile{poporange}{Méthodes}{et exercices résolus}\end{minipage}\hfill
  \begin{minipage}[t]{0.235\linewidth}\poptile{poppink}{Exercices}{type examen + corrigés}\end{minipage}\hfill
  \begin{minipage}[t]{0.235\linewidth}\poptile{popgreen}{Fiche bilan}{l'essentiel à retenir}\end{minipage}\par}

% Sommaire
\newtcolorbox{sommaire}{enhanced,colback=white,colframe=popink,boxrule=2.2pt,arc=10pt,
  drop shadow=popink,left=18pt,right=18pt,top=13pt,bottom=13pt,before skip=6pt,after skip=20pt,
  before upper={\poppill{Sommaire}{poppurple}{white}\par\vspace{9pt}\popsemi\fontsize{10.6}{14.5}\selectfont\color{popink}}}

% Signature de fin de cours — poussée tout en bas de la dernière page
\newcommand{\findecours}{%
  \par\vfill\nopagebreak
  \begin{center}\popsquiggle{poporange}\end{center}\vspace{4pt}
  \signatureonbuch
}
\AtBeginDocument{\def\llbracket{\mathopen{[\mkern-3.5mu[}}\def\rrbracket{\mathclose{]\mkern-3.5mu]}}} % intervalles d'entiers (glyphe ⟦ absent de la police)
% Glyphes absents de Fira Math : redéfinis à partir de glyphes disponibles
\AtBeginDocument{%
  \def\setminus{\mathbin{\backslash}}\let\smallsetminus\setminus
  \def\vdots{\mathord{\vbox{\baselineskip3.2pt\lineskiplimit0pt\kern4pt\hbox{.}\hbox{.}\hbox{.}}}}%
  \def\ddots{\mathinner{\mkern1mu\raise6.5pt\hbox{.}\mkern2mu\raise3.6pt\hbox{.}\mkern2mu\raise0.7pt\hbox{.}\mkern1mu}}%
  \def\triangle{\mathord{\scalebox{0.9}{$\symup{\Delta}$}}}%
  \def\nabla{\mathord{\raisebox{\depth}{\scalebox{1}[-1]{$\symup{\Delta}$}}}}%
  \def\checkmark{\popcheck{popdark}}%
  \def\star{\mathbin{\ast}}\def\bigstar{\mathord{\ast}}%
  \def\diamond{\mathbin{\scalebox{0.6}{$\Diamond$}}}%
  \def\bigcup{\mathop{\mathchoice{\raisebox{-0.3em}{\scalebox{1.7}{$\cup$}}}{\scalebox{1.25}{$\cup$}}{\cup}{\cup}}\displaylimits}%
  \def\bigcap{\mathop{\mathchoice{\raisebox{-0.3em}{\scalebox{1.7}{$\cap$}}}{\scalebox{1.25}{$\cap$}}{\cap}{\cap}}\displaylimits}%
  \def\lesseqqgtr{\mathrel{\lessgtr}}\def\bigoplus{\mathop{\oplus}\displaylimits}%
}
\providecommand{\ssi}{si et seulement si} % abréviation fréquente
\AtBeginDocument{\providecommand{\wideparen}{\overparen}} % arc de cercle

\begin{document}

\pagegarde{Transformations affines du plan}{Mathématiques}{1re Tech}{Mathématiques – classe de Première technique}{N15}{2 séances (fiche de progression harmonisée)}{Cette leçon étudie la composition des transformations affines usuelles du plan : composée de deux rotations de même sens, de deux homothéties de même centre, et d'une homothétie suivie d'une translation. L'élève apprend à identifier la nature de la transformation composée, à déterminer ses éléments caractéristiques et à exploiter ces résultats dans des situations concrètes de la vie courante.
\vspace{4pt}
\begin{multicols}{2}\small
\begin{itemize}
\item À la fin de cette leçon, je sais définir la composée de deux transformations du plan et déterminer l'image d'un point par une composée
\item À la fin de cette leçon, je sais démontrer que la composée de deux rotations de même centre est une rotation (ou l'identité) et déterminer son angle
\item À la fin de cette leçon, je sais déterminer la nature et les éléments caractéristiques de la composée de deux rotations de centres distincts
\item À la fin de cette leçon, je sais démontrer que la composée de deux homothéties de même centre est une homothétie de même centre (ou l'identité) et calculer son rapport
\item À la fin de cette leçon, je sais déterminer la nature de la composée d'une homothétie et d'une translation selon la valeur du rapport
\item À la fin de cette leçon, je sais construire l'image d'une figure simple par une composée de transformations
\end{itemize}\end{multicols}}

\begin{sommaire}
\begin{multicols}{2}
\begin{enumerate}
\item Composée de deux transformations : généralités
\item Composée de deux rotations de même centre
\item Composée de deux rotations de centres distincts
\item Composée de deux homothéties de même centre
\item Composée d'une homothétie et d'une translation
\item Synthèse, méthodes et applications concrètes
\item Activité d'intégration
\item Méthodes et exercices résolus
\item Exercices
\item Corrigés des exercices
\item Fiche bilan
\end{enumerate}
\end{multicols}
\end{sommaire}

\begin{objectifs}
\begin{itemize}
\item À la fin de cette leçon, je sais définir la composée de deux transformations du plan et déterminer l'image d'un point par une composée
\item À la fin de cette leçon, je sais démontrer que la composée de deux rotations de même centre est une rotation (ou l'identité) et déterminer son angle
\item À la fin de cette leçon, je sais déterminer la nature et les éléments caractéristiques de la composée de deux rotations de centres distincts
\item À la fin de cette leçon, je sais démontrer que la composée de deux homothéties de même centre est une homothétie de même centre (ou l'identité) et calculer son rapport
\item À la fin de cette leçon, je sais déterminer la nature de la composée d'une homothétie et d'une translation selon la valeur du rapport
\item À la fin de cette leçon, je sais construire l'image d'une figure simple par une composée de transformations
\item À la fin de cette leçon, je sais appliquer les composées de transformations à des situations concrètes (agrandissements, réductions, pavages, mécanique)
\item À la fin de cette leçon, je sais rédiger et justifier une démarche complète : identification de la transformation, calcul des éléments caractéristiques, conclusion
\end{itemize}
\end{objectifs}

\begin{prerequis}
\begin{itemize}
\item Notion de transformation du plan : translation, rotation, homothétie (cours précédents de Première)
\item Éléments caractéristiques : centre et angle d'une rotation, centre et rapport d'une homothétie, vecteur d'une translation
\item Relation de Chasles sur les vecteurs et calcul vectoriel (Seconde / début de Première)
\item Égalité de vecteurs, colinéarité et coordonnées dans un repère
\item Propriétés de conservation : alignement, distances (isométries), rapports de distances (homothéties)
\end{itemize}
\end{prerequis}

\newpage

\coursec{Composée de deux transformations : généralités}

\textbf{Situation d'accroche.} Un imprimeur de Douala doit reproduire un logo sur des affiches de tailles différentes : il commence par \emph{agrandir} le logo, puis il le \emph{décale} sur chaque support pour bien le positionner. Un tailleur de Yaoundé agrandit un patron de 150\,\% puis le \emph{translate} pour l'adapter à un nouveau tissu. Un mécanicien fait tourner deux engrenages solidaires l'un de l'autre : le mouvement d'une dent résulte de \emph{deux rotations successives}. Dans chaque cas, on \cle{enchaîne} deux transformations du plan.

\textbf{Problématique.} Que devient la figure finale quand on applique deux transformations l'une après l'autre ? Est-ce encore une transformation simple (translation, rotation, homothétie) ? L'ordre dans lequel on enchaîne les transformations a-t-il de l'importance ? Cette leçon répond à ces questions en étudiant la \cle{composée} de deux rotations, de deux homothéties, et d'une homothétie avec une translation. Cette première section pose les bases : la définition de la composée, la construction de l'image d'un point, et les pièges à éviter.

\courssub{Rappels : les transformations usuelles du plan}

Avant d'enchaîner des transformations, rappelons brièvement les trois transformations étudiées dans les leçons précédentes, leurs \cle{éléments caractéristiques} et leurs \cle{propriétés de conservation}.

\begin{definition}[Transformations usuelles]
Soit $M$ un point du plan.
\begin{itemize}
\item La \cle{translation} de vecteur $\vec{u}$, notée $t_{\vec{u}}$, associe à $M$ le point $M'$ tel que $\overrightarrow{MM'} = \vec{u}$. Elle est caractérisée par son \textbf{vecteur}.
\item La \cle{rotation} de centre $O$ et d'angle $\alpha$, notée $r_{(O,\alpha)}$, associe à $M$ (distinct de $O$) le point $M'$ tel que $OM' = OM$ et $\left(\overrightarrow{OM}, \overrightarrow{OM'}\right) = \alpha$. Elle est caractérisée par son \textbf{centre} et son \textbf{angle}.
\item L'\cle{homothétie} de centre $O$ et de rapport $k$ ($k \neq 0$), notée $h_{(O,k)}$, associe à $M$ le point $M'$ tel que $\overrightarrow{OM'} = k\,\overrightarrow{OM}$. Elle est caractérisée par son \textbf{centre} et son \textbf{rapport}.
\end{itemize}
\end{definition}

\begin{propriete}[Propriétés de conservation]
\begin{itemize}
\item La translation et la rotation sont des \cle{isométries} : elles conservent les \textbf{distances}, les \textbf{angles}, les \textbf{aires} et l'\textbf{alignement}.
\item L'homothétie de rapport $k$ \textbf{multiplie les distances par $|k|$} et les \textbf{aires par $k^2$} ; elle conserve les angles, l'alignement et transforme une droite en une droite \textbf{parallèle}.
\end{itemize}
\end{propriete}

\begin{attention}[Bien distinguer les trois transformations]
Ne confonds pas les éléments caractéristiques : une translation n'a \textbf{pas de centre}, une homothétie de rapport $1$ est l'identité, et une homothétie de rapport $-1$ est une symétrie centrale (donc aussi une rotation d'angle $180\degres$). Au Probatoire, l'examinateur vérifie d'abord que tu sais identifier correctement chaque transformation.
\end{attention}

\courssub{Définition de la composée de deux transformations}

Revenons à l'imprimeur de Douala. Pour obtenir son affiche, il n'applique pas une transformation, mais \emph{deux à la suite} : d'abord l'agrandissement (une homothétie), ensuite le décalage (une translation). Mathématiquement, cet enchaînement s'appelle une \cle{composée}.

L'idée est simple : pour connaître l'image finale d'un point $M$, on applique d'abord la première transformation $g$, ce qui donne un point intermédiaire $M' = g(M)$, puis on applique la deuxième transformation $f$ à ce point intermédiaire, ce qui donne le point final $M'' = f(M')$.

\begin{definition}[Composée de deux transformations]
Soient $f$ et $g$ deux transformations du plan. La \cle{composée} de $g$ par $f$, notée $f \circ g$ (lire « $f$ rond $g$ »), est la transformation qui à tout point $M$ du plan associe le point $f\big(g(M)\big)$ :
\[ f \circ g \,:\, M \longmapsto f\big(g(M)\big). \]
Autrement dit, dans l'écriture $f \circ g$, c'est \textbf{$g$ qui agit en premier}, puis $f$ agit sur le résultat.
\end{definition}

\begin{aretenir}[Lecture de la notation]
La notation $f \circ g$ se lit de \textbf{droite à gauche} dans l'ordre d'application : $g$ d'abord, $f$ ensuite. On dit aussi « $f$ rond $g$ » ou « $f$ composée avec $g$ ». Schéma de l'enchaînement :
\[ M \xrightarrow{\;g\;} M' = g(M) \xrightarrow{\;f\;} M'' = f(M') = f \circ g\,(M). \]
\end{aretenir}

\begin{attention}[Erreur fréquente : le mauvais ordre]
L'erreur la plus fréquente au Probatoire est d'appliquer les transformations dans le mauvais ordre. Dans $f \circ g$, on applique \textbf{d'abord $g$} (celle qui est écrite à droite, au plus près du point), \textbf{puis $f$}. Retiens l'image d'une voiture à deux étages de lavage : dans $f \circ g$, le point $M$ rencontre d'abord $g$, car $g$ est « collé » à $M$ dans l'écriture $f\big(g(M)\big)$.
\end{attention}

\begin{propriete}[La composée est une transformation]
La composée de deux transformations du plan est encore une \cle{transformation du plan} : tout point $M$ possède une image unique $M'' = f \circ g\,(M)$, et tout point possède un antécédent unique (on peut « remonter » en appliquant les transformations réciproques dans l'ordre inverse).
\end{propriete}

\courssub{Construction de l'image d'un point par une composée}

\begin{methode}[Déterminer l'image d'un point par $f \circ g$]
Pour construire ou calculer l'image d'un point $M$ par la composée $f \circ g$ :
\begin{enumerate}
\item \textbf{Étape 1} : construire (ou calculer) l'image intermédiaire $M' = g(M)$ en appliquant la \emph{première} transformation $g$ ;
\item \textbf{Étape 2} : construire (ou calculer) l'image finale $M'' = f(M')$ en appliquant la \emph{deuxième} transformation $f$ au point $M'$ ;
\item \textbf{Étape 3} : conclure : $M'' = f \circ g\,(M)$.
\end{enumerate}
Pour une figure entière, on applique ces étapes à ses points remarquables (sommets), puis on relie les images.
\end{methode}

\begin{popfigure}
\begin{center}
\begin{tikzpicture}[scale=1.1,>=Stealth]
\coordinate (M) at (0,0);
\coordinate (Mp) at (2.6,0.9);
\coordinate (Mpp) at (4.6,2.6);
\fill[popblue] (M) circle (2pt);
\fill[poporange] (Mp) circle (2pt);
\fill[poppurple] (Mpp) circle (2pt);
\node[below left, popblue] at (M) {$M$};
\node[below, poporange] at (Mp) {$M' = g(M)$};
\node[above right, poppurple] at (Mpp) {$M'' = f(M')$};
\draw[->, very thick, poporange] (M) to[bend left=25] node[above]{$g$ (en premier)} (Mp);
\draw[->, very thick, poppurple] (Mp) to[bend left=25] node[above right]{$f$ (ensuite)} (Mpp);
\draw[->, thick, popink, dashed] (M) to[bend right=30] node[below]{$f \circ g$} (Mpp);
\end{tikzpicture}
\end{center}
\legende{Image d'un point par une composée : on applique d'abord $g$, puis $f$. La flèche en pointillés représente la transformation composée $f \circ g$.}
\end{popfigure}

\begin{exoresolu}[Translation suivie d'une homothétie]
Le plan est muni d'un repère $(O\,;\vec{\imath},\vec{\jmath})$. Soit $t$ la translation de vecteur $\vec{u}\begin{pmatrix} 3 \\ 1 \end{pmatrix}$ et $h$ l'homothétie de centre $O$ et de rapport $2$. On considère le point $A(1\,;2)$.

Déterminer les coordonnées de $A''$, image de $A$ par la composée $h \circ t$.
\tcblower
\textbf{Solution détaillée.}

Dans $h \circ t$, c'est la translation $t$ qui agit \textbf{en premier}.

\textbf{Étape 1 : image intermédiaire $A' = t(A)$.} La translation de vecteur $\vec{u}(3\,;1)$ ajoute $3$ à l'abscisse et $1$ à l'ordonnée :
\[ A' = \big(1 + 3\,;\, 2 + 1\big) = (4\,;\,3). \]

\textbf{Étape 2 : image finale $A'' = h(A')$.} L'homothétie de centre $O$ et de rapport $2$ multiplie les coordonnées par $2$ (car $\overrightarrow{OA''} = 2\,\overrightarrow{OA'}$) :
\[ A'' = \big(2 \times 4\,;\, 2 \times 3\big) = (8\,;\,6). \]

\textbf{Conclusion.} $h \circ t\,(A) = A''(8\,;\,6)$.

\emph{Vérification graphique mentale :} le point $A(1;2)$ est d'abord décalé de 3 unités vers la droite et 1 vers le haut, puis « éloigné » de $O$ en doublant ses distances. Le résultat est cohérent.
\end{exoresolu}

\courssub{L'ordre des transformations : en général $f \circ g \neq g \circ f$}

Voici un point capital de la leçon. Reprenons le tailleur de Yaoundé : agrandir le patron \emph{puis} le déplacer ne donne pas le même résultat que le déplacer \emph{puis} l'agrandir. En effet, si l'on déplace d'abord puis qu'on agrandit, le déplacement lui-même est agrandi !

\begin{propriete}[Non-commutativité de la composition]
En général, la composition des transformations n'est \textbf{pas commutative} :
\[ f \circ g \;\neq\; g \circ f \quad \text{en général.} \]
Il existe des cas particuliers où l'égalité a lieu (par exemple deux translations, ou deux homothéties de même centre), mais on ne peut \emph{jamais} l'affirmer sans justification.
\end{propriete}

\begin{exemplebox}[Contre-exemple : rotation et translation]
Soit $O$ un point fixe, $r$ la rotation de centre $O$ et d'angle $90\degres$ (sens direct), et $t$ la translation de vecteur $\vec{u}$ horizontal dirigé vers la droite. Prenons un point $M$ situé à droite de $O$.
\begin{itemize}
\item \textbf{Par $t \circ r$} : la rotation envoie $M$ en $M_1$ au-dessus de $O$, puis la translation décale $M_1$ horizontalement en $M_2$.
\item \textbf{Par $r \circ t$} : la translation décale d'abord $M$ horizontalement en $N_1$, puis la rotation envoie $N_1$ en $N_2$, au-dessus d'un point \emph{décalé}.
\end{itemize}
Les points finaux $M_2$ et $N_2$ sont \textbf{distincts} : $t \circ r \neq r \circ t$.
\end{exemplebox}

\begin{popfigure}
\begin{center}
\begin{tikzpicture}[scale=1.35,>=Stealth]
\coordinate (O) at (0,0);
\coordinate (M) at (1.6,0);
\coordinate (M1) at (0,1.6);
\coordinate (M2) at (1.1,1.6);
\coordinate (N1) at (2.7,0);
\coordinate (N2) at (0,2.7);
\fill[popdark] (O) circle (1.8pt); \node[below left] at (O) {$O$};
\fill[popblue] (M) circle (2pt); \node[below, popblue] at (M) {$M$};
% chemin t o r : rotation puis translation
\fill[poporange] (M1) circle (2pt); \node[left, poporange] at (M1) {$M_1=r(M)$};
\fill[poporange] (M2) circle (2pt); \node[above, poporange] at (M2) {$M_2=t(M_1)$};
\draw[->, thick, poporange] (M) arc[start angle=0, end angle=90, radius=1.6] node[midway, above left=-2pt]{$r$};
\draw[->, thick, poporange] (M1) -- (M2) node[midway, above]{$t$};
% chemin r o t : translation puis rotation
\fill[poppurple] (N1) circle (2pt); \node[below, poppurple] at (N1) {$N_1=t(M)$};
\fill[poppurple] (N2) circle (2pt); \node[above left, poppurple] at (N2) {$N_2=r(N_1)$};
\draw[->, thick, poppurple] (M) -- (N1) node[midway, below]{$t$};
\draw[->, thick, poppurple] (N1) arc[start angle=0, end angle=90, radius=2.7] node[midway, above left=-2pt]{$r$};
\draw[dashed, popmuted] (M2) -- (N2);
\node[popink, align=center] at (2.9,2.2) {$M_2 \neq N_2$};
\end{tikzpicture}
\end{center}
\legende{Contre-exemple : en orange, le chemin $t \circ r$ (rotation puis translation) aboutit en $M_2$ ; en violet, le chemin $r \circ t$ (translation puis rotation) aboutit en $N_2$. Les deux résultats sont différents.}
\end{popfigure}

\courssub{Composée d'une transformation et de sa réciproque}

Chaque transformation étudiée possède une \cle{transformation réciproque} qui « défait » son action :
\begin{itemize}
\item la réciproque de la translation $t_{\vec{u}}$ est la translation $t_{-\vec{u}}$ ;
\item la réciproque de la rotation $r_{(O,\alpha)}$ est la rotation $r_{(O,-\alpha)}$ (même centre, angle opposé) ;
\item la réciproque de l'homothétie $h_{(O,k)}$ est l'homothétie $h_{(O,\,1/k)}$ (même centre, rapport inverse).
\end{itemize}

\begin{definition}[Transformation identique]
La \cle{transformation identique} du plan, notée $\mathrm{Id}$, est la transformation qui laisse \textbf{tous les points invariants} : pour tout point $M$, $\mathrm{Id}(M) = M$.
\end{definition}

\begin{propriete}[Composée d'une transformation et de sa réciproque]
Si $f$ est une transformation du plan de réciproque $f^{-1}$, alors :
\[ f \circ f^{-1} = f^{-1} \circ f = \mathrm{Id}. \]
Autrement dit, appliquer une transformation puis sa réciproque (ou l'inverse) ramène chaque point à sa position de départ.
\end{propriete}

\begin{exemplebox}[Vérification sur un exemple]
Soit $h$ l'homothétie de centre $O$ et de rapport $3$. Sa réciproque est $h^{-1}$, homothétie de centre $O$ et de rapport $\dfrac{1}{3}$. Pour un point $M$ quelconque :
\[ \overrightarrow{OM'} = 3\,\overrightarrow{OM} \quad \text{puis} \quad \overrightarrow{OM''} = \frac{1}{3}\,\overrightarrow{OM'} = \frac{1}{3}\times 3\,\overrightarrow{OM} = \overrightarrow{OM}, \]
donc $M'' = M$ : on retrouve bien $h^{-1} \circ h = \mathrm{Id}$.
\end{exemplebox}

\courssub{Lien avec la vie courante et les métiers techniques}

Les composées de transformations sont partout autour de nous, en particulier dans les métiers techniques :

\begin{itemize}
\item \textbf{Imprimerie (Douala).} Reproduire un logo sur des affiches de tailles différentes : on applique une \emph{homothétie} (agrandissement ou réduction) \emph{puis} une \emph{translation} pour positionner le logo sur le support. L'affiche finale est l'image du logo par $t \circ h$.
\item \textbf{Couture (Yaoundé).} Le tailleur agrandit un patron de 150\,\% (homothétie de rapport $1{,}5$) puis le translate pour l'adapter à la laize du nouveau tissu : la pièce de tissu découpée est l'image du patron initial par une composée.
\item \textbf{Mécanique.} Deux engrenages solidaires tournent l'un après l'autre : le déplacement d'une dent résulte de la composée de deux rotations. Nous verrons dans les sections suivantes que cette composée est encore une rotation (ou une translation), ce qui permet de prévoir la position finale des pièces.
\item \textbf{CAO et dessin industriel.} Dans les logiciels de conception, chaque pièce dessinée subit des enchaînements de rotations, homothéties et translations avant d'être envoyée à la machine-outil.
\end{itemize}

\begin{savaistu}[Les logiciels de retouche d'image]
Quand tu tournes une photo de $30\degres$ puis que tu la redimensionnes sur ton téléphone, le logiciel applique en réalité une \textbf{composée de transformations} à \emph{chaque pixel} de l'image : d'abord la rotation, puis l'homothétie (redimensionnement). Une photo de 12 millions de pixels subit donc 12 millions de fois le calcul $f\big(g(M)\big)$ ! Les mathématiciens parlent de « transformations affines », et c'est exactement le même principe qui fait fonctionner les effets spéciaux au cinéma et les jeux vidéo.
\end{savaistu}

\begin{exercice}[À toi de jouer]
Le plan est muni d'un repère $(O\,;\vec{\imath},\vec{\jmath})$. Soit $t$ la translation de vecteur $\vec{u}\begin{pmatrix} 3 \\ 1 \end{pmatrix}$ et $h$ l'homothétie de centre $O$ et de rapport $2$. On reprend le point $A(1\,;2)$.
\begin{enumerate}
\item Déterminer les coordonnées de l'image de $A$ par $t \circ h$.
\item Comparer avec le résultat de $h \circ t$ obtenu dans l'exercice résolu. Conclure.
\end{enumerate}
\end{exercice}

\begin{corrige}[Exercice]
\begin{enumerate}
\item Dans $t \circ h$, c'est $h$ qui agit en premier : $A' = h(A) = (2\times 1\,;\,2\times 2) = (2\,;\,4)$, puis $A'' = t(A') = (2+3\,;\,4+1) = (5\,;\,5)$. Donc $t \circ h\,(A) = (5\,;\,5)$.
\item On a obtenu $h \circ t\,(A) = (8\,;\,6)$ et $t \circ h\,(A) = (5\,;\,5)$ : les résultats sont différents. On vérifie ainsi, par le calcul, que $h \circ t \neq t \circ h$ : l'ordre des transformations est essentiel.
\end{enumerate}
\end{corrige}

\begin{aretenir}[L'essentiel de la section]
\begin{itemize}
\item La composée $f \circ g$ applique \textbf{d'abord $g$, puis $f$} : $M \mapsto f\big(g(M)\big)$.
\item La composée de deux transformations du plan est encore une transformation du plan.
\item En général $f \circ g \neq g \circ f$ : \textbf{l'ordre compte}.
\item Si $f^{-1}$ est la réciproque de $f$, alors $f \circ f^{-1} = f^{-1} \circ f = \mathrm{Id}$.
\item Pour construire l'image d'un point : image intermédiaire d'abord, image finale ensuite.
\end{itemize}
\end{aretenir}

\coursec{Composée de deux rotations de même centre}

\textbf{Situation.} Un tourneur sur bois à Bafoussam réalise une pièce symétrique : il programme une première rotation de $60\degres$ sur la commande numérique, puis une seconde rotation de $40\degres$ sur le même centre (l'axe de la pièce). Quel est le mouvement global équivalent ?

\courssub{Résultat : une rotation de même centre, angle somme}

\begin{propriete}[Composée de deux rotations de même centre]
Soient $r_1 = r_{(O,\alpha)}$ et $r_2 = r_{(O,\beta)}$ deux rotations de \textbf{même centre} $O$ et d'angles $\alpha$ et $\beta$ (orientés dans le même sens, c'est-à-dire tous deux directs ou tous deux indirects). Leur composée $r_2 \circ r_1$ est la rotation de centre $O$ et d'angle $\alpha + \beta$ :
\[ r_{(O,\beta)} \circ r_{(O,\alpha)} = r_{(O,\,\alpha+\beta)}. \]
\end{propriete}

\begin{experience}[Démonstration guidée]
Soit $M$ un point distinct de $O$.
\begin{enumerate}
\item $M' = r_1(M)$ : par définition, $OM' = OM$ et $\left(\overrightarrow{OM}, \overrightarrow{OM'}\right) = \alpha$.
\item $M'' = r_2(M')$ : par définition, $OM'' = OM' = OM$ et $\left(\overrightarrow{OM'}, \overrightarrow{OM''}\right) = \beta$.
\item Par relation de Chasles sur les angles orientés : $\left(\overrightarrow{OM}, \overrightarrow{OM''}\right) = \left(\overrightarrow{OM}, \overrightarrow{OM'}\right) + \left(\overrightarrow{OM'}, \overrightarrow{OM''}\right) = \alpha + \beta$.
\item $M''$ vérifie donc $OM'' = OM$ et $\left(\overrightarrow{OM}, \overrightarrow{OM''}\right) = \alpha+\beta$, ce qui est exactement la définition de l'image de $M$ par la rotation $r_{(O,\alpha+\beta)}$.
\end{enumerate}
Le centre $O$ est invariant par les deux rotations, donc invariant par la composée.
\end{experience}

\begin{attention}[Sens des angles]
La formule $\alpha+\beta$ suppose que les angles sont \textbf{orientés dans le même sens} (tous deux directs $>0$ ou tous deux indirects $<0$). Si les sens sont opposés, on fait une \textbf{somme algébrique} (ex: $+60\degres$ et $-40\degres$ donnent $+20\degres$). Au Probatoire, précise toujours le sens (direct/indirect) ou utilise les angles orientés.
\end{attention}

\begin{exemplebox}[Cas particulier : angle somme multiple de $360\degres$]
Si $\alpha + \beta = k \times 360\degres$ ($k \in \mathbb{Z}$), la composée est la \textbf{transformation identique} $\mathrm{Id}$ (chaque point retrouve sa place initiale). C'est le cas de deux rotations de $180\degres$ de même centre (deux symétries centrales successives).
\end{exemplebox}

\begin{exoresolu}[Application numérique]
Soient $r_1$ la rotation de centre $O$ et d'angle $+70\degres$ (sens direct) et $r_2$ la rotation de centre $O$ et d'angle $+110\degres$ (sens direct). Déterminer la nature et les caractéristiques de $r_2 \circ r_1$.
\tcblower
\textbf{Solution.}
Les deux rotations ont même centre $O$ et même sens (direct). D'après la propriété :
\[ r_2 \circ r_1 = r_{(O,\,70\degres+110\degres)} = r_{(O,\,180\degres)}. \]
La composée est une rotation de centre $O$ et d'angle $180\degres$, c'est-à-dire une \textbf{symétrie centrale} de centre $O$.
\end{exoresolu}

\coursec{Composée de deux rotations de centres distincts}

\textbf{Situation.} Dans une boîte de vitesses, un engrenage $A$ (centre $O_1$) entraîne un engrenage $B$ (centre $O_2$). Une dent de $A$ tourne autour de $O_1$, puis son contact fait tourner une dent de $B$ autour de $O_2$. Le mouvement global de la dent de $B$ par rapport au châssis est une composée de deux rotations de centres distincts.

\courssub{Théorème : nature de la composée}

\begin{propriete}[Composée de deux rotations de centres distincts]
Soient $r_1 = r_{(O_1,\alpha)}$ et $r_2 = r_{(O_2,\beta)}$ deux rotations de centres \textbf{distincts} $O_1 \neq O_2$ et d'angles orientés $\alpha$, $\beta$ (même sens).
\begin{itemize}
\item Si $\alpha + \beta \not\equiv 0 \pmod{360\degres}$, la composée $r_2 \circ r_1$ est une \textbf{rotation} d'angle $\alpha + \beta$. Son centre $O$ est l'unique point invariant (point fixe) de la composée.
\item Si $\alpha + \beta \equiv 0 \pmod{360\degres}$ (c'est-à-dire $\alpha + \beta = k \times 360\degres$), la composée $r_2 \circ r_1$ est une \textbf{translation} de vecteur $\overrightarrow{O_1O_2}$ tourné de l'angle $\alpha$ (ou équivalentement $2\,\overrightarrow{O_1O_2}$ pour deux demi-tours).
\end{itemize}
\end{propriete}

\begin{aretenir}[Construction du centre de la rotation composée (cas $\alpha+\beta \neq 0$)]
Le centre $O$ de la rotation composée $r_2 \circ r_1$ est l'intersection de deux droites :
\begin{enumerate}
\item La médiatrice du segment $\big[O_1\, r_2(O_1)\big]$ (ou la bissectrice de l'angle $\widehat{O_1 O O_2}$).
\item La médiatrice du segment $\big[O_2\, r_1^{-1}(O_2)\big]$.
\end{enumerate}
Plus concrètement : $O$ est le sommet du triangle isocèle $O_1 O O_2$ tel que $\widehat{O_1 O O_2} = \frac{\alpha+\beta}{2}$ (angles orientés).
\end{aretenir}

\begin{experience}[Découverte : composée de deux demi-tours (angles $180\degres$)]
Soient $s_1$ la symétrie centrale de centre $O_1$ (rotation $180\degres$) et $s_2$ la symétrie centrale de centre $O_2$ (rotation $180\degres$). Montrons que $s_2 \circ s_1$ est une translation.
\begin{enumerate}
\item Soit $M$ un point quelconque. $M' = s_1(M)$ est le symétrique de $M$ par rapport à $O_1$ : $O_1$ est le milieu de $[MM']$.
\item $M'' = s_2(M')$ est le symétrique de $M'$ par rapport à $O_2$ : $O_2$ est le milieu de $[M'M'']$.
\item Dans le triangle $MM''M'$, $O_1$ et $O_2$ sont les milieux de $[MM']$ et $[M'M'']$. D'après la propriété du segment joignant les milieux, $\overrightarrow{MM''} = 2\,\overrightarrow{O_1O_2}$.
\item Ce vecteur ne dépend pas de $M$ : la transformation est la translation de vecteur $2\,\overrightarrow{O_1O_2}$.
\end{enumerate}
\end{experience}

\begin{exemplebox}[Construction graphique du centre (exemple)]
Soient $r_1$ rotation de centre $O_1$, angle $+60\degres$ et $r_2$ rotation de centre $O_2$, angle $+60\degres$ ($O_1 \neq O_2$). La composée est une rotation d'angle $+120\degres$. Son centre $O$ est tel que le triangle $O_1 O O_2$ est isocèle en $O$ avec un angle au sommet de $60\degres$ (donc équilatéral). $O$ se construit en traçant des arcs de cercle de $60\degres$ de part et d'autre de $[O_1O_2]$.
\end{exemplebox}

\begin{exoresolu}[Cas translation : deux rotations de $90\degres$ et $-90\degres$]
Soit $r_1$ rotation de centre $O_1$, angle $+90\degres$ et $r_2$ rotation de centre $O_2$, angle $-90\degres$. Déterminer $r_2 \circ r_1$.
\tcblower
\textbf{Solution.}
$\alpha + \beta = 90\degres + (-90\degres) = 0\degres \equiv 0 \pmod{360\degres}$.
La composée est une \textbf{translation}.
Son vecteur est $\vec{v} = \overrightarrow{O_1O_2}$ tourné de l'angle $\alpha = +90\degres$ (sens direct).
Si $\overrightarrow{O_1O_2}$ est horizontal vers la droite, $\vec{v}$ est vertical vers le haut.
\end{exoresolu}

\begin{attention}[Piège Probatoire : centres distincts, angles opposés]
Ne confonds pas « centres distincts, angles opposés » (donne une translation) et « même centre, angles opposés » (donne l'identité). Vérifie toujours si $O_1 = O_2$ ou $O_1 \neq O_2$ avant de conclure.
\end{attention}

\coursec{Composée de deux homothéties de même centre}

\textbf{Situation.} Un architecte utilise un logiciel de CAO : il applique une homothétie de rapport $2$ (agrandissement) puis une homothétie de rapport $3$ (encore un agrandissement), toutes deux centrées sur l'origine du repère du projet. Quel est l'effet global ?

\courssub{Résultat : une homothétie de même centre, rapport produit}

\begin{propriete}[Composée de deux homothéties de même centre]
Soient $h_1 = h_{(O,k_1)}$ et $h_2 = h_{(O,k_2)}$ deux homothéties de \textbf{même centre} $O$ et de rapports $k_1, k_2 \in \mathbb{R}^*$. Leur composée $h_2 \circ h_1$ est l'homothétie de centre $O$ et de rapport $k_1 k_2$ :
\[ h_{(O,k_2)} \circ h_{(O,k_1)} = h_{(O,\,k_1 k_2)}. \]
\end{propriete}

\begin{experience}[Démonstration par les vecteurs]
Soit $M$ un point quelconque.
\begin{enumerate}
\item $M' = h_1(M) \iff \overrightarrow{OM'} = k_1 \overrightarrow{OM}$.
\item $M'' = h_2(M') \iff \overrightarrow{OM''} = k_2 \overrightarrow{OM'} = k_2 (k_1 \overrightarrow{OM}) = (k_1 k_2) \overrightarrow{OM}$.
\item Donc $M'' = h_{(O,k_1 k_2)}(M)$. Le centre $O$ est invariant ($k_1 k_2 \cdot \vec{0} = \vec{0}$).
\end{enumerate}
\end{experience}

\begin{aretenir}[Conséquences importantes]
\begin{itemize}
\item La composée de deux homothéties de même centre est \textbf{commutative} : $h_2 \circ h_1 = h_1 \circ h_2 = h_{(O,k_1 k_2)}$.
\item Si $k_1 k_2 = 1$, la composée est l'\textbf{identité} ($h_2$ est la réciproque de $h_1$).
\item Si $k_1 k_2 = -1$, la composée est une \textbf{symétrie centrale} de centre $O$ (rotation de $180\degres$).
\end{itemize}
\end{aretenir}

\begin{exoresolu}[Calcul de rapport composé]
Une pièce mécanique subit deux agrandissements successifs par homothétie de centre $O$ : d'abord rapport $1{,}5$, puis rapport $4$. Quel est le rapport global ? Si la longueur initiale est $L_0 = 10\,\text{mm}$, quelle est la longueur finale ?
\tcblower
\textbf{Solution.}
Rapport global $k = 1{,}5 \times 4 = 6$.
Longueur finale $L = |k| \times L_0 = 6 \times 10 = 60\,\text{mm}$.
La composée est l'homothétie $h_{(O,6)}$.
\end{exoresolu}

\coursec{Composée d'une homothétie et d'une translation}

\textbf{Situation.} Reprenons l'imprimeur de Douala : il agrandit le logo (homothétie $h$ de centre $O$, rapport $k$) puis le déplace sur l'affiche (translation $t$ de vecteur $\vec{u}$). La transformation globale est $t \circ h$. Est-ce encore une homothétie ? Une translation ?

\courssub{Cas 1 : Homothétie puis Translation ($t \circ h$)}

\begin{propriete}[Composée $t_{\vec{u}} \circ h_{(O,k)}$ avec $k \neq 1$]
Soit $h = h_{(O,k)}$ ($k \neq 1$) et $t = t_{\vec{u}}$. La composée $t \circ h$ est une \textbf{homothétie} de rapport $k$ et de centre $O'$ image de $O$ par la translation $t$ \textbf{composée avec l'homothétie de rapport $\frac{1}{1-k}$}... Non, plus simple :
Le centre $O'$ de l'homothétie composée est le point tel que $\overrightarrow{OO'} = \frac{1}{1-k}\,\vec{u}$.
Autrement dit : $O'$ est l'image de $O$ par l'homothétie de centre $O$ et de rapport $\frac{1}{1-k}$ suivie de la translation... Non, standard :
$t_{\vec{u}} \circ h_{(O,k)} = h_{(O',k)}$ où $O'$ est défini par $\overrightarrow{OO'} = \frac{\vec{u}}{1-k}$.
\end{propriete}

\begin{experience}[Démonstration vectorielle]
Cherchons $O'$ tel que $t_{\vec{u}} \circ h_{(O,k)} = h_{(O',k)}$.
Pour tout $M$ : $\overrightarrow{OM'} = k\overrightarrow{OM}$ puis $\overrightarrow{OM''} = \overrightarrow{OM'} + \vec{u} = k\overrightarrow{OM} + \vec{u}$.
On veut $\overrightarrow{O'M''} = k\overrightarrow{O'M}$.
$\overrightarrow{O'M''} = \overrightarrow{O'O} + \overrightarrow{OM''} = -\overrightarrow{OO'} + k\overrightarrow{OM} + \vec{u}$.
$k\overrightarrow{O'M} = k(\overrightarrow{O'O} + \overrightarrow{OM}) = -k\overrightarrow{OO'} + k\overrightarrow{OM}$.
Égalité $\iff -\overrightarrow{OO'} + \vec{u} = -k\overrightarrow{OO'} \iff \vec{u} = (1-k)\overrightarrow{OO'} \iff \overrightarrow{OO'} = \frac{\vec{u}}{1-k}$.
\end{experience}

\begin{attention}[Condition $k \neq 1$]
Si $k=1$, $h$ est l'identité, la composée est la translation $t_{\vec{u}}$. La formule $\frac{\vec{u}}{1-k}$ a un dénominateur nul : c'est normal car ce n'est plus une homothétie.
\end{attention}

\courssub{Cas 2 : Translation puis Homothétie ($h \circ t$)}

\begin{propriete}[Composée $h_{(O,k)} \circ t_{\vec{u}}$ avec $k \neq 1$]
Soit $h = h_{(O,k)}$ ($k \neq 1$) et $t = t_{\vec{u}}$. La composée $h \circ t$ est une \textbf{homothétie} de rapport $k$ et de centre $O''$ tel que $\overrightarrow{OO''} = \frac{k}{1-k}\,\vec{u}$.
\[ h_{(O,k)} \circ t_{\vec{u}} = h_{(O'',k)} \quad \text{avec} \quad \overrightarrow{OO''} = \frac{k}{1-k}\,\vec{u}. \]
\end{propriete}

\begin{aretenir}[Comparaison des deux ordres]
\begin{center}
\begin{tabularx}{\linewidth}{|l|X|X|}\hline
\rowcolor{popblueL}
\textbf{Composée} & \textbf{Nature} & \textbf{Centre (si $k \neq 1$)} \\ \hline
$t_{\vec{u}} \circ h_{(O,k)}$ & Homothétie rapport $k$ & $O'$ tel que $\overrightarrow{OO'} = \dfrac{\vec{u}}{1-k}$ \\ \hline
$h_{(O,k)} \circ t_{\vec{u}}$ & Homothétie rapport $k$ & $O''$ tel que $\overrightarrow{OO''} = \dfrac{k}{1-k}\,\vec{u}$ \\ \hline
\end{tabularx}
\end{center}
Les centres $O'$ et $O''$ sont \textbf{différents} (sauf si $\vec{u}=\vec{0}$ ou $k=0$ impossible). Cela prouve encore que $t \circ h \neq h \circ t$.
\end{aretenir}

\begin{exemplebox}[Cas particulier $k = -1$ (symétrie centrale)]
$h_{(O,-1)} \circ t_{\vec{u}}$ est une homothétie de rapport $-1$, c'est-à-dire une symétrie centrale. Son centre $O''$ vérifie $\overrightarrow{OO''} = \frac{-1}{2}\vec{u} = -\frac{1}{2}\vec{u}$. Le centre est le milieu de $[O, O - \vec{u}]$...
\end{exemplebox}

\begin{exoresolu}[Application complète : Imprimeur]
Un logo est défini par le point $A(2\,;1)$ dans un repère $(O\,;\vec{\imath},\vec{\jmath})$. L'imprimeur applique une homothétie $h$ de centre $O$ et rapport $3$, puis une translation $t$ de vecteur $\vec{u}(4\,;-2)$.
\begin{enumerate}
\item Déterminer les coordonnées de $A'' = t \circ h(A)$.
\item Déterminer le centre $O'$ de l'homothétie globale $t \circ h$.
\end{enumerate}
\tcblower
\textbf{Solution.}
\begin{enumerate}
\item $A' = h(A) = (6\,;3)$. $A'' = t(A') = (6+4\,;\,3-2) = (10\,;\,1)$.
\item $k=3 \neq 1$. $\overrightarrow{OO'} = \frac{\vec{u}}{1-k} = \frac{(4\,;-2)}{1-3} = \frac{(4\,;-2)}{-2} = (-2\,;1)$.
Donc $O'(-2\,;1)$. Vérification : $h_{(O',3)}(A) \to \overrightarrow{O'A} = (4\,;0) \to \overrightarrow{O'A''} = (12\,;0) \to A'' = O' + (12\,;0) = (10\,;1)$. Correct.
\end{enumerate}
\end{exoresolu}

\coursec{Synthèse, méthodes et applications concrètes}

\courssub{Tableau récapitulatif des composées (Programme Première Tech)}

\begin{center}
\begin{tabularx}{\linewidth}{|l|X|X|}\hline
\rowcolor{popblueL}
\textbf{Composée} & \textbf{Condition} & \textbf{Résultat} \\ \hline
$r_{(O,\beta)} \circ r_{(O,\alpha)}$ & Même centre & Rotation $r_{(O,\alpha+\beta)}$ \\ \hline
$r_{(O_2,\beta)} \circ r_{(O_1,\alpha)}$ & $O_1 \neq O_2$, $\alpha+\beta \not\equiv 0\ [360\degres]$ & Rotation angle $\alpha+\beta$, centre à construire \\ \hline
$r_{(O_2,\beta)} \circ r_{(O_1,\alpha)}$ & $O_1 \neq O_2$, $\alpha+\beta \equiv 0\ [360\degres]$ & Translation (vecteur $2\overrightarrow{O_1O_2}$ si $\alpha=\beta=180\degres$) \\ \hline
$h_{(O,k_2)} \circ h_{(O,k_1)}$ & Même centre & Homothétie $h_{(O,k_1 k_2)}$ (commutative) \\ \hline
$t_{\vec{u}} \circ h_{(O,k)}$ & $k \neq 1$ & Homothétie rapport $k$, centre $O'$ : $\overrightarrow{OO'}=\frac{\vec{u}}{1-k}$ \\ \hline
$h_{(O,k)} \circ t_{\vec{u}}$ & $k \neq 1$ & Homothétie rapport $k$, centre $O''$ : $\overrightarrow{OO''}=\frac{k\vec{u}}{1-k}$ \\ \hline
$t_{\vec{u}} \circ h_{(O,1)}$ ou $h_{(O,1)} \circ t_{\vec{u}}$ & $k=1$ (identité) & Translation $t_{\vec{u}}$ \\ \hline
\end{tabularx}
\end{center}

\courssub{Méthode unifiée : « Comment déterminer la nature d'une composée ? »}

\begin{methode}[Stratégie de résolution au Probatoire]
Face à une composée $f \circ g$ :
\begin{enumerate}
\item \textbf{Identifier} $f$ et $g$ (nature, éléments caractéristiques : centre, angle, vecteur, rapport).
\item \textbf{Vérifier l'ordre} : $g$ d'abord, $f$ ensuite.
\item \textbf{Repérer le cas} dans le tableau ci-dessus (mêmes centres ? centres distincts ? homothétie + translation ?).
\item \textbf{Appliquer la formule} correspondante (somme d'angles, produit de rapports, calcul du nouveau centre).
\item \textbf{Conclure} en énonçant clairement la nature et les caractéristiques de la transformation composée.
\item \textbf{(Optionnel) Vérifier} sur un point particulier (ex: le centre d'une rotation, l'origine pour une homothétie).
\end{enumerate}
\end{methode}

\courssub{Applications techniques et vie courante}

\begin{exemplebox}[Robotique : Bras articulé à deux segments]
Un bras robotique a deux articulations $O_1$ (épaule) et $O_2$ (coude). Le premier segment tourne de $\alpha$ autour de $O_1$, le second de $\beta$ autour de $O_2$ (solidaire du premier). La position finale de l'outil (bout du bras) par rapport à la base $O_1$ résulte de la composée de deux rotations de centres distincts. Si $\alpha+\beta \neq 0$, l'outil décrit un arc de cercle autour d'un centre virtuel $O$ (calculé comme vu en Section 3). Cela permet de programmer la trajectoire sans simuler chaque articulation.
\end{exemplebox}

\begin{exemplebox}[Optique : Système de deux lentilles minces]
En optique géométrique (approximation de Gauss), une lentille mince réalise une homothétie (grandissement) centrée sur son centre optique. Deux lentilles successives de centres $O_1, O_2$ et rapports $g_1, g_2$ (grandissements) réalisent une composée. Si les centres coïncident (lentilles collées), le grandissement total est $g_1 g_2$ (produit des rapports). Si les centres sont distincts, la composée n'est plus une homothétie pure mais une transformation affine plus générale (homothétie + translation), base du calcul des systèmes optiques.
\end{exemplebox}

\begin{savaistu}[Géométrie affine et informatique graphique]
L'ensemble des transformations étudiées (translations, rotations, homothéties, et leurs composées) forme le groupe des \textbf{similitudes} (si on ajoute les symétries axiales, on a les isométries ; avec homothéties de rapport quelconque, les similitudes). En infographie 2D/3D, on représente ces transformations par des \textbf{matrices $3 \times 3$ (coordonnées homogènes)}. La composée devient un simple \textbf{produit de matrices}. C'est ainsi que les GPU (cartes graphiques) calculent l'affichage de millions de triangles par seconde dans les jeux vidéo ou la CAO.
\end{savaistu}

\begin{exercice}[Entraînement Probatoire - Type Bac]
Le plan est muni d'un repère orthonormé $(O\,;\vec{\imath},\vec{\jmath})$.
Soient $h$ l'homothétie de centre $O$ et de rapport $k = -2$, et $t$ la translation de vecteur $\vec{u}\begin{pmatrix} 6 \\ -3 \end{pmatrix}$.
On considère la composée $f = t \circ h$.
\begin{enumerate}
\item Déterminer la nature de $f$ et ses caractéristiques (centre, rapport).
\item Calculer les coordonnées de l'image de $A(1\,;2)$ par $f$ de deux façons : par la définition ($h$ puis $t$) et par la formule de l'homothétie trouvée en 1.
\item Soit $g = h \circ t$. Déterminer la nature et les caractéristiques de $g$. Comparer $f$ et $g$.
\end{enumerate}
\end{exercice}

\begin{corrige}[Exercice]
\begin{enumerate}
\item $f = t \circ h$ avec $h = h_{(O,-2)}$, $k=-2 \neq 1$. C'est une homothétie de rapport $k=-2$.
Centre $O'$ : $\overrightarrow{OO'} = \frac{\vec{u}}{1-k} = \frac{(6\,;-3)}{1-(-2)} = \frac{(6\,;-3)}{3} = (2\,;-1)$.
Donc $f = h_{(O',-2)}$ avec $O'(2\,;-1)$.
\item \textbf{Par définition :} $A' = h(A) = (-2 \times 1\,;\, -2 \times 2) = (-2\,;-4)$.
$A'' = t(A') = (-2+6\,;\,-4-3) = (4\,;-7)$.
\textbf{Par formule :} $\overrightarrow{O'A} = (1-2\,;\,2-(-1)) = (-1\,;3)$.
$\overrightarrow{O'A''} = -2 \times (-1\,;3) = (2\,;-6)$.
$A'' = O' + (2\,;-6) = (2+2\,;\,-1-6) = (4\,;-7)$. \textbf{Identique.}
\item $g = h \circ t$. $k=-2 \neq 1$. C'est une homothétie de rapport $-2$.
Centre $O''$ : $\overrightarrow{OO''} = \frac{k}{1-k}\vec{u} = \frac{-2}{3}(6\,;-3) = (-4\,;2)$.
Donc $O''(-4\,;2)$.
$f$ a pour centre $O'(2\,;-1)$, $g$ a pour centre $O''(-4\,;2)$. $O' \neq O''$, donc $f \neq g$.
\end{enumerate}
\end{corrige}

\begin{aretenir}[L'essentiel de la leçon (Fiche de révision Probatoire)]
\begin{itemize}
\item \textbf{Composée $f \circ g$} : $g$ puis $f$. Ordre crucial ($f \circ g \neq g \circ f$ en général).
\item \textbf{2 Rotations même centre} $\to$ Rotation même centre, angle somme.
\item \textbf{2 Rotations centres distincts, même sens} :
  \begin{itemize}
  \item $\alpha+\beta \neq 0\ [360\degres]$ $\to$ Rotation angle $\alpha+\beta$, centre à construire (point fixe).
  \item $\alpha+\beta = 0\ [360\degres]$ $\to$ Translation.
  \end{itemize}
\item \textbf{2 Homothéties même centre} $\to$ Homothétie même centre, rapport produit (commutatif).
\item \textbf{Homothétie $k \neq 1$ + Translation} $\to$ Homothétie rapport $k$, centre déplacé.
  \begin{itemize}
  \item $t_{\vec{u}} \circ h_{(O,k)}$ : centre $O'$ tel que $\overrightarrow{OO'} = \frac{\vec{u}}{1-k}$.
  \item $h_{(O,k)} \circ t_{\vec{u}}$ : centre $O''$ tel que $\overrightarrow{OO''} = \frac{k\vec{u}}{1-k}$.
  \end{itemize}
\item \textbf{Réciproque} : $f \circ f^{-1} = \mathrm{Id}$.
\end{itemize}
\end{aretenir}

\coursec{Composée de deux rotations de même centre}

Dans la section précédente, nous avons vu ce qu'est la \cle{composée} de deux transformations : appliquer d'abord $f$, puis $g$, c'est définir une nouvelle transformation notée $g \circ f$. Nous avons aussi rappelé que toute rotation est une \cle{isométrie} (elle conserve les distances). Nous allons maintenant étudier le cas le plus simple et le plus utile : composer deux rotations qui ont le \emph{même centre}. Cette situation est partout autour de nous : une roue que l'on fait tourner deux fois de suite, une aiguille d'horloge qui avance pendant deux intervalles de temps, une vis que l'on serre en deux coups de clé.

\courssub{Situation et conjecture}

\begin{definition}[Situation étudiée]
On considère deux rotations de \emph{même centre} $O$ et de \emph{même sens} (toutes deux directes, c'est-à-dire trigonométriques, ou toutes deux indirectes, c'est-à-dire horaires) :
\begin{itemize}
\item $r_1$ : rotation de centre $O$ et d'angle $\alpha$ ;
\item $r_2$ : rotation de centre $O$ et d'angle $\beta$.
\end{itemize}
On s'intéresse à la composée $r_2 \circ r_1$ : pour tout point $M$, on pose $M' = r_1(M)$ puis $M'' = r_2(M')$, c'est-à-dire $M'' = (r_2 \circ r_1)(M)$.
\end{definition}

Prenons un point $M$ distinct de $O$. La première rotation $r_1$ envoie $M$ sur $M'$ : le point tourne de l'angle $\alpha$ autour de $O$, sans changer de distance à $O$. La deuxième rotation $r_2$ envoie $M'$ sur $M''$ : le point tourne encore, de l'angle $\beta$, toujours autour du même centre $O$. Au total, le point $M$ a tourné autour de $O$ d'un angle $\alpha + \beta$, et sa distance à $O$ n'a jamais changé. On conjecture donc :

\begin{quote}
\emph{La composée $r_2 \circ r_1$ semble être la rotation de centre $O$ et d'angle $\alpha + \beta$.}
\end{quote}

\begin{popfigure}
\begin{center}
\begin{tikzpicture}[scale=1.1]
\coordinate (O) at (0,0);
\coordinate (M) at (3,0);
\coordinate (Mp) at (40:3);
\coordinate (Mpp) at (105:3);
\draw[popline] (O) circle (3);
\fill[popink] (O) circle (1.8pt) node[below left] {$O$};
\fill[popblue] (M) circle (1.8pt) node[right] {$M$};
\fill[poporange] (Mp) circle (1.8pt) node[above right] {$M'=r_1(M)$};
\fill[poppurple] (Mpp) circle (1.8pt) node[above left] {$M''=r_2(M')$};
\draw[popline] (O)--(M);
\draw[popline] (O)--(Mp);
\draw[popline] (O)--(Mpp);
\draw[->,popblue,thick] (1.1,0) arc (0:40:1.1) node[midway,right] {$\alpha$};
\draw[->,poporange,thick] (40:1.5) arc (40:105:1.5) node[midway,above right] {$\beta$};
\draw[->,poppurple,thick,dashed] (2.2,0) arc (0:105:2.2) node[midway,above] {$\alpha+\beta$};
\end{tikzpicture}
\end{center}
\legende{Composée de deux rotations de même centre $O$ : $M$ tourne d'abord de $\alpha$, puis de $\beta$ ; au total, il a tourné de $\alpha+\beta$.}
\end{popfigure}

\courssub{Théorème fondamental}

\begin{propriete}[Composée de deux rotations de même centre]
Soient $r_1$ et $r_2$ deux rotations de même centre $O$, d'angles respectifs $\alpha$ et $\beta$, de même sens. Alors la composée $r_2 \circ r_1$ est la \cle{rotation} de centre $O$ et d'angle $\alpha + \beta$ :
\[ r_2 \circ r_1 = r\bigl(O,\ \alpha + \beta\bigr). \]
\emph{Cas particulier :} si $\alpha + \beta \equiv 0 \pmod{360\degres}$ (notamment $\alpha + \beta = 0$ ou $360\degres$), alors $r_2 \circ r_1$ est la \cle{transformation identique} $\mathrm{Id}$ : tout point est envoyé sur lui-même.
\end{propriete}

\begin{experience}[Démonstration guidée du théorème]
On veut montrer que $r_2 \circ r_1$ est la rotation de centre $O$ et d'angle $\alpha + \beta$. Pour cela, on vérifie les deux conditions qui caractérisent une rotation : la conservation de la distance au centre et la mesure de l'angle orienté.

\textbf{Étape 1 : le centre $O$ est invariant.} $r_1(O) = O$ car le centre d'une rotation est invariant, puis $r_2(O) = O$ pour la même raison. Donc $(r_2 \circ r_1)(O) = O$.

\textbf{Étape 2 : conservation de la distance à $O$.} Soit $M$ un point du plan, $M' = r_1(M)$ et $M'' = r_2(M')$. La rotation $r_1$ est une isométrie, donc $OM' = OM$. La rotation $r_2$ est aussi une isométrie, donc $OM'' = OM'$. Par conséquent :
\[ OM'' = OM' = OM. \]
Le point $M''$ est donc sur le cercle de centre $O$ passant par $M$.

\textbf{Étape 3 : additivité des angles orientés.} Par définition de $r_1$, l'angle orienté $\left(\overrightarrow{OM}, \overrightarrow{OM'}\right)$ mesure $\alpha$. Par définition de $r_2$, l'angle orienté $\left(\overrightarrow{OM'}, \overrightarrow{OM''}\right)$ mesure $\beta$. La relation de Chasles pour les angles orientés donne :
\[ \left(\overrightarrow{OM}, \overrightarrow{OM''}\right) = \left(\overrightarrow{OM}, \overrightarrow{OM'}\right) + \left(\overrightarrow{OM'}, \overrightarrow{OM''}\right) = \alpha + \beta. \]
Cette égalité est valable quel que soit le sens commun des deux rotations, dès lors que les angles sont pris avec leur signe (positif en sens trigonométrique, négatif en sens horaire).

\textbf{Conclusion.} Le point $M''$ vérifie $OM'' = OM$ et $\left(\overrightarrow{OM}, \overrightarrow{OM''}\right) = \alpha + \beta$ : c'est exactement dire que $M''$ est l'image de $M$ par la rotation de centre $O$ et d'angle $\alpha + \beta$. Ceci étant vrai pour tout point $M$, on a bien $r_2 \circ r_1 = r(O, \alpha + \beta)$.

\textbf{Étape 4 : cas particulier.} Si $\alpha + \beta \equiv 0 \pmod{360\degres}$, alors chaque point $M$ est envoyé sur un point $M''$ tel que $OM'' = OM$ et $\left(\overrightarrow{OM}, \overrightarrow{OM''}\right) = 0$, c'est-à-dire $M'' = M$. La composée est l'identité.
\end{experience}

\begin{aretenir}[L'idée à retenir]
Composer deux rotations de même centre revient à \cle{additionner les angles orientés}. Le centre ne change pas. C'est l'une des rares situations où la composée de deux transformations se devine sans calcul : « tourner de $\alpha$ puis de $\beta$, c'est tourner de $\alpha + \beta$ ».
\end{aretenir}

\courssub{Méthode et exemple chiffré}

\begin{methode}[Déterminer la composée de deux rotations de même centre]
\begin{enumerate}
\item Vérifier que les deux rotations ont le \emph{même centre} $O$. Si ce n'est pas le cas, le théorème ne s'applique pas directement (voir section suivante).
\item Repérer le \emph{sens} de chaque rotation. Par convention, le sens trigonométrique (inverse des aiguilles d'une montre) est le sens positif ; le sens horaire correspond à un angle \emph{négatif}.
\item Additionner les angles en tenant compte des signes : $\alpha + \beta$.
\item Conclure : $r_2 \circ r_1$ est la rotation de centre $O$ et d'angle $\alpha + \beta$ ; si la somme vaut $0$ ou un multiple de $360\degres$, c'est l'identité.
\end{enumerate}
\end{methode}

\begin{exemplebox}[Exemple chiffré]
Soient $r_1$ la rotation de centre $O$ et d'angle $40\degres$, et $r_2$ la rotation de centre $O$ et d'angle $65\degres$, toutes deux dans le sens trigonométrique. Déterminons $r_2 \circ r_1$.

Les deux rotations ont le même centre $O$ et le même sens. D'après le théorème :
\[ r_2 \circ r_1 = r\bigl(O,\ 40\degres + 65\degres\bigr) = r(O,\ 105\degres). \]
La composée est la rotation de centre $O$ et d'angle $105\degres$.
\end{exemplebox}

\begin{exoresolu}[Vissage en deux temps]
Un mécanicien serre un écrou en deux coups de clé : au premier coup, l'écrou tourne de $50\degres$ autour de son axe $O$ ; au deuxième coup, il tourne encore de $70\degres$ dans le même sens. De quel angle total l'écrou a-t-il tourné ? Quelle transformation globale a-t-il subie ?
\tcblower
\textbf{Solution.} Le premier coup de clé réalise la rotation $r_1$ de centre $O$ et d'angle $50\degres$ ; le deuxième réalise la rotation $r_2$ de centre $O$ et d'angle $70\degres$, de même sens. L'écrou subit donc la composée $r_2 \circ r_1$. Comme les deux rotations ont le même centre $O$ :
\[ r_2 \circ r_1 = r\bigl(O,\ 50\degres + 70\degres\bigr) = r(O,\ 120\degres). \]
L'écrou a tourné au total de $120\degres$ autour de $O$ : la transformation globale est la rotation de centre $O$ et d'angle $120\degres$.
\end{exoresolu}

\courssub{Rotations réciproques}

\begin{definition}[Rotations réciproques]
Deux rotations de même centre $O$ et d'angles \emph{opposés} $\alpha$ et $-\alpha$ sont dites \cle{réciproques} l'une de l'autre : chacune « annule » l'effet de l'autre.
\end{definition}

En effet, si $r_1 = r(O, \alpha)$ et $r_2 = r(O, -\alpha)$, alors d'après le théorème :
\[ r_2 \circ r_1 = r\bigl(O,\ \alpha + (-\alpha)\bigr) = r(O, 0) = \mathrm{Id}. \]
Composer une rotation avec sa réciproque redonne la transformation identique : tout point revient à sa position de départ. C'est ce qui se passe quand on visse puis on dévisse du même angle, ou quand on tourne une clé de $30\degres$ vers la droite puis de $30\degres$ vers la gauche.

\begin{attention}[Erreur fréquente : oublier le sens de rotation]
L'addition des angles n'est valable que si l'on \emph{respecte les signes}. Par convention, un angle compté dans le sens trigonométrique (anti-horaire) est positif, et un angle compté dans le sens horaire est \emph{négatif}. Ainsi, composer une rotation de $80\degres$ (sens trigonométrique) avec une rotation de $30\degres$ dans le sens horaire donne une rotation d'angle $80\degres - 30\degres = 50\degres$, et non $110\degres$. Toujours écrire les angles avec leur signe avant de les additionner.
\end{attention}

\courssub{Applications concrètes}

\begin{popfigure}
\begin{center}
\begin{tikzpicture}[scale=0.9]
\coordinate (O) at (0,0);
\draw[fill=popblueL,draw=popblue] (O) circle (1.6);
\foreach \a in {0,45,...,315} {
  \draw[popblue,fill=popblueL] (\a:1.6) -- (\a-8:1.95) -- (\a+8:1.95) -- cycle;
}
\draw[fill=white,draw=popdark] (O) circle (0.5);
\fill[popdark] (O) circle (2pt) node[below=6pt] {$O$};
\draw[->,poporange,very thick] (2.6,0) arc (0:60:2.6) node[midway,right] {$\alpha$};
\draw[->,poppurple,very thick] (60:2.9) arc (60:130:2.9) node[midway,above] {$\beta$};
\draw[popdark,thick] (O) -- (60:1.6);
\draw[popdark,thick,dashed] (O) -- (130:1.6);
\node[popmuted] at (0,-2.6) {Rotation totale : $\alpha + \beta$};
\end{tikzpicture}
\end{center}
\legende{Une roue dentée subit deux rotations successives de même axe $O$ : l'angle total de rotation est la somme des deux angles.}
\end{popfigure}

\begin{savaistu}[En horlogerie]
En horlogerie, la position d'une aiguille après deux intervalles de temps s'obtient en additionnant les angles de rotation. L'aiguille des minutes tourne de $6\degres$ par minute (dans le sens horaire, donc $-6\degres$ par minute avec notre convention). Si elle avance de $10$ minutes puis encore de $15$ minutes, elle a tourné au total de $-60\degres + (-90\degres) = -150\degres$, c'est-à-dire $150\degres$ dans le sens horaire : c'est exactement la composée de deux rotations de même centre (l'axe de l'aiguille). Les horlogers et les concepteurs d'engrenages utilisent en permanence cette additivité des angles pour calculer les positions des roues dentées.
\end{savaistu}

\begin{exercice}[À faire]
Dans chaque cas, déterminer la nature et les éléments caractéristiques de $r_2 \circ r_1$.
\begin{enumerate}
\item $r_1 = r(O, 35\degres)$ et $r_2 = r(O, 55\degres)$, même sens trigonométrique.
\item $r_1 = r(O, 110\degres)$ et $r_2 = r(O, -110\degres)$.
\item $r_1 = r(O, 200\degres)$ et $r_2 = r(O, 160\degres)$, même sens.
\end{enumerate}
\end{exercice}

\begin{corrige}[Exercices proposés]
\begin{enumerate}
\item Les deux rotations ont le même centre $O$ et le même sens : $r_2 \circ r_1 = r(O, 35\degres + 55\degres) = r(O, 90\degres)$. C'est le quart de tour direct de centre $O$.
\item Les angles sont opposés : $r_2 \circ r_1 = r(O, 110\degres - 110\degres) = r(O, 0) = \mathrm{Id}$. Les deux rotations sont réciproques.
\item $r_2 \circ r_1 = r(O, 200\degres + 160\degres) = r(O, 360\degres) = \mathrm{Id}$ : un tour complet ramène chaque point à sa place.
\end{enumerate}
\end{corrige}

Que se passe-t-il maintenant si les deux rotations n'ont \emph{pas} le même centre ? C'est ce que nous étudierons dans la section suivante : le résultat reste une rotation d'angle $\alpha + \beta$, mais son centre doit être déterminé avec soin.

\coursec{Composée de deux rotations de centres distincts}

\courssub{Introduction : et si les centres ne sont plus les mêmes ?}

Dans la section précédente, nous avons composé deux rotations de \emph{même centre} $O$ : le résultat était simplement la rotation de centre $O$ et d'angle $\alpha+\beta$. La situation était confortable : le centre commun restait fixe, et tout se résumait à une addition d'angles.

Dans la pratique des ateliers et des chantiers, les rotations successives n'ont presque jamais le même centre. Un bras articulé de robot pivote autour d'une première articulation $O_1$, puis autour d'une seconde articulation $O_2$. Une manivelle entraîne une pièce qui tourne autour d'un autre axe. Que vaut la transformation obtenue en enchaînant deux rotations de centres \emph{distincts} ? C'est l'objet de cette section, et la réponse contient une surprise : dans certains cas, on n'obtient pas une rotation, mais une \cle{translation}.

\courssub{La situation et les notations}

\begin{definition}[Composée de deux rotations de centres distincts]
On considère deux rotations du plan :
\begin{itemize}
\item $r_1$ de centre $O_1$ et d'angle $\alpha$ ;
\item $r_2$ de centre $O_2$ et d'angle $\beta$, avec $O_1 \neq O_2$.
\end{itemize}
La \cle{composée} $r_2 \circ r_1$ est la transformation qui à tout point $M$ associe le point $M''$ obtenu en appliquant d'abord $r_1$ (on obtient $M' = r_1(M)$), puis $r_2$ (on obtient $M'' = r_2(M')$).
\end{definition}

\begin{attention}[Ordre de lecture]
Dans l'écriture $r_2 \circ r_1$, c'est $r_1$ qui agit \textbf{en premier}. On lit de droite à gauche : $M \xrightarrow{r_1} M' \xrightarrow{r_2} M''$.
\end{attention}

\courssub{Le théorème fondamental}

\begin{propriete}[Théorème (admis) : composée de deux rotations de centres distincts]
Soit $r_1$ la rotation de centre $O_1$ et d'angle $\alpha$, et $r_2$ la rotation de centre $O_2$ et d'angle $\beta$, avec $O_1 \neq O_2$. Alors :
\begin{itemize}
\item si $\alpha + \beta \not\equiv 0 \pmod{360°}$, la composée $r_2 \circ r_1$ est une \cle{rotation} d'angle $\alpha + \beta$, dont le centre $\Omega$ se construit géométriquement ;
\item si $\alpha + \beta \equiv 0 \pmod{360°}$, la composée $r_2 \circ r_1$ est une \cle{translation}.
\end{itemize}
Ce théorème est admis au programme de Première technique ; la démonstration n'est pas exigible au Probatoire, mais l'idée est présentée ci-dessous.
\end{propriete}

\begin{experience}[Idée de la démonstration : décomposer en symétries axiales]
Rappelons un fait essentiel : toute rotation de centre $O$ et d'angle $\theta$ est la composée de deux symétries axiales dont les axes passent par $O$ et font entre eux un angle $\theta/2$ (orienté).
Choisissons comme axe commun la droite $(O_1O_2)$, notée $d$.
\begin{enumerate}
\item Écrivons $r_1 = s_d \circ s_{\Delta_1}$ où $\Delta_1$ est une droite passant par $O_1$ telle que l'angle orienté $(\Delta_1, d) = \alpha/2$. On applique donc $s_{\Delta_1}$ puis $s_d$.
\item Écrivons $r_2 = s_{\Delta_2} \circ s_d$ où $\Delta_2$ est une droite passant par $O_2$ telle que l'angle orienté $(d, \Delta_2) = \beta/2$. On applique donc $s_d$ puis $s_{\Delta_2}$.
\item En composant : $r_2 \circ r_1 = (s_{\Delta_2} \circ s_d) \circ (s_d \circ s_{\Delta_1}) = s_{\Delta_2} \circ (s_d \circ s_d) \circ s_{\Delta_1} = s_{\Delta_2} \circ s_{\Delta_1}$ car $s_d \circ s_d = \text{Id}$.
\item La composée de deux symétries axiales $s_{\Delta_2} \circ s_{\Delta_1}$ est bien connue :
  \begin{itemize}
  \item Si les axes $\Delta_1$ et $\Delta_2$ se coupent en un point $\Omega$, la composée est une \textbf{rotation} de centre $\Omega$ et d'angle $2(\Delta_1, \Delta_2) = 2(\alpha/2 + \beta/2) = \alpha + \beta$.
  \item Si les axes $\Delta_1$ et $\Delta_2$ sont parallèles (ce qui arrive exactement quand $\alpha + \beta \equiv 0 \pmod{360°}$, car alors $(\Delta_1, \Delta_2) \equiv 0 \pmod{180°}$), la composée est une \textbf{translation}.
  \end{itemize}
\end{enumerate}
C'est cette alternative « axes sécants ou axes parallèles » qui explique le double visage du théorème.
\end{experience}

\courssub{Construction géométrique du centre \texorpdfstring{$\Omega$}{Omega} (cas $\alpha + \beta \not\equiv 0 \pmod{360°}$)}

L'idée de la démonstration fournit directement la construction du centre $\Omega$ de la rotation composée : c'est le point d'intersection des deux axes $\Delta_1$ et $\Delta_2$, c'est-à-dire le point depuis lequel on « voit » le segment $[O_1O_2]$ sous des angles liés à $\alpha/2$ et $\beta/2$.

\begin{methode}[Construire le centre $\Omega$ de la rotation composée]
\begin{enumerate}
\item Tracer le segment $[O_1O_2]$ reliant les deux centres.
\item En $O_1$, tracer la demi-droite $\Delta_1$ telle que l'angle orienté $\big(\overrightarrow{O_1\Omega}, \overrightarrow{O_1O_2}\big) = \dfrac{\alpha}{2}$. (Autrement dit : on tourne le rayon $\overrightarrow{O_1O_2}$ de $\alpha/2$ vers la gauche si $\alpha>0$, vers la droite si $\alpha<0$, pour obtenir $\overrightarrow{O_1\Omega}$).
\item En $O_2$, tracer la demi-droite $\Delta_2$ telle que $\big(\overrightarrow{O_2O_1}, \overrightarrow{O_2\Omega}\big) = \dfrac{\beta}{2}$.
\item Le point d'intersection de $\Delta_1$ et $\Delta_2$ est le centre $\Omega$ cherché.
\item Vérification : choisir un point $M$, construire $M' = r_1(M)$ puis $M'' = r_2(M')$, et contrôler que $\Omega M = \Omega M''$ et que l'angle orienté $\big(\overrightarrow{\Omega M}, \overrightarrow{\Omega M''}\big)$ vaut $\alpha + \beta$.
\end{enumerate}
\end{methode}

\begin{popfigure}
\begin{center}
\begin{tikzpicture}[scale=1.1, every node/.style={font=\small}]
% Centres
\coordinate (O1) at (0,0);
\coordinate (O2) at (4,0);
\coordinate (Om) at (2,2.6);
% Trajectoire M -> M' -> M''
\coordinate (M) at (-1.6,1.2);
\coordinate (Mp) at (0.9,2.1);
\coordinate (Mpp) at (5.3,2.4);
% Points
\fill[popblue] (O1) circle (1.6pt) node[below left] {$O_1$};
\fill[poporange] (O2) circle (1.6pt) node[below right] {$O_2$};
\fill[popgreen] (Om) circle (1.8pt) node[above] {$\Omega$};
\fill[popdark] (M) circle (1.4pt) node[left] {$M$};
\fill[popdark] (Mp) circle (1.4pt) node[above left] {$M'$};
\fill[popdark] (Mpp) circle (1.4pt) node[right] {$M''$};
% Segment O1O2
\draw[popmuted, dashed] (O1) -- (O2);
% Axes de construction
\draw[popblue, dashed] (O1) -- (Om);
\draw[poporange, dashed] (O2) -- (Om);
% Arcs de rotation
\draw[popblue, -{Stealth}] (M) .. controls (-1.2,2.4) .. (Mp) node[midway, above left, popblue] {$r_1$};
\draw[poporange, -{Stealth}] (Mp) .. controls (3.2,3.3) .. (Mpp) node[midway, above, poporange] {$r_2$};
% Rotation composée
\draw[popgreen, -{Stealth}, thick] (M) .. controls (0.6,3.9) and (3.6,4.1) .. (Mpp) node[midway, above, popgreen] {$r_2\circ r_1$ : rotation d'angle $\alpha+\beta$};
% Rayons depuis Omega
\draw[popgreen!60, thin] (Om) -- (M);
\draw[popgreen!60, thin] (Om) -- (Mpp);
% Angles moitiés
\draw[popblue] (0.9,0) arc (0:52:0.9) node[midway, right, popblue] {$\frac{\alpha}{2}$};
\draw[poporange] (3.1,0) arc (180:128:0.9) node[midway, left, poporange] {$\frac{\beta}{2}$};
\end{tikzpicture}
\end{center}
\legende{Deux rotations de centres distincts $O_1$ et $O_2$ : la trajectoire $M \mapsto M' \mapsto M''$ coïncide avec une unique rotation de centre $\Omega$, intersection des axes faisant les angles $\alpha/2$ et $\beta/2$ avec $(O_1O_2)$.}
\end{popfigure}

\begin{exemplebox}[Composée de deux quarts de tour]
Soit $r_1$ la rotation de centre $O_1$ et d'angle $90°$, et $r_2$ la rotation de centre $O_2$ et d'angle $90°$, avec $O_1 \neq O_2$.
\begin{itemize}
\item Somme des angles : $\alpha + \beta = 90° + 90° = 180° \not\equiv 0 \pmod{360°}$.
\item Donc $r_2 \circ r_1$ est une \textbf{rotation d'angle $180°$}, c'est-à-dire une \cle{symétrie centrale}.
\item Son centre $\Omega$ se construit avec les angles moitiés $\alpha/2 = 45°$ et $\beta/2 = 45°$ : le triangle $O_1O_2\Omega$ est rectangle isocèle en $\Omega$ (les angles en $O_1$ et $O_2$ valent $45°$). Il y a donc deux positions possibles pour $\Omega$, de part et d'autre de $[O_1O_2]$ ; le sens des rotations (direct ici) départage : $\Omega$ est du côté gauche du vecteur $\overrightarrow{O_1O_2}$.
\end{itemize}
\end{exemplebox}

\courssub{Cas particulier fondamental : $\alpha + \beta \equiv 0 \pmod{360°}$, la composée est une translation}

Voici la surprise annoncée. Si les deux rotations ont des angles \emph{opposés} (par exemple $60°$ et $-60°$, ou $90°$ et $-90°$), les deux axes $\Delta_1$ et $\Delta_2$ de la démonstration deviennent parallèles : il n'y a plus de centre $\Omega$, et la composée est une translation.

\begin{propriete}[Composée de deux rotations d'angles opposés]
Si $\alpha + \beta \equiv 0 \pmod{360°}$, alors $r_2 \circ r_1$ est une \cle{translation}. Son vecteur $\vec{u}$ se détermine en calculant l'image d'un point particulier : comme $r_1(O_1) = O_1$ (le centre d'une rotation est invariant), on a
\[ r_2 \circ r_1 (O_1) = r_2(O_1), \]
donc le vecteur de la translation est $\vec{u} = \overrightarrow{O_1\, r_2(O_1)}$.
\end{propriete}

\begin{methode}[Déterminer la nature et les éléments caractéristiques de $r_2 \circ r_1$]
\begin{enumerate}
\item Calculer la somme $\alpha + \beta$ des angles (réduite entre $-180°$ et $180°$ ou modulo $360°$).
\item Si $\alpha + \beta \not\equiv 0 \pmod{360°}$ : conclure que $r_2 \circ r_1$ est une \textbf{rotation d'angle $\alpha + \beta$}, puis construire son centre $\Omega$ avec les angles moitiés $\alpha/2$ et $\beta/2$.
\item Si $\alpha + \beta \equiv 0 \pmod{360°}$ : conclure que $r_2 \circ r_1$ est une \textbf{translation}, puis calculer l'image de $O_1$ : comme $r_1(O_1) = O_1$, il suffit de construire $r_2(O_1)$ ; le vecteur de translation est $\overrightarrow{O_1\, r_2(O_1)}$.
\item Rédiger la conclusion en citant le théorème du cours.
\end{enumerate}
\end{methode}

\begin{exoresolu}[Une composée qui est une translation]
\textbf{Énoncé.} Soit $O_1$ et $O_2$ deux points distincts tels que $O_1O_2 = 4$ cm. On note $r_1$ la rotation de centre $O_1$ et d'angle $90°$ (sens direct), et $r_2$ la rotation de centre $O_2$ et d'angle $-90°$. Déterminer la nature et les éléments caractéristiques de $f = r_2 \circ r_1$.
\tcblower
\textbf{Solution détaillée.}
\begin{enumerate}
\item \textbf{Somme des angles.} $\alpha + \beta = 90° + (-90°) = 0$. D'après le théorème du cours, $f$ est une \textbf{translation}.
\item \textbf{Recherche du vecteur.} On calcule l'image de $O_1$ par $f$ :
\[ f(O_1) = r_2\big(r_1(O_1)\big) = r_2(O_1), \]
car $O_1$ est invariant par sa propre rotation $r_1$.
\item Notons $B = r_2(O_1)$. Par définition de $r_2$ : $O_2B = O_2O_1 = 4$ cm et $\big(\overrightarrow{O_2O_1}, \overrightarrow{O_2B}\big) = -90°$.
Le triangle $O_1O_2B$ est rectangle isocèle en $O_2$. L'hypoténuse $O_1B$ mesure $4\sqrt{2}$ cm.
Le vecteur $\vec{u} = \overrightarrow{O_1B}$ forme un angle de $+45°$ avec $\overrightarrow{O_1O_2}$ (sens direct).
\item \textbf{Conclusion.} $f$ est la translation de vecteur $\vec{u} = \overrightarrow{O_1B}$, où $B$ est l'image de $O_1$ par la rotation de centre $O_2$ et d'angle $-90°$. Sa longueur est $4\sqrt{2}$ cm, sa direction fait un angle de $45°$ avec $\overrightarrow{O_1O_2}$.
\end{enumerate}
\end{exoresolu}

\begin{popfigure}
\begin{center}
\begin{tikzpicture}[scale=1.0, every node/.style={font=\small}]
\coordinate (O1) at (0,0);
\coordinate (O2) at (4,0);
% B = r2(O1) : rotation -90 deg around O2(4,0). O1 rel O2 = (-4,0). Rot -90 -> (0,4). Abs (4,4).
\coordinate (B) at (4,4);
% Points
\fill[popblue] (O1) circle (1.6pt) node[below left] {$O_1$};
\fill[poporange] (O2) circle (1.6pt) node[below right] {$O_2$};
\fill[poppurple] (B) circle (1.6pt) node[above right] {$B = r_2(O_1)$};
\draw[popmuted, dashed] (O1) -- (O2);
\draw[poporange, dashed] (O2) -- (B);
\draw[poporange] (3.2,0.2) arc (180:90:0.8) node[midway, left, poporange] {$-90°$};
% Vecteur de translation
\draw[-{Stealth}, very thick, popgreen] (O1) -- (B) node[midway, above left, popgreen] {$\vec{u}$};
% Un point M et son image M''
\coordinate (M) at (-1.5,1.8);
\coordinate (Mpp) at (2.5,5.8); % M + u
\fill[popdark] (M) circle (1.4pt) node[left] {$M$};
\fill[popdark] (Mpp) circle (1.4pt) node[right] {$M''$};
\draw[-{Stealth}, thick, popgreen] (M) -- (Mpp) node[midway, above, popgreen] {$\vec{u}$};
% Trajectoire M -> M' -> M''
\coordinate (Mp) at (-1.8,-1.5); % r1(M) center O1 angle 90
\fill[popdark] (Mp) circle (1.2pt) node[below] {$M'$};
\draw[popblue, -{Stealth}] (M) .. controls (-2.5,1.0) .. (Mp) node[midway, left, popblue] {$r_1$};
\draw[poporange, -{Stealth}] (Mp) .. controls (1.0,2.0) .. (Mpp) node[midway, above right, poporange] {$r_2$};
\end{tikzpicture}
\end{center}
\legende{Cas $\alpha + \beta = 0$ : la composée est une translation. Tous les points se déplacent du même vecteur $\vec{u} = \overrightarrow{O_1B}$, où $B = r_2(O_1)$. Ici $\vec{u}$ fait $45°$ avec $\overrightarrow{O_1O_2}$.}
\end{popfigure}

\begin{exemplebox}[Composée d'une rotation de $60°$ et d'une rotation de $-60°$]
Soit $r_1$ de centre $O_1$ et d'angle $60°$, et $r_2$ de centre $O_2$ et d'angle $-60°$. Comme $60° + (-60°) = 0$, la composée $r_2 \circ r_1$ est une translation. Son vecteur est $\overrightarrow{O_1B}$ où $B = r_2(O_1)$ : le triangle $O_1O_2B$ est ici \emph{équilatéral} (car $O_2B = O_2O_1$ et l'angle en $O_2$ vaut $60°$ en valeur absolue), donc la translation a pour longueur $O_1O_2$ et pour direction celle qui fait un angle de $+60°$ avec $\overrightarrow{O_1O_2}$ (sens direct).
\end{exemplebox}

\begin{popfigure}
\begin{center}
\begin{tikzpicture}[scale=1.05, every node/.style={font=\small}]
\coordinate (O1) at (0,0);
\coordinate (O2) at (3,0);
% B = r2(O1) : rotation -60 deg around O2(3,0). O1 rel O2 = (-3,0). Rot -60 -> (-1.5, 2.598). Abs (1.5, 2.598).
\coordinate (B) at (1.5,2.598);
% Points
\fill[popblue] (O1) circle (1.6pt) node[below left] {$O_1$};
\fill[poporange] (O2) circle (1.6pt) node[below right] {$O_2$};
\fill[poppurple] (B) circle (1.6pt) node[above] {$B = r_2(O_1)$};
\draw[popmuted, dashed] (O1) -- (O2);
\draw[poporange, dashed] (O2) -- (B);
\draw[poporange] (2.3,0.2) arc (180:120:0.7) node[midway, left, poporange] {$-60°$};
% Vecteur de translation
\draw[-{Stealth}, very thick, popgreen] (O1) -- (B) node[midway, above left, popgreen] {$\vec{u}$};
% Un point M et son image M''
\coordinate (M) at (-1.5,1.8);
\coordinate (Mpp) at (0,4.398); % M + u
\fill[popdark] (M) circle (1.4pt) node[left] {$M$};
\fill[popdark] (Mpp) circle (1.4pt) node[right] {$M''$};
\draw[-{Stealth}, thick, popgreen] (M) -- (Mpp) node[midway, above, popgreen] {$\vec{u}$};
% Trajectoire M -> M' -> M''
\coordinate (Mp) at (-2.30,-0.40); % approx r1(M) center O1 angle 60
\fill[popdark] (Mp) circle (1.2pt) node[below left] {$M'$};
\draw[popblue, -{Stealth}] (M) .. controls (-2.0,2.5) .. (Mp) node[midway, left, popblue] {$r_1$};
\draw[poporange, -{Stealth}] (Mp) .. controls (0.5,2.5) .. (Mpp) node[midway, above, poporange] {$r_2$};
\end{tikzpicture}
\end{center}
\legende{Cas $\alpha + \beta = 0$ (ici $60°$ et $-60°$) : la composée est une translation de vecteur $\vec{u} = \overrightarrow{O_1B}$. Le triangle $O_1O_2B$ est équilatéral, $\vec{u}$ a même longueur que $O_1O_2$ et fait $60°$ avec $\overrightarrow{O_1O_2}$.}
\end{popfigure}

\begin{attention}[Erreur fréquente au Probatoire]
Beaucoup d'élèves écrivent automatiquement : « la composée de deux rotations est une rotation d'angle $\alpha + \beta$ ». C'est \textbf{faux} quand $\alpha + \beta \equiv 0 \pmod{360°}$ : dans ce cas, la composée est une \textbf{translation}, et il n'y a \emph{pas de centre}. Le réflexe obligatoire est : \textbf{calculer d'abord $\alpha + \beta$ modulo $360°$}, puis seulement conclure. Le cas de la translation est explicitement au programme et tombe régulièrement au Probatoire technique. Autre piège : oublier que $r_1(O_1) = O_1$, ce qui simplifie énormément le calcul du vecteur.
\end{attention}

\courssub{Application concrète : déplacement d'une pièce rigide}

\begin{exemplebox}[Manivelle et volant en atelier]
Un technicien fait pivoter un support rigide autour d'un axe $O_1$ d'un angle de $120°$ (rotation $r_1$), puis la pièce solidaire du support pivote autour d'un second axe $O_2$ d'un angle de $-120°$ (rotation $r_2$). Comme $120° + (-120°) = 0$, le déplacement global de la pièce est une \textbf{translation} : la pièce a changé de position mais \emph{pas d'orientation} — une marque peinte sur la pièce pointe toujours dans la même direction. C'est exactement le principe utilisé dans certains mécanismes de manutention où l'on veut déplacer une charge sans la faire tourner (nacelles, plates-formes, parallélogrammes articulés). À l'inverse, avec deux rotations de $90°$ et $90°$ autour d'axes distincts, on obtient un demi-tour (symétrie centrale) : la pièce est retournée.
\end{exemplebox}

\begin{savaistu}[Tout déplacement est une rotation ou une translation]
Un théorème célèbre (dû à Chasles, mathématicien français du XIX\textsuperscript{e} siècle) affirme que \emph{tout} déplacement du plan qui conserve les orientations — c'est-à-dire toute manière de bouger un solide rigide sans le retourner comme une crêpe — est soit une \textbf{translation}, soit une \textbf{rotation}. La composée de deux rotations étudiée ici en est une illustration parfaite : quels que soient les centres et les angles, on retombe toujours sur l'un de ces deux types. C'est ce principe qui permet aux concepteurs de robots et de machines (fraiseuses numériques, bras de soudure dans les usines d'assemblage) de décrire n'importe quel mouvement complexe par une seule rotation ou une seule translation équivalente.
\end{savaistu}

\begin{exercice}[Pour s'entraîner]
Soit $O_1$ et $O_2$ deux points distincts. On note $r_1$ la rotation de centre $O_1$ et d'angle $45°$, et $r_2$ la rotation de centre $O_2$ et d'angle $135°$.
\begin{enumerate}
\item Déterminer la nature de $r_2 \circ r_1$ et préciser son angle.
\item Expliquer comment construire son centre $\Omega$.
\item Que se passerait-il si l'angle de $r_2$ était $-45°$ au lieu de $135°$ ?
\end{enumerate}
\end{exercice}

\begin{corrige}[Exercice 1]
\begin{enumerate}
\item $\alpha + \beta = 45° + 135° = 180° \not\equiv 0 \pmod{360°}$ : la composée est une rotation d'angle $180°$, c'est-à-dire une symétrie centrale.
\item On trace $[O_1O_2]$, puis en $O_1$ la demi-droite faisant un angle $\alpha/2 = 22{,}5°$ avec $(O_1O_2)$ (sens direct), et en $O_2$ la demi-droite faisant un angle $\beta/2 = 67{,}5°$ avec $(O_2O_1)$ (sens direct). Leur intersection est $\Omega$.
\item Si $\beta = -45°$, alors $\alpha + \beta = 0$ : la composée serait une translation, de vecteur $\overrightarrow{O_1\, r_2(O_1)}$.
\end{enumerate}
\end{corrige}

\begin{aretenir}[L'essentiel de la section]
\begin{itemize}
\item Composée de deux rotations de centres distincts $O_1$, $O_2$ et d'angles $\alpha$, $\beta$ : \textbf{on calcule d'abord $\alpha + \beta$ modulo $360°$}.
\item Si $\alpha + \beta \not\equiv 0 \pmod{360°}$ : c'est une \textbf{rotation d'angle $\alpha + \beta$}, dont le centre $\Omega$ s'obtient comme intersection des demi-droites faisant les angles $\alpha/2$ et $\beta/2$ avec $(O_1O_2)$ (selon la convention orientée du cours).
\item Si $\alpha + \beta \equiv 0 \pmod{360°}$ : c'est une \textbf{translation}, de vecteur $\overrightarrow{O_1\, r_2(O_1)}$ (car $r_1(O_1) = O_1$).
\item Cas remarquables : $90°$ et $90°$ donnent une symétrie centrale ; $60°$ et $-60°$ donnent une translation dont le vecteur se lit sur un triangle équilatéral.
\end{itemize}
\end{aretenir}

\coursec{Composée de deux homothéties de même centre}

Dans la section précédente, nous avons composé des rotations : nous avons vu que les angles s'\emph{additionnent}. Nous passons maintenant à un autre type de transformation fondamentale en atelier et en bureau d'études : l'\cle{homothétie}, c'est-à-dire l'agrandissement ou la réduction depuis un point fixe. Question naturelle : si l'on agrandit une figure, puis qu'on agrandit (ou réduit) le résultat \emph{depuis le même centre}, obtient-on encore une homothétie ? Et si oui, quel est son rapport ? C'est toute la question du réglage successif d'une photocopieuse, d'un zoom d'appareil photo ou de l'agrandissement d'un patron de couture en deux étapes.

\courssub{Situation et intuition}

\begin{definition}[Rappel : homothétie]
Soit $O$ un point du plan et $k$ un nombre réel non nul. L'\cle{homothétie} de centre $O$ et de rapport $k$, notée $h(O,k)$, est la transformation qui à tout point $M$ associe le point $M'$ tel que :
\[ \overrightarrow{OM'} = k\,\overrightarrow{OM}. \]
Si $k>0$, $M'$ est du même côté de $O$ que $M$ ; si $k<0$, $M'$ est de l'autre côté ; si $|k|>1$ c'est un agrandissement, si $0<|k|<1$ c'est une réduction.
\end{definition}

Considérons maintenant deux homothéties de \textbf{même centre} $O$ :
\begin{itemize}
\item $h_1 = h(O, k_1)$ qui envoie $M$ sur $M'$ : $\overrightarrow{OM'} = k_1\,\overrightarrow{OM}$ ;
\item $h_2 = h(O, k_2)$ qui envoie $M'$ sur $M''$ : $\overrightarrow{OM''} = k_2\,\overrightarrow{OM'}$.
\end{itemize}

La composée $h_2 \circ h_1$ envoie directement $M$ sur $M''$. Enchaînons les égalités vectorielles :
\[ \overrightarrow{OM''} = k_2\,\overrightarrow{OM'} = k_2\bigl(k_1\,\overrightarrow{OM}\bigr) = k_1 k_2\,\overrightarrow{OM}. \]
On reconnaît exactement la définition d'une homothétie de centre $O$ et de rapport $k_1 k_2$. L'intuition est donc limpide : \textbf{agrandir deux fois de suite revient à agrandir une seule fois, avec un rapport égal au produit des rapports}.

\courssub{Théorème fondamental}

\begin{propriete}[Composée de deux homothéties de même centre — démonstration exigible]
Soient $h_1 = h(O, k_1)$ et $h_2 = h(O, k_2)$ deux homothéties de \textbf{même centre} $O$ et de rapports non nuls $k_1$ et $k_2$. Alors la composée $h_2 \circ h_1$ est l'homothétie de centre $O$ et de rapport $k_1 \times k_2$ :
\[ h(O, k_2) \circ h(O, k_1) = h(O, k_1 k_2). \]
En particulier, si $k_1 k_2 = 1$, la composée est la \cle{transformation identique} : chaque point est envoyé sur lui-même. On dit alors que $h_1$ et $h_2$ sont des \cle{homothéties réciproques}.
\end{propriete}

\begin{experience}[Démonstration guidée du théorème]
On note $M$ un point quelconque du plan, $M' = h_1(M)$ et $M'' = h_2(M') = (h_2 \circ h_1)(M)$.
\begin{enumerate}
\item \textbf{Traduction vectorielle de $h_1$.} Par définition de $h_1 = h(O,k_1)$ : $\overrightarrow{OM'} = k_1\,\overrightarrow{OM}$.
\item \textbf{Traduction vectorielle de $h_2$.} Par définition de $h_2 = h(O,k_2)$ appliquée à $M'$ : $\overrightarrow{OM''} = k_2\,\overrightarrow{OM'}$.
\item \textbf{Substitution.} On remplace $\overrightarrow{OM'}$ par son expression :
\[ \overrightarrow{OM''} = k_2\,\overrightarrow{OM'} = k_2 \times k_1\,\overrightarrow{OM} = k_1 k_2\,\overrightarrow{OM}. \]
(Les vecteurs $\overrightarrow{OM}$, $\overrightarrow{OM'}$, $\overrightarrow{OM''}$ sont tous colinéaires, portés par la droite $(OM)$ : c'est ce qui rend le calcul si simple.)
\item \textbf{Conclusion.} L'égalité $\overrightarrow{OM''} = k_1 k_2\,\overrightarrow{OM}$ signifie exactement que $M''$ est l'image de $M$ par $h(O, k_1 k_2)$. Comme $M$ est quelconque, $h_2 \circ h_1 = h(O, k_1 k_2)$. $\blacksquare$
\item \textbf{Cas $k_1 k_2 = 1$.} Alors $\overrightarrow{OM''} = \overrightarrow{OM}$, donc $M'' = M$ pour tout point $M$ : c'est l'identité.
\end{enumerate}
\end{experience}

\begin{attention}[Le centre doit être le même]
Ce théorème n'est valable que si les deux homothéties ont le \textbf{même centre}. Si les centres sont distincts et $k_1 k_2 \neq 1$, la composée est encore une homothétie de rapport $k_1 k_2$ mais son centre n'est \emph{pas} $O$ : ce cas n'est pas au programme de Première technique. Retenez : même centre $\Rightarrow$ centre conservé et rapports multipliés.
\end{attention}

\courssub{Méthode et exemples chiffrés}

\begin{methode}[Déterminer la composée de deux homothéties de même centre]
\begin{enumerate}
\item Vérifier que les deux homothéties ont le \textbf{même centre} $O$.
\item \textbf{Multiplier} les rapports en tenant compte de leurs \textbf{signes} : $k = k_1 \times k_2$.
\item Conclure : $h_2 \circ h_1 = h(O, k_1 k_2)$.
\item Interpréter le résultat :
\begin{itemize}
\item si $k_1 k_2 = 1$ : c'est l'\textbf{identité} (on revient au point de départ) ;
\item si $k_1 k_2 = -1$ : c'est la \cle{symétrie centrale} de centre $O$ ;
\item si $k_1 k_2 > 0$ : homothétie de rapport positif (image du même côté de $O$) ;
\item si $k_1 k_2 < 0$ : homothétie de rapport négatif (image de l'autre côté de $O$).
\end{itemize}
\end{enumerate}
\end{methode}

\begin{exemplebox}[Une symétrie centrale cachée]
Soient $h_1 = h(O, 3)$ et $h_2 = h\!\left(O, -\dfrac{1}{3}\right)$. Déterminons $h_2 \circ h_1$.

Les deux homothéties ont le même centre $O$. Le rapport de la composée est :
\[ k = k_1 \times k_2 = 3 \times \left(-\frac{1}{3}\right) = -1. \]
Donc $h_2 \circ h_1 = h(O, -1)$, c'est-à-dire la \textbf{symétrie centrale de centre $O$}. Vérification vectorielle : $\overrightarrow{OM''} = -\dfrac{1}{3}\,\overrightarrow{OM'} = -\dfrac{1}{3}\times 3\,\overrightarrow{OM} = -\overrightarrow{OM}$, donc $O$ est le milieu de $[MM'']$.
\end{exemplebox}

\begin{exoresolu}[Retour à la figure initiale]
On considère un triangle $ABC$ et le point $O$. On applique $h_1 = h(O, 2)$ puis $h_2 = h\!\left(O, \dfrac{1}{2}\right)$. Déterminer la nature de $h_2 \circ h_1$ et interpréter géométriquement.
\tcblower
\textbf{Solution.} Les deux homothéties ont le même centre $O$, donc :
\[ h_2 \circ h_1 = h\!\left(O,\; 2 \times \frac{1}{2}\right) = h(O, 1). \]
Or $h(O,1)$ envoie tout point $M$ sur $M'$ tel que $\overrightarrow{OM'} = \overrightarrow{OM}$, c'est-à-dire $M' = M$ : c'est la \textbf{transformation identique}.

\textbf{Interprétation.} Le triangle $A'B'C'$, image de $ABC$ par $h_1$, a des dimensions doubles. Puis $h_2$ réduit $A'B'C'$ de moitié : on retombe exactement sur le triangle $ABC$ initial. Les homothéties $h(O,2)$ et $h\!\left(O,\tfrac12\right)$ sont \textbf{réciproques} : l'une « défait » ce que l'autre a fait. C'est le principe du zoom avant puis zoom arrière sur un plan de chantier numérique.
\end{exoresolu}

\begin{popfigure}
\begin{center}
\begin{tikzpicture}[scale=0.85, every node/.style={font=\small}]
\coordinate (O) at (0,0);
\coordinate (A) at (1.5,0.6);
\coordinate (B) at (2.4,2.0);
\coordinate (C) at (0.7,1.9);
\coordinate (A1) at (3,1.2);
\coordinate (B1) at (4.8,4.0);
\coordinate (C1) at (1.4,3.8);
% rayons depuis O
\draw[popmuted, dashed] (O) -- (A1);
\draw[popmuted, dashed] (O) -- (B1);
\draw[popmuted, dashed] (O) -- (C1);
% triangle initial
\draw[thick, popblue, fill=popblueL] (A) -- (B) -- (C) -- cycle;
% triangle intermédiaire
\draw[thick, poporange, fill=poporangeL, opacity=0.85] (A1) -- (B1) -- (C1) -- cycle;
% points
\fill[popdark] (O) circle (2pt) node[below left] {$O$};
\fill[popblue] (A) circle (1.8pt) node[below] {$A$};
\fill[popblue] (B) circle (1.8pt) node[right] {$B$};
\fill[popblue] (C) circle (1.8pt) node[left] {$C$};
\fill[poporange] (A1) circle (1.8pt) node[below right] {$A'$};
\fill[poporange] (B1) circle (1.8pt) node[right] {$B'$};
\fill[poporange] (C1) circle (1.8pt) node[above left] {$C'$};
% flèches de transformation
\draw[-{Stealth}, thick, popgreen] (2.1,3.1) arc[start angle=120, end angle=60, radius=1.1] node[midway, above] {$h_1 = h(O,2)$};
\draw[-{Stealth}, thick, poppurple] (4.6,2.6) arc[start angle=-30, end angle=-90, radius=1.6] node[midway, right] {$h_2 = h(O,\frac12)$};
\node[popdark, font=\small\itshape] at (2.6,-0.9) {$h_2 \circ h_1 = h(O,1) = \mathrm{Id}$ : retour au triangle initial};
\end{tikzpicture}
\end{center}
\legende{Le triangle $ABC$ (bleu) est agrandi par $h_1 = h(O,2)$ en $A'B'C'$ (orange), puis $h_2 = h(O,\tfrac12)$ ramène $A'B'C'$ sur $ABC$ : la composée est l'identité.}
\end{popfigure}

\courssub{Signe du rapport produit}

Le signe de $k_1 k_2$ se détermine par la règle des signes, et il a une signification géométrique précise :

\begin{center}
\begin{tabularx}{\linewidth}{|l|X|X|X|}
\hline
\rowcolor{popblueL} \textbf{Signes de $k_1$ et $k_2$} & \textbf{Signe de $k_1k_2$} & \textbf{Position de $M''$ par rapport à $O$} & \textbf{Exemple} \\
\hline
$k_1>0$ et $k_2>0$ & $+$ & même côté que $M$ & $h(O,2)\circ h(O,3)=h(O,6)$ \\
\hline
$k_1<0$ et $k_2<0$ & $+$ & même côté que $M$ & $h(O,-2)\circ h(O,-3)=h(O,6)$ \\
\hline
$k_1$ et $k_2$ de signes contraires & $-$ & côté opposé à $M$ & $h(O,3)\circ h(O,-\tfrac13)=h(O,-1)$ \\
\hline
\end{tabularx}
\end{center}

\begin{aretenir}[À retenir]
\begin{itemize}
\item Composée de deux homothéties de même centre : on \textbf{multiplie} les rapports, le centre est \textbf{conservé}.
\[ h(O,k_2) \circ h(O,k_1) = h(O, k_1 k_2) \]
\item Si $k_1 k_2 = 1$ : la composée est l'\textbf{identité} (homothéties réciproques).
\item Si $k_1 k_2 = -1$ : la composée est la \cle{symétrie centrale} de centre $O$.
\item Le signe de $k_1 k_2$ indique si l'image finale est du même côté de $O$ que le point initial ou du côté opposé.
\end{itemize}
\end{aretenir}

\begin{attention}[Erreur fréquente : additionner au lieu de multiplier]
L'erreur classique au Probatoire consiste à écrire $h(O,2) \circ h(O,3) = h(O, 5)$ en \textbf{additionnant} les rapports. C'est faux ! Avec les rotations, les \emph{angles} s'additionnent ; avec les homothéties, les \emph{rapports} se \textbf{multiplient}. Vérification concrète : agrandir une longueur de \SI{1}{cm} par $h(O,2)$ donne \SI{2}{cm}, puis par $h(O,3)$ donne \SI{6}{cm} : le rapport global est bien $2 \times 3 = 6$, et non $5$.
\end{attention}

\courssub{Effet sur les longueurs et les aires}

\begin{propriete}[Effet d'une homothétie sur les aires]
Une homothétie de rapport $k$ multiplie toutes les longueurs par $|k|$ et toutes les \textbf{aires} par $k^2$.
\end{propriete}

Par conséquent, pour la composée $h_2 \circ h_1 = h(O, k_1 k_2)$ :
\begin{itemize}
\item les longueurs sont multipliées par $|k_1 k_2| = |k_1| \times |k_2|$ ;
\item les aires sont multipliées par $(k_1 k_2)^2 = k_1^2 \times k_2^2$.
\end{itemize}

\begin{exemplebox}[Aire après deux agrandissements]
Un atelier de couture agrandit un patron d'aire \SI{200}{cm^2} par $h_1 = h(O, 3)$, puis le résultat par $h_2 = h(O, 2)$. L'aire finale est :
\[ \mathcal{A} = 200 \times 3^2 \times 2^2 = 200 \times 9 \times 4 = \SI{7200}{cm^2}. \]
On retrouve le même résultat directement : le rapport global est $k = 6$, donc l'aire est multipliée par $6^2 = 36$, et $200 \times 36 = \SI{7200}{cm^2}$.
\end{exemplebox}

\courssub{Application concrète : agrandissements successifs}

\begin{exoresolu}[La photocopieuse de l'atelier]
Un technicien photocopie un plan en réglant la machine sur \SI{150}{\percent}, puis il photocopie le résultat en réglant sur \SI{80}{\percent}.
\begin{enumerate}
\item Modéliser chaque étape par une homothétie et déterminer la transformation globale.
\item Le plan final est-il un agrandissement ou une réduction du plan initial ?
\item Le plan initial a une aire de \SI{600}{cm^2}. Quelle est l'aire du plan final ?
\end{enumerate}
\tcblower
\textbf{Solution.}
\begin{enumerate}
\item Un réglage à \SI{150}{\percent} correspond à $h_1 = h(O, \num{1,5})$ ; un réglage à \SI{80}{\percent} correspond à $h_2 = h(O, \num{0,8})$. Les deux agrandissements se font depuis le même coin de la feuille (même centre $O$), donc :
\[ h_2 \circ h_1 = h(O,\; \num{1,5} \times \num{0,8}) = h(O, \num{1,2}). \]
\item Le rapport global est $\num{1,2} > 1$ : le plan final est un \textbf{agrandissement de \SI{120}{\percent}} du plan initial. Les longueurs ont été multipliées par $\num{1,2}$.
\item L'aire est multipliée par $(\num{1,2})^2 = \num{1,44}$ :
\[ \mathcal{A} = 600 \times \num{1,44} = \SI{864}{cm^2}. \]
\end{enumerate}
\end{exoresolu}

\begin{popfigure}
\begin{center}
\begin{tikzpicture}[every node/.style={font=\small}]
% photo initiale
\draw[thick, popblue, fill=popblueL] (0,0) rectangle (2,1.4);
\node[popblue] at (1,0.7) {Plan initial};
\node[popblue, below] at (1,0) {\SI{100}{\percent}};
% flèche 1
\draw[-{Stealth}, very thick, poporange] (2.3,0.7) -- (3.7,0.7) node[midway, above] {$\times \num{1,5}$};
% photo intermédiaire
\draw[thick, poporange, fill=poporangeL] (4,-0.35) rectangle (7,1.75);
\node[poporange] at (5.5,0.7) {Copie intermédiaire};
\node[poporange, below] at (5.5,-0.35) {\SI{150}{\percent}};
% flèche 2
\draw[-{Stealth}, very thick, popgreen] (7.3,0.7) -- (8.7,0.7) node[midway, above] {$\times \num{0,8}$};
% photo finale
\draw[thick, popgreen, fill=popgreenL] (9,-0.14) rectangle (11.4,1.54);
\node[popgreen] at (10.2,0.7) {Copie finale};
\node[popgreen, below] at (10.2,-0.14) {\SI{120}{\percent}};
% bilan
\node[popdark, font=\small\itshape] at (5.7,-1.3) {$\num{1,5} \times \num{0,8} = \num{1,2}$ : au final, un agrandissement de \SI{120}{\percent}};
\end{tikzpicture}
\end{center}
\legende{Agrandissements successifs d'un plan à la photocopieuse : les rapports se multiplient ($\num{1,5} \times \num{0,8} = \num{1,2}$), ils ne s'additionnent pas.}
\end{popfigure}

\begin{savaistu}[Le piège de la photocopieuse]
Beaucoup d'utilisateurs pensent qu'une photocopie à \SI{150}{\percent} suivie d'une photocopie à \SI{80}{\percent} « compense » presque et redonne à peu près la taille initiale, car $150 - 80 = 70$... Erreur ! Les pourcentages d'agrandissement se \textbf{multiplient} : $\num{1,5} \times \num{0,8} = \num{1,2}$, soit un agrandissement final de \SI{120}{\percent}. Pour revenir exactement à la taille initiale après un réglage à \SI{150}{\percent}, il faudrait régler la machine sur $\dfrac{100}{\num{1,5}} \approx \SI{66,7}{\percent}$, car $\num{1,5} \times \dfrac{2}{3} = 1$ : les deux réglages correspondent alors à des homothéties \textbf{réciproques}. Ce calcul de pourcentages successifs est partout : remises commerciales successives au marché, taux d'intérêt composés à la banque, zooms d'appareils photo.
\end{savaistu}

\begin{exercice}[À faire]
\begin{enumerate}
\item Déterminer la nature et les éléments caractéristiques de $h_2 \circ h_1$ dans chaque cas :
\begin{enumerate}
\item $h_1 = h(O, 4)$ et $h_2 = h\!\left(O, \dfrac{1}{4}\right)$ ;
\item $h_1 = h(O, -2)$ et $h_2 = h(O, -5)$ ;
\item $h_1 = h\!\left(O, \dfrac{1}{2}\right)$ et $h_2 = h(O, -2)$.
\end{enumerate}
\item Un négatif photo de format \SI{24}{mm} $\times$ \SI{36}{mm} est agrandi par $h(O, 5)$, puis le tirage obtenu est reproduit par $h(O, \num{0,4})$. Quelles sont les dimensions du document final ? Quel est le rapport global ?
\end{enumerate}
\end{exercice}

\begin{corrige}[Exercice]
\begin{enumerate}
\item \begin{enumerate}
\item $k = 4 \times \dfrac{1}{4} = 1$ : la composée est la \textbf{transformation identique} (homothéties réciproques).
\item $k = (-2) \times (-5) = 10$ : la composée est $h(O, 10)$, homothétie de rapport positif (l'image est du même côté de $O$).
\item $k = \dfrac{1}{2} \times (-2) = -1$ : la composée est $h(O,-1)$, c'est-à-dire la \textbf{symétrie centrale} de centre $O$.
\end{enumerate}
\item Rapport global : $k = 5 \times \num{0,4} = 2$. Dimensions finales : $24 \times 2 = \SI{48}{mm}$ et $36 \times 2 = \SI{72}{mm}$. Le document final mesure donc $48 \times 72$ (en mm), soit le double du négatif initial.
\end{enumerate}
\end{corrige}

En résumé, la composée de deux homothéties de même centre est d'une grande simplicité : le centre ne change pas et les rapports se multiplient. La situation se complique lorsque l'on mélange les types de transformations : que se passe-t-il lorsqu'on compose une homothétie avec une translation ? C'est l'objet de la section suivante.

\coursec{Composée d'une homothétie et d'une translation}

Dans la section précédente, nous avons étudié la composée de deux homothéties de même centre $O$ : nous avons vu qu'on obtient encore une homothétie de centre $O$, dont le rapport est le produit des rapports. Mais que se passe-t-il lorsqu'on compose une homothétie avec une transformation d'un autre type, la \cle{translation} ? C'est exactement la situation qu'on rencontre dans les métiers du textile, de l'imprimerie ou de la signalétique : on agrandit un motif, \emph{puis} on le déplace pour le placer au bon endroit. Le résultat est surprenant : on obtient encore une homothétie, mais son centre a \emph{déménagé}. C'est l'objet de cette section.

\courssub{Situation et mise en place du problème}

\begin{definition}[Composée d'une homothétie et d'une translation]
Soit $h$ l'homothétie de centre $O$ et de rapport $k$ ($k \neq 0$), et $t$ la translation de vecteur $\vec{u}$. On peut former deux composées :
\begin{itemize}
\item $f = t \circ h$ : on applique \textbf{d'abord} l'homothétie, \textbf{puis} la translation ;
\item $g = h \circ t$ : on applique \textbf{d'abord} la translation, \textbf{puis} l'homothétie.
\end{itemize}
Pour tout point $M$, on note $M' = h(M)$ puis $M'' = t(M')$ dans le premier cas, et $M' = t(M)$ puis $M'' = h(M')$ dans le second.
\end{definition}

Traduisons ces composées avec des vecteurs. Rappelons que $M' = h(M)$ signifie $\overrightarrow{OM'} = k\,\overrightarrow{OM}$, et que $M'' = t(M')$ signifie $\overrightarrow{M'M''} = \vec{u}$.

Pour $f = t \circ h$, l'image $M''$ de $M$ vérifie :
\[ \overrightarrow{OM''} = \overrightarrow{OM'} + \overrightarrow{M'M''} = k\,\overrightarrow{OM} + \vec{u}. \]

Pour $g = h \circ t$, avec $M' = t(M)$ :
\[ \overrightarrow{OM''} = k\,\overrightarrow{OM'} = k\left(\overrightarrow{OM} + \vec{u}\right) = k\,\overrightarrow{OM} + k\vec{u}. \]

On observe déjà deux choses importantes : les deux composées \textbf{ne donnent pas le même résultat} (les vecteurs ajoutés sont $\vec{u}$ et $k\vec{u}$, différents en général), et dans les deux cas l'image est obtenue en multipliant par $k$ puis en ajoutant un vecteur fixe. C'est exactement la forme vectorielle d'une homothétie de rapport $k$... à condition de trouver le bon centre.

\courssub{Théorème fondamental}

\begin{propriete}[Théorème (admis) : composée d'une homothétie et d'une translation]
Soit $h$ l'homothétie de centre $O$ et de rapport $k$, et $t$ la translation de vecteur $\vec{u}$.
\begin{itemize}
\item Si $k \neq 1$, alors $t \circ h$ et $h \circ t$ sont des \textbf{homothéties de rapport $k$}, dont les centres sont en général \textbf{différents de $O$} (et différents l'un de l'autre).
\item Si $k = 1$, alors $h$ est l'identité du plan, et les deux composées sont égales à la \textbf{translation} $t$.
\end{itemize}
Ce théorème est admis au Probatoire technique ; sa justification vectorielle est donnée ci-dessous à titre formateur.
\end{propriete}

\begin{experience}[Pourquoi obtient-on une homothétie ? Justification vectorielle guidée]
Étudions $f = t \circ h$ avec $k \neq 1$. Cherchons un point $\Omega$ qui serait \textbf{invariant} (fixe) par $f$, c'est-à-dire tel que $f(\Omega) = \Omega$.
\begin{enumerate}
\item La condition $f(\Omega) = \Omega$ s'écrit vectoriellement : $\overrightarrow{O\Omega} = k\,\overrightarrow{O\Omega} + \vec{u}$.
\item On isole : $(1 - k)\,\overrightarrow{O\Omega} = \vec{u}$, d'où, puisque $k \neq 1$ :
\[ \overrightarrow{O\Omega} = \frac{1}{1 - k}\,\vec{u}. \]
Ce vecteur est parfaitement défini : il existe donc un \textbf{unique} point fixe $\Omega$.
\item Montrons maintenant que $f$ est l'homothétie de centre $\Omega$ et de rapport $k$. Pour tout point $M$ d'image $M'' = f(M)$ :
\begin{align*}
\overrightarrow{\Omega M''} &= \overrightarrow{\Omega O} + \overrightarrow{OM''} = \overrightarrow{\Omega O} + k\,\overrightarrow{OM} + \vec{u} \\
&= \overrightarrow{\Omega O} + k\left(\overrightarrow{O\Omega} + \overrightarrow{\Omega M}\right) + \vec{u} \\
&= (1 - k)\,\overrightarrow{\Omega O} + \vec{u} + k\,\overrightarrow{\Omega M}.
\end{align*}
Or, d'après l'étape 2, $(1-k)\,\overrightarrow{O\Omega} = \vec{u}$, donc $(1-k)\,\overrightarrow{\Omega O} + \vec{u} = \vec{0}$. Il reste :
\[ \overrightarrow{\Omega M''} = k\,\overrightarrow{\Omega M}. \]
C'est exactement la définition de l'homothétie de centre $\Omega$ et de rapport $k$.
\end{enumerate}
Le même raisonnement s'applique à $h \circ t$ : le point fixe $\Omega'$ vérifie $\overrightarrow{O\Omega'} = \dfrac{k}{1-k}\,\vec{u}$, différent de $\Omega$ en général.
\end{experience}

\begin{aretenir}[Ce qu'il faut retenir]
Composer une homothétie de rapport $k \neq 1$ avec une translation donne \textbf{toujours une homothétie de même rapport $k$}. La translation ne change pas le rapport, elle \textbf{déplace le centre}. Le nouveau centre est l'unique point fixe de la composée.
\end{aretenir}

\courssub{Détermination du nouveau centre : la méthode du point fixe}

\begin{methode}[Trouver le centre de l'homothétie composée]
Pour déterminer le centre $\Omega$ de $f = t \circ h$ (ou de $g = h \circ t$) :
\begin{enumerate}
\item On se place dans un repère et on note $\Omega(x\,;\,y)$ le point cherché.
\item On écrit la condition de point fixe : $\Omega = f(\Omega)$, c'est-à-dire $\Omega = t\big(h(\Omega)\big)$.
\item On traduit cette égalité par deux équations sur les coordonnées (une pour $x$, une pour $y$).
\item On résout le système obtenu : la solution unique (car $k \neq 1$) donne les coordonnées de $\Omega$.
\item On conclut : la composée est l'homothétie de centre $\Omega$ et de rapport $k$.
\end{enumerate}
On peut aussi travailler directement avec les vecteurs : pour $t \circ h$, $\overrightarrow{O\Omega} = \dfrac{1}{1-k}\,\vec{u}$ ; pour $h \circ t$, $\overrightarrow{O\Omega'} = \dfrac{k}{1-k}\,\vec{u}$.
\end{methode}

\begin{exoresolu}[Centre de $t \circ h$ avec $h(O\,;\,2)$ et $\vec{u}(3\,;\,0)$]
\textbf{Énoncé.} Dans un repère orthonormé, $O$ est l'origine, $h$ est l'homothétie de centre $O$ et de rapport $2$, et $t$ est la translation de vecteur $\vec{u}(3\,;\,0)$. Déterminer la nature et les éléments caractéristiques de $f = t \circ h$.
\tcblower
\textbf{Solution détaillée.}

\emph{Étape 1 : nature de $f$.} Le rapport est $k = 2 \neq 1$ : d'après le théorème, $f$ est une homothétie de rapport $2$, de centre $\Omega$ à déterminer.

\emph{Étape 2 : recherche du point fixe.} Soit $\Omega(x\,;\,y)$. Son image par $h$ est $\Omega_1 = h(\Omega)$, avec $\overrightarrow{O\Omega_1} = 2\,\overrightarrow{O\Omega}$, donc $\Omega_1(2x\,;\,2y)$. Son image par $f$ est $\Omega_2 = t(\Omega_1)$, de coordonnées $(2x + 3\,;\,2y)$.

La condition $\Omega = f(\Omega)$ s'écrit :
\[
\begin{cases}
x = 2x + 3 \\
y = 2y
\end{cases}
\quad\Longleftrightarrow\quad
\begin{cases}
-x = 3 \\
-y = 0
\end{cases}
\quad\Longleftrightarrow\quad
\begin{cases}
x = -3 \\
y = 0
\end{cases}
\]

\emph{Étape 3 : conclusion.} Le point fixe est $\Omega(-3\,;\,0)$. Donc $f = t \circ h$ est l'\textbf{homothétie de centre $\Omega(-3\,;\,0)$ et de rapport $2$}.

\emph{Vérification vectorielle :} $\overrightarrow{O\Omega} = \dfrac{1}{1-2}\,\vec{u} = -\vec{u}$, donc $\Omega = O - \vec{u}$, c'est-à-dire $\Omega(-3\,;\,0)$. Le résultat est confirmé.
\end{exoresolu}

\begin{attention}[Erreur fréquente : « le centre reste $O$ »]
Beaucoup d'élèves affirment que la composée $t \circ h$ est l'homothétie de centre $O$ et de rapport $k$. C'est \textbf{faux} : la translation \textbf{déplace le centre}. Dans l'exemple ci-dessus, le centre est passé de $O(0\,;\,0)$ à $\Omega(-3\,;\,0)$. Le seul cas où le centre ne bouge pas est celui où $\vec{u} = \vec{0}$ (translation nulle). Autre piège : oublier que le centre de $t \circ h$ et celui de $h \circ t$ ne sont \textbf{pas les mêmes} en général.
\end{attention}

\courssub{Comparaison de $t \circ h$ et $h \circ t$}

Reprenons l'exemple précédent et calculons le centre $\Omega'$ de $g = h \circ t$. Pour $\Omega'(x\,;\,y)$, on a d'abord $t(\Omega') = (x+3\,;\,y)$, puis $h\big(t(\Omega')\big) = \big(2(x+3)\,;\,2y\big)$. La condition de point fixe donne :
\[
\begin{cases}
x = 2x + 6 \\
y = 2y
\end{cases}
\quad\Longleftrightarrow\quad
\begin{cases}
x = -6 \\
y = 0
\end{cases}
\]
Donc $g = h \circ t$ est l'homothétie de centre $\Omega'(-6\,;\,0)$ et de rapport $2$.

\begin{propriete}[Comparaison des deux composées]
$t \circ h$ et $h \circ t$ sont deux homothéties de \textbf{même rapport $k$}, mais de \textbf{centres différents en général} : la composition d'une homothétie et d'une translation \textbf{n'est pas commutative}. Avec les notations précédentes :
\[ \overrightarrow{O\Omega} = \frac{1}{1-k}\,\vec{u} \qquad \text{et} \qquad \overrightarrow{O\Omega'} = \frac{k}{1-k}\,\vec{u}. \]
\end{propriete}

\begin{exemplebox}[Illustration numérique de la non-commutativité]
Appliquons les deux composées au point $A(1\,;\,0)$ avec $h(O\,;\,2)$ et $\vec{u}(3\,;\,0)$ :
\begin{itemize}
\item Par $f = t \circ h$ : $h(A) = (2\,;\,0)$, puis $t\big(h(A)\big) = (5\,;\,0)$. Vérification avec le centre $\Omega(-3\,;\,0)$ : $\overrightarrow{\Omega A} = (4\,;\,0)$, et $2\,\overrightarrow{\Omega A} = (8\,;\,0)$, d'où l'image $\Omega + (8\,;\,0) = (5\,;\,0)$. Correct.
\item Par $g = h \circ t$ : $t(A) = (4\,;\,0)$, puis $h\big(t(A)\big) = (8\,;\,0)$. Vérification avec $\Omega'(-6\,;\,0)$ : $2\,\overrightarrow{\Omega' A} = 2\,(7\,;\,0) = (14\,;\,0)$, d'où l'image $(8\,;\,0)$. Correct.
\end{itemize}
Les deux images de $A$ sont distinctes : $(5\,;\,0) \neq (8\,;\,0)$. L'ordre des transformations compte.
\end{exemplebox}

\begin{popfigure}
\begin{center}
\begin{tikzpicture}[scale=0.85,>=Stealth]
% Repère
\draw[popline,->] (-4.5,0) -- (6.5,0) node[below right] {$x$};
\draw[popline,->] (0,-1.5) -- (0,3.5) node[above left] {$y$};
% Points O et Omega
\fill[popdark] (0,0) circle (2pt) node[below left] {$O$};
\fill[poporange] (-3,0) circle (2pt) node[below left] {$\Omega(-3\,;\,0)$};
% Point A et ses images
\fill[popblue] (1,0.8) circle (2pt) node[above left] {$A(1\,;\,0{,}8)$};
\fill[popblue!60] (2,1.6) circle (2pt) node[above] {$A_1 = h(A)$};
\fill[poppurple] (5,1.6) circle (2pt) node[above right] {$A_2 = t(A_1) = f(A)$};
% Flèches
\draw[->,popblue,thick] (1,0.8) -- (2,1.6) node[midway,above,sloped] {\small $h(O\,;\,2)$};
\draw[->,poppurple,thick] (2,1.6) -- (5,1.6) node[midway,above] {\small $t$, $\vec{u}(3\,;\,0)$};
% Vecteur u depuis O
\draw[->,popgreen,thick] (0,-0.6) -- (3,-0.6) node[midway,below] {$\vec{u}$};
% Relation Omega -> A2
\draw[poporange,dashed,thick] (-3,0) -- (5,1.6);
\node[poporange] at (1.6,-1.1) {\small $\overrightarrow{\Omega A_2} = 2\,\overrightarrow{\Omega A}$};
\draw[poporange,dashed] (-3,0) -- (1,0.8);
\end{tikzpicture}
\end{center}
\legende{L'homothétie $h(O\,;\,2)$ suivie de la translation de vecteur $\vec{u}(3\,;\,0)$ : la composée est l'homothétie de centre $\Omega(-3\,;\,0)$, distinct de $O$, et de rapport $2$.}
\end{popfigure}

\begin{popfigure}
\begin{center}
\begin{tikzpicture}[scale=0.7,>=Stealth]
% ---- Gauche : t o h ----
\begin{scope}
\draw[popline,->] (-2.5,0) -- (7,0);
\draw[popline,->] (0,-0.5) -- (0,5);
\fill[popdark] (0,0) circle (1.5pt) node[below left] {\small $O$};
\fill[poporange] (-2,0) circle (1.5pt) node[below] {\small $\Omega$};
% Triangle ABC
\coordinate (A) at (1,1); \coordinate (B) at (2,1); \coordinate (C) at (1.5,2.5);
\draw[popblue,thick] (A) -- (B) -- (C) -- cycle;
\node[popblue] at (1.5,0.5) {\small $ABC$};
% Images par h (rapport 2)
\coordinate (A1) at (2,2); \coordinate (B1) at (4,2); \coordinate (C1) at (3,5);
\draw[popblue!50,thick,dashed] (A1) -- (B1) -- (C1) -- cycle;
\node[popblue!70] at (3.5,4.5) {\small $h(ABC)$};
% Images par t o h (vecteur (2,0))
\coordinate (A2) at (4,2); \coordinate (B2) at (6,2); \coordinate (C2) at (5,5);
\draw[poppurple,very thick] (A2) -- (B2) -- (C2) -- cycle;
\node[poppurple] at (5.5,4.8) {\small $t\circ h(ABC)$};
% Centre Omega = -2,0 (car u=(2,0), k=2 -> 1/(1-2)*u = -u)
\draw[poporange,dashed] (-2,0) -- (A) node[midway,left] {\small $\overrightarrow{\Omega A}$};
\draw[poporange,dashed] (-2,0) -- (A2) node[midway,right] {\small $\overrightarrow{\Omega A_2}=2\overrightarrow{\Omega A}$};
\node at (2.5,-1) {\small \textbf{$t \circ h$ : homothétie de centre $\Omega(-2,0)$}};
\end{scope}
% ---- Droite : h o t ----
\begin{scope}[xshift=11cm]
\draw[popline,->] (-4.5,0) -- (5,0);
\draw[popline,->] (0,-0.5) -- (0,5);
\fill[popdark] (0,0) circle (1.5pt) node[below left] {\small $O$};
\fill[popgreen] (-4,0) circle (1.5pt) node[below] {\small $\Omega'$};
\coordinate (A) at (1,1); \coordinate (B) at (2,1); \coordinate (C) at (1.5,2.5);
\draw[popblue,thick] (A) -- (B) -- (C) -- cycle;
\node[popblue] at (1.5,0.5) {\small $ABC$};
% Images par t (vecteur (2,0))
\coordinate (A1) at (3,1); \coordinate (B1) at (4,1); \coordinate (C1) at (3.5,2.5);
\draw[popblue!50,thick,dashed] (A1) -- (B1) -- (C1) -- cycle;
\node[popblue!70] at (3.5,2.8) {\small $t(ABC)$};
% Images par h o t
\coordinate (A2) at (6,2); \coordinate (B2) at (8,2); \coordinate (C2) at (7,5);
\draw[popgreen!80!black,very thick] (A2) -- (B2) -- (C2) -- cycle;
\node[popgreen!80!black] at (7.5,4.8) {\small $h\circ t(ABC)$};
% Centre Omega' = -4,0 (car k/(1-k)*u = 2/-1*(2,0) = -4,0)
\draw[popgreen!80!black,dashed] (-4,0) -- (A) node[midway,left] {\small $\overrightarrow{\Omega' A}$};
\draw[popgreen!80!black,dashed] (-4,0) -- (A2) node[midway,right] {\small $\overrightarrow{\Omega' A_2}=2\overrightarrow{\Omega' A}$};
\node at (0.5,-1) {\small \textbf{$h \circ t$ : homothétie de centre $\Omega'(-4,0)$}};
\end{scope}
\end{tikzpicture}
\end{center}
\legende{Comparaison sur un même triangle $ABC$ ($k=2$, $\vec{u}=(2,0)$) : à gauche $t \circ h$ (centre $\Omega(-2,0)$), à droite $h \circ t$ (centre $\Omega'(-4,0)$). Les deux images sont des triangles deux fois plus grands (même rapport $k=2$), mais placés différemment car les centres sont distincts.}
\end{popfigure}

\courssub{Cas particulier $k = 1$ et applications concrètes}

\begin{propriete}[Cas $k = 1$]
Si $k = 1$, l'homothétie $h(O\,;\,1)$ est l'\textbf{identité} du plan (chaque point est sa propre image). Alors :
\[ t \circ h = h \circ t = t. \]
La composée est simplement la translation $t$. Dans ce cas, il n'y a pas de nouveau centre à chercher.
\end{propriete}

\begin{exemplebox}[Application : redimensionnement et placement d'un motif sur un tissu]
Une couturière de Yaoundé veut reproduire un motif losangique sur un pagne. Sur sa feuille de patron, le motif de référence est dessiné dans un repère, et elle veut obtenir un motif \textbf{deux fois plus grand}, \textbf{décalé de \SI{12}{cm}} vers la droite pour poursuivre la frise.

Situation mathématique : $h$ est l'homothétie de centre $O$ (coin du patron) et de rapport $2$, $t$ est la translation de vecteur $\vec{u}(12\,;\,0)$ (en cm). La transformation à appliquer est $f = t \circ h$.

D'après le théorème, $f$ est une homothétie de rapport $2$ dont le centre $\Omega$ vérifie $\overrightarrow{O\Omega} = \dfrac{1}{1-2}\,\vec{u} = -\vec{u}$, soit $\Omega(-12\,;\,0)$. Concrètement : au lieu d'agrandir puis de déplacer en deux opérations, la couturière peut tracer directement le motif agrandi \emph{en visant depuis le point $\Omega$ situé à \SI{12}{cm} à gauche du coin du patron} : chaque point du nouveau motif est sur la demi-droite $[\Omega M)$, à distance double. Un seul point de visée suffit, ce qui limite les erreurs de report.
\end{exemplebox}

\begin{exemplebox}[Application : placement d'un logo sur une affiche]
Un imprimeur de Douala doit placer le logo d'une entreprise sur une affiche : le logo doit être agrandi d'un facteur $1{,}5$ puis déplacé vers le coin supérieur droit, translation de vecteur $\vec{u}(20\,;\,30)$ (en mm). La transformation est $f = t \circ h$ avec $k = 1{,}5 \neq 1$ : c'est une homothétie de rapport $1{,}5$ dont le centre vérifie $\overrightarrow{O\Omega} = \dfrac{1}{1 - 1{,}5}\,\vec{u} = -2\,\vec{u}$, soit $\Omega(-40\,;\,-60)$. Le logiciel de PAO effectue exactement ce calcul pour afficher le logo à la bonne position en une seule opération.
\end{exemplebox}

\begin{savaistu}[Frises et pavages : les maths des pagnes et des façades]
Les frises qui ornent les pagnes traditionnels (ndop, kaba ngondo), les façades des cases et les carrelages combinent en permanence des \textbf{homothéties} (pour faire varier la taille des motifs) et des \textbf{translations} (pour les répéter le long d'une bande). En mathématiques, on dit que ces frises sont invariantes par un groupe de transformations. Le théorème de cette section explique pourquoi un motif agrandi puis décalé reste parfaitement « aligné » avec les autres : la composée reste une homothétie, qui conserve les directions, les angles et les proportions. C'est aussi le principe utilisé par les machines à broder numériques et les logiciels de dessin industriel.
\end{savaistu}

\begin{exercice}[À faire]
Dans un repère orthonormé, on considère l'homothétie $h$ de centre $O$ et de rapport $-2$, et la translation $t$ de vecteur $\vec{u}(6\,;\,3)$.
\begin{enumerate}
\item Quelle est la nature de $f = t \circ h$ ? Préciser son rapport.
\item Déterminer les coordonnées du centre $\Omega$ de $f$.
\item Déterminer les coordonnées du centre $\Omega'$ de $g = h \circ t$. Les centres $\Omega$ et $\Omega'$ sont-ils confondus ?
\item Calculer l'image du point $A(1\,;\,2)$ par $f$ de deux façons différentes.
\end{enumerate}
\end{exercice}

\begin{corrige}[Exercice]
\begin{enumerate}
\item Le rapport est $k = -2 \neq 1$ : d'après le théorème, $f$ est une \textbf{homothétie de rapport $-2$}.
\item Soit $\Omega(x\,;\,y)$. On a $h(\Omega) = (-2x\,;\,-2y)$ puis $f(\Omega) = (-2x + 6\,;\,-2y + 3)$. La condition $\Omega = f(\Omega)$ donne :
\[
\begin{cases}
x = -2x + 6 \\
y = -2y + 3
\end{cases}
\Longleftrightarrow
\begin{cases}
3x = 6 \\
3y = 3
\end{cases}
\Longleftrightarrow
\begin{cases}
x = 2 \\
y = 1
\end{cases}
\]
Donc $\Omega(2\,;\,1)$. Vérification vectorielle : $\overrightarrow{O\Omega} = \dfrac{1}{1-(-2)}\,\vec{u} = \dfrac{1}{3}\vec{u} = (2\,;\,1)$. Correct.
\item Pour $g = h \circ t$ : $t(\Omega') = (x+6\,;\,y+3)$ puis $g(\Omega') = (-2x-12\,;\,-2y-6)$. La condition de point fixe donne $3x = -12$ et $3y = -6$, soit $\Omega'(-4\,;\,-2)$. Les centres $\Omega(2\,;\,1)$ et $\Omega'(-4\,;\,-2)$ sont \textbf{distincts}.
\item Première façon : $h(A) = (-2\,;\,-4)$, puis $f(A) = t\big(h(A)\big) = (4\,;\,-1)$. Deuxième façon : $\overrightarrow{\Omega A} = (1-2\,;\,2-1) = (-1\,;\,1)$, donc $\overrightarrow{\Omega f(A)} = -2\,\overrightarrow{\Omega A} = (2\,;\,-2)$, d'où $f(A) = \Omega + (2\,;\,-2) = (4\,;\,-1)$. Les deux méthodes concordent.
\end{enumerate}
\end{corrige}

\coursec{Synthèse, méthodes et applications concrètes}

Nous avons étudié successivement la composée de deux rotations (même centre puis centres distincts), la composée de deux homothéties de même centre, et la composée d'une homothétie et d'une translation. Toutes ces transformations font partie de la grande famille des \textbf{similitudes directes} : elles conservent les angles, l'alignement et le parallélisme, et transforment une droite en une droite. Cette section synthétise les résultats, donne une méthode unique pour résoudre tout exercice de composée, illustre la démarche sur des problèmes techniques concrets, et ouvre une perspective vers l'écriture complexe qui unifiera tout cela en Terminale.

\courssub{Tableau récapitulatif : nature de la composée}

Le tableau suivant résume l'ensemble des cas rencontrés dans cette leçon. C'est votre « arbre de décision » : devant un exercice, identifiez les deux transformations, relevez leurs paramètres, et lisez la nature de la composée dans la bonne ligne.

\begin{popfigure}
\centering
\begin{tikzpicture}[
    node distance=1.0cm and 1.2cm,
    decision/.style={diamond, draw=popblue, thick, fill=popblue!10, text width=3.8cm, align=center, inner sep=4pt},
    process/.style={rectangle, draw=popdark, thick, fill=popcream, text width=4.5cm, align=center, inner sep=4pt, rounded corners=4pt},
    start/.style={ellipse, draw=popgreen, thick, fill=popgreen!15, text width=3cm, align=center, inner sep=4pt},
    arrow/.style={-Stealth, thick, popdark},
    labelstyle/.style={font=\small, text=popdark}
]
% Nœuds
\node[start] (start) {Deux transformations $T_2 \circ T_1$};
\node[decision, below=of start] (type1) {$T_1$ et $T_2$ sont-elles deux rotations ?};
\node[decision, below=of type1] (type2) {$T_1$ et $T_2$ sont-elles deux homothéties de même centre ?};
\node[decision, below=of type2] (type3) {L'une est homothétie, l'autre translation ?};
% Branche Rotations
\node[process, left=of type1, xshift=-2.5cm] (rot1) {\textbf{Rotations même centre $\Omega$}\newline Angles $\alpha_1, \alpha_2$\newline $\Rightarrow$ Rotation centre $\Omega$, angle $\alpha_1+\alpha_2$};
\node[process, right=of type1, xshift=2.5cm] (rot2) {\textbf{Rotations centres distincts $\Omega_1 \neq \Omega_2$}\newline Angles $\alpha_1, \alpha_2$\newline $\alpha_1+\alpha_2 \not\equiv 0 [2\pi]$ $\Rightarrow$ Rotation angle $\alpha_1+\alpha_2$\newline Centre : intersection de deux médiatrices (ou composition de symétries)\newline $\alpha_1+\alpha_2 \equiv 0 [2\pi]$ $\Rightarrow$ Translation\newline Vecteur $\vec{v} = \overrightarrow{\Omega_1\Omega_2} - R_{\alpha_1}(\overrightarrow{\Omega_1\Omega_2})$};
% Branche Homothéties
\node[process, left=of type2, xshift=-2.5cm] (hom1) {\textbf{Homothéties même centre $\Omega$}\newline Rapports $k_1, k_2$\newline $\Rightarrow$ Homothétie centre $\Omega$, rapport $k_1 k_2$};
% Branche Homothétie + Translation
\node[process, left=of type3, xshift=-2.5cm] (ht1) {\textbf{Homothétie $H(\Omega,k)$ + Translation $T_{\vec{v}}$}\newline ($k \neq 1$) $\Rightarrow$ Homothétie rapport $k$\newline Ordre $T \circ H$ : Centre $\Omega' = \Omega - \frac{\vec{v}}{k-1}$\newline Ordre $H \circ T$ : Centre $\Omega'' = \Omega - \frac{k\vec{v}}{k-1}$\newline Si $k=1$ : Translation $\vec{v}$};
% Flèches
\draw[arrow] (start) -- (type1);
\draw[arrow] (type1) -- node[labelstyle, left] {Oui} (rot1);
\draw[arrow] (type1) -- node[labelstyle, right] {Non} (type2);
\draw[arrow] (type2) -- node[labelstyle, left] {Oui} (hom1);
\draw[arrow] (type2) -- node[labelstyle, right] {Non} (type3);
\draw[arrow] (type3) -- node[labelstyle, left] {Oui} (ht1);
\draw[arrow] (type3) -- node[labelstyle, right] {Non} ++(3,0) node[process, anchor=west] {Autres cas (non au programme 1re Tech)};
\end{tikzpicture}
\legende{Arbre de décision pour identifier la nature de la composée de deux transformations affines directes.}
\end{popfigure}

\begin{aretenir}[Récapitulatif des composées au programme]
\begin{itemize}
    \item \textbf{Rotation + Rotation (même centre $\Omega$) :} Rotation de centre $\Omega$, angle somme des angles ($\alpha_1+\alpha_2$).
    \item \textbf{Rotation + Rotation (centres distincts $\Omega_1 \neq \Omega_2$) :}
    \begin{itemize}
        \item Si $\alpha_1 + \alpha_2 \not\equiv 0 \pmod{2\pi}$ : Rotation d'angle $\alpha_1+\alpha_2$. Le centre se construit géométriquement (intersection des médiatrices de $[M M']$ pour deux points $M$ distincts).
        \item Si $\alpha_1 + \alpha_2 \equiv 0 \pmod{2\pi}$ : Translation de vecteur $\vec{v} = \overrightarrow{\Omega_1\Omega_2} - R_{\alpha_1}(\overrightarrow{\Omega_1\Omega_2})$.
    \end{itemize}
    \item \textbf{Homothétie + Homothétie (même centre $\Omega$) :} Homothétie de centre $\Omega$, rapport produit des rapports ($k_1 k_2$).
    \item \textbf{Homothétie $H(\Omega,k)$ + Translation $T_{\vec{v}}$ (dans n'importe quel ordre, $k \neq 1$) :}
    \begin{itemize}
        \item Composée $T_{\vec{v}} \circ H_{\Omega,k}$ : Homothétie de rapport $k$, centre $\Omega' = \Omega - \frac{\vec{v}}{k-1}$.
        \item Composée $H_{\Omega,k} \circ T_{\vec{v}}$ : Homothétie de rapport $k$, centre $\Omega'' = \Omega - \frac{k\vec{v}}{k-1}$.
        \item Si $k = 1$ : Translation de vecteur $\vec{v}$ (cas trivial).
    \end{itemize}
\end{itemize}
\textbf{Point commun :} Toutes ces composées sont des \textbf{similitudes directes} (conservation des angles, de l'alignement, transformation d'une droite en droite).
\end{aretenir}

\courssub{Stratégie générale et rédaction type}

Au Probatoire, on ne vous demande pas seulement le résultat final, mais une \textbf{justification rédigée}. Une réponse du type « C'est une rotation » sans préciser le centre et l'angle, ou « C'est une homothétie » sans donner le rapport et le centre, \textbf{perd la majorité des points}.

\begin{methode}[Méthode en 4 étapes pour traiter une composée]
\textbf{Étape 1 : Identifier les deux transformations.}
Lire l'énoncé, repérer les mots-clés : « rotation », « homothétie », « translation ». Noter $T_1$ la première appliquée, $T_2$ la seconde. La composée cherchée est $T = T_2 \circ T_1$ (on applique $T_1$ puis $T_2$).

\textbf{Étape 2 : Relever les éléments caractéristiques de chacune.}
\begin{itemize}
    \item Rotation : Centre ($\Omega$), Angle ($\alpha$, orienté, en degrés ou radians), Sens (souvent trigonométrique par défaut).
    \item Homothétie : Centre ($\Omega$), Rapport ($k \in \mathbb{R}^*$). Attention : si $k<0$, c'est une homothétie de rapport $|k|$ composée d'une rotation de $180^\circ$ (similitude directe).
    \item Translation : Vecteur ($\vec{v}$).
\end{itemize}

\textbf{Étape 3 : Appliquer le théorème adapté (selon l'arbre de décision).}
Cas par cas, écrire l'égalité de transformations (ex: $R_{\Omega_2,\alpha_2} \circ R_{\Omega_1,\alpha_1} = R_{\Omega,\alpha}$) et calculer les paramètres de la composée (centre, angle, rapport, vecteur).

\textbf{Étape 4 : Conclure par une phrase complète.}
\textit{« La transformation $T$ est une [nature] de [paramètres complets]. »}
Exemple : « La transformation $T$ est une rotation de centre $\Omega$ et d'angle $120^\circ$. » ou « $T$ est une homothétie de centre $I$ et de rapport $-3$. »
\end{methode}

\begin{attention}[Piège majeur au Probatoire : réponse incomplète]
\begin{itemize}
    \item \textbf{Ne jamais écrire :} « C'est une rotation. » $\rightarrow$ \textbf{0 point} sur la conclusion.
    \item \textbf{Écrire obligatoirement :} « C'est une rotation de centre $\Omega$ (coordonnées ou construction) et d'angle $60^\circ$. »
    \item Pour une homothétie : centre \textbf{ET} rapport.
    \item Pour une translation : vecteur (coordonnées ou définition géométrique).
    \item \textbf{Justifier chaque paramètre :} « car la composée de deux rotations de même centre... » ou « car le rapport est le produit $k_1 \times k_2 = ...$ ».
    \item \textbf{Attention au rapport négatif :} Une homothétie de rapport $k<0$ est une similitude directe d'angle $180^\circ$ et de rapport $|k|$. La composée avec une rotation ajoute les angles.
\end{itemize}
\end{attention}

\begin{exemplebox}[Rédaction type complète — Modèle à reproduire]
\textbf{Énoncé :} Soit $R_1$ la rotation de centre $O(0;0)$ et d'angle $+60^\circ$. Soit $H$ l'homothétie de centre $O$ et de rapport $k=2$. Déterminer la nature et les caractéristiques de la composée $T = H \circ R_1$.

\textbf{Solution rédigée :}
\begin{enumerate}
    \item \textbf{Identification :} $T_1 = R_1$ est une rotation de centre $O$ et d'angle $\alpha_1 = 60^\circ$. $T_2 = H$ est une homothétie de centre $O$ et de rapport $k=2$.
    \item \textbf{Théorème :} La composée d'une rotation et d'une homothétie de \textbf{même centre} $O$ est une similitude directe de centre $O$, d'angle l'angle de la rotation et de rapport le rapport de l'homothétie. (C'est une rotation-homothétie, ou similitude directe de centre $O$).
    \item \textbf{Calcul des paramètres :} Centre : $O$. Angle : $60^\circ$. Rapport : $2$.
    \item \textbf{Conclusion :} La transformation $T = H \circ R_1$ est une \textbf{similitude directe de centre $O$, d'angle $60^\circ$ et de rapport $2$}. (On peut aussi dire : c'est la rotation-homothétie $R_{O,60^\circ} \circ H_{O,2}$).
\end{enumerate}
\end{exemplebox}

\courssub{Applications concrètes : problèmes techniques résolus}

Nous appliquons la méthode à trois situations professionnelles typiques : topographie/urbanisme (double agrandissement), mécanique (engrenages), communication visuelle (logo).

\begin{exoresolu}[Problème 1 : Double agrandissement d'un plan de concession (Topographie / Urbanisme)]
\textbf{Énoncé :} Un géomètre-topographe lève le plan d'une concession à l'échelle $1/100$ (1 cm sur le plan = 1 m sur le terrain). Pour l'affichage en mairie, il réalise un premier agrandissement par homothétie de centre $O$ (coin du bureau) et de rapport $k_1 = 3$. Le plan est encore trop petit : il fait un second agrandissement, toujours de centre $O$, de rapport $k_2 = 2$. Quelle est la transformation globale subie par le plan initial ? Quelle est l'échelle finale du plan affiché ?

\textbf{Solution détaillée}
\begin{enumerate}
    \item \textbf{Identification :} $T_1 = H_1$ homothétie centre $O$, rapport $k_1=3$. $T_2 = H_2$ homothétie centre $O$, rapport $k_2=2$. Composée $T = H_2 \circ H_1$.
    \item \textbf{Théorème :} Composée de deux homothéties de même centre $O$ $\Rightarrow$ Homothétie de centre $O$, rapport $k = k_1 \times k_2$.
    \item \textbf{Calcul :} $k = 3 \times 2 = 6$.
    \item \textbf{Conclusion :} La transformation globale est une \textbf{homothétie de centre $O$ et de rapport $6$}.
    \item \textbf{Échelle finale :} L'échelle initiale était $1/100$. Un rapport $6$ multiplie les longueurs par $6$. L'échelle finale est $\frac{1}{100} \times 6 = \frac{6}{100} = \frac{3}{50}$. Soit \textbf{1 cm = 50 cm sur le terrain} (échelle $1/50$).
\end{enumerate}
\end{exoresolu}

\begin{exoresolu}[Problème 2 : Rotation d'un engrenage puis d'un autre (Mécanique)]
\textbf{Énoncé :} Dans une boîte de vitesses, un arbre d'entrée $A_1$ entraîne un premier engrenage (pignon) par une rotation $R_1$ de centre $\Omega_1$ et d'angle $+120^\circ$ (sens trigonométrique). Ce pignon en entraîne un second, fixe sur l'arbre de sortie, par une rotation $R_2$ de centre $\Omega_2$ (distinct de $\Omega_1$) et d'angle $-120^\circ$ (sens horaire). Déterminer la transformation globale qui relie la position angulaire de l'arbre d'entrée à celle de l'arbre de sortie.

\textbf{Solution détaillée}
\begin{enumerate}
    \item \textbf{Identification :} $R_1(\Omega_1, +120^\circ)$, $R_2(\Omega_2, -120^\circ)$. Centres distincts $\Omega_1 \neq \Omega_2$. Composée $T = R_2 \circ R_1$.
    \item \textbf{Somme des angles :} $\alpha = \alpha_1 + \alpha_2 = 120^\circ + (-120^\circ) = 0^\circ \equiv 0 \pmod{360^\circ}$.
    \item \textbf{Théorème (cas particulier) :} Composée de deux rotations de centres distincts dont la somme des angles est un multiple de $360^\circ$ $\Rightarrow$ \textbf{Translation}.
    \item \textbf{Vecteur de la translation :} $\vec{v} = \overrightarrow{\Omega_1\Omega_2} - R_{\alpha_1}(\overrightarrow{\Omega_1\Omega_2})$.
        Soit $\overrightarrow{\Omega_1\Omega_2} = \vec{u}$. $R_{+120^\circ}(\vec{u})$ tourne $\vec{u}$ de $120^\circ$.
        $\vec{v} = \vec{u} - R_{120^\circ}(\vec{u})$.
        \textit{(Si par exemple $\vec{u} = (4;0)$ en cm, alors $R_{120^\circ}(\vec{u}) = (-2; 2\sqrt{3})$ et $\vec{v} = (6; -2\sqrt{3})$ cm).}
    \item \textbf{Conclusion :} La transformation globale est une \textbf{translation de vecteur $\vec{v} = \overrightarrow{\Omega_1\Omega_2} - R_{+120^\circ}(\overrightarrow{\Omega_1\Omega_2})$}.
    \item \textbf{Interprétation mécanique :} L'arbre de sortie subit un déplacement pur (pas de rotation nette) par rapport à l'entrée : c'est un mouvement de translation. Les deux arbres sont parallèles et l'angle relatif est constant.
\end{enumerate}
\end{exoresolu}

\begin{exoresolu}[Problème 3 : Agrandissement puis décalage d'un logo publicitaire (Communication visuelle / PAO)]
\textbf{Énoncé :} Un graphiste crée un logo centré sur l'origine $O(0;0)$ d'un repère. Il applique d'abord une homothétie $H$ de centre $O$ et de rapport $k = 1,5$ pour l'agrandir. Ensuite, il déplace le logo vers le coin supérieur droit de l'affiche par une translation $T$ de vecteur $\vec{v} = \begin{pmatrix} 8 \\ 5 \end{pmatrix}$ (en cm). Déterminer la transformation globale $T \circ H$. Si le graphiste avait fait l'inverse (translation puis homothétie), le résultat serait-il le même ?

\textbf{Solution détaillée}
\begin{enumerate}
    \item \textbf{Cas 1 : $T \circ H$ (Homothétie puis Translation).}
    \begin{itemize}
        \item $H$ : centre $O(0;0)$, rapport $k=1,5 \neq 1$. $T$ : vecteur $\vec{v} = (8;5)$.
        \item Théorème : Composée $T \circ H$ = Homothétie de rapport $k=1,5$ et de centre $\Omega_1 = O - \frac{\vec{v}}{k-1}$.
        \item Calcul : $k-1 = 0,5$. $\Omega_1 = (0;0) - \frac{(8;5)}{0,5} = (0;0) - (16;10) = (-16;-10)$.
        \item Conclusion 1 : \textbf{Homothétie de centre $\Omega_1(-16;-10)$ et de rapport $1,5$}.
        \item \textit{Note :} Le logo final est centré sur $T(H(O)) = T(O) = (8;5)$, mais le \textbf{centre de la similitude globale} (point fixe) est $\Omega_1(-16;-10)$.
    \end{itemize}
    \item \textbf{Cas 2 : $H \circ T$ (Translation puis Homothétie).}
    \begin{itemize}
        \item On applique $T$ puis $H$. $T$ : vecteur $\vec{v} = (8;5)$. $H$ : centre $O$, rapport $k=1,5$.
        \item Théorème : Composée $H \circ T$ = Homothétie de rapport $k=1,5$ et de centre $\Omega_2 = O - \frac{k\vec{v}}{k-1}$.
        \item Calcul : $\Omega_2 = (0;0) - \frac{1,5 \times (8;5)}{0,5} = -3 \times (8;5) = (-24;-15)$.
        \item Conclusion 2 : \textbf{Homothétie de centre $\Omega_2(-24;-15)$ et de rapport $1,5$}.
        \item \textit{Note :} Le logo final est centré sur $H(T(O)) = H(8;5) = (12;7,5)$.
    \end{itemize}
    \item \textbf{Comparaison :}
    \begin{itemize}
        \item $T \circ H$ : Centre de la similitude $\Omega_1(-16;-10)$, position finale du logo $(8;5)$.
        \item $H \circ T$ : Centre de la similitude $\Omega_2(-24;-15)$, position finale du logo $(12;7,5)$.
    \end{itemize}
    \textbf{Les centres de la transformation globale sont différents, et les positions finales du logo sont différentes.} L'ordre des opérations change le résultat (non-commutativité). Le graphiste doit choisir l'ordre selon l'emplacement final souhaité.
\end{enumerate}
\end{exoresolu}

\begin{popfigure}
\centering
\begin{tikzpicture}[scale=0.65, every node/.style={font=\small}]
    % --- Cas 1 : T o H ---
    \begin{scope}
    % Repère
    \draw[->, popmuted] (-2,0) -- (12,0) node[right] {$x$ (cm)};
    \draw[->, popmuted] (0,-2) -- (0,10) node[above] {$y$ (cm)};
    \foreach \x in {0,2,...,10} \draw (\x,-0.1) -- (\x,0.1) node[below=2pt] {\x};
    \foreach \y in {0,2,...,8} \draw (-0.1,\y) -- (0.1,\y) node[left=2pt] {\y};
    % Logo initial (carré centré sur O)
    \draw[popblue, thick, fill=popblue!10] (-1,-1) rectangle (1,1);
    \node[popblue] at (0, -1.5) {Logo initial $L_0$};
    % Flèche H
    \draw[poporange, thick, -Stealth] (1.5, 0) -- (3.5, 0) node[midway, above] {$H(O, 1,5)$};
    % Logo après H (centré sur O, plus grand)
    \draw[poporange, thick, fill=poporange!10] (-1.5,-1.5) rectangle (1.5,1.5);
    \node[poporange] at (0, -2) {$L_1 = H(L_0)$};
    % Flèche T
    \draw[popgreen, thick, -Stealth] (2, 2) -- (7, 6) node[midway, above left] {$T(\vec{v})$};
    % Logo final T o H (centré sur (8,5))
    \draw[popgreen, thick, fill=popgreen!10] (6.5, 3.5) rectangle (9.5, 6.5);
    \node[popgreen] at (8, 3) {$L_2 = T(L_1)$};
    \node[popgreen, align=center] at (8, 7) {Position finale\\$(8;5)$};
    \end{scope}
    % --- Cas 2 : H o T ---
    \begin{scope}[xshift=16cm]
        \draw[->, popmuted] (-2,0) -- (16,0) node[right] {$x$};
        \draw[->, popmuted] (0,-2) -- (0,12) node[above] {$y$};
        \foreach \x in {0,2,...,14} \draw (\x,-0.1) -- (\x,0.1) node[below=2pt] {\x};
        \foreach \y in {0,2,...,10} \draw (-0.1,\y) -- (0.1,\y) node[left=2pt] {\y};
        % Logo initial
        \draw[popblue, thick, fill=popblue!10] (-1,-1) rectangle (1,1);
        \node[popblue] at (0, -1.5) {Logo initial $L_0$};
        % Flèche T
        \draw[popgreen, thick, -Stealth] (1.5, 0) -- (6, 4) node[midway, above right] {$T(\vec{v})$};
        % Logo après T (centré sur (8,5))
        \draw[popgreen, thick, fill=popgreen!10] (7, 4) rectangle (9, 6);
        \node[popgreen] at (8, 3.5) {$L_1' = T(L_0)$};
        \node[popgreen] at (8, 6.5) {Position $(8;5)$};
        % Flèche H
        \draw[poporange, thick, -Stealth] (9, 5) -- (12, 8) node[midway, above right] {$H(O, 1,5)$};
        % Logo final H o T (centré sur (12, 7.5))
        \draw[poporange, thick, fill=poporange!10] (10.5, 6) rectangle (13.5, 9);
        \node[poporange] at (12, 5.5) {$L_2' = H(L_1')$};
        \node[poporange, align=center] at (12, 9.5) {Position finale\\$(12;7,5)$};
    \end{scope}
\end{tikzpicture}
\legende{Comparaison des deux ordres : $T \circ H$ (gauche) place le logo en $(8;5)$, $H \circ T$ (droite) le place en $(12;7,5)$. Les centres mathématiques des similitudes globales ($\Omega_1, \Omega_2$) sont hors cadre (négatifs). L'ordre change la position finale.}
\end{popfigure}

\courssub{Problème inverse : retrouver les transformations}

Parfois, l'énoncé donne la composée et demande de retrouver deux transformations possibles. Il y a une infinité de solutions, mais on cherche les plus simples.

\begin{exemplebox}[Problème inverse : Décomposer une similitude donnée]
\textbf{Énoncé :} Soit $S$ la similitude directe qui transforme le triangle $ABC$ en $A'B'C'$ tel que $A(0;0)$, $B(1;0)$, $C(0;1)$ et $A'(2;2)$, $B'(5;2)$, $C'(2;5)$. Décomposer $S$ en une homothétie et une rotation de même centre, puis en une homothétie et une translation.

\textbf{Résolution}
\begin{enumerate}
    \item \textbf{Analyse de $S$ :} 
    $\overrightarrow{AB} = (1;0) \to \overrightarrow{A'B'} = (3;0)$. Longueur $\times 3$, angle $0^\circ$.
    $\overrightarrow{AC} = (0;1) \to \overrightarrow{A'C'} = (0;3)$. Longueur $\times 3$, angle $0^\circ$.
    $S$ est une \textbf{homothétie de rapport $k=3$} (rotation angle $0^\circ$). Ce n'est pas une rotation-homothétie non triviale, mais une homothétie pure. Il faut trouver son centre $\Omega$.
    \item \textbf{Décomposition 1 : Rotation + Homothétie de même centre (ici rotation nulle).}
    $S = H_{\Omega, 3} \circ Id$. On cherche $\Omega$ tel que $H_{\Omega,3}(A)=A'$.
    $\overrightarrow{\Omega A'} = 3 \overrightarrow{\Omega A} \iff \vec{A'} - \vec{\Omega} = 3(\vec{A} - \vec{\Omega}) \iff 2\vec{\Omega} = 3\vec{A} - \vec{A'}$.
    $A(0;0), A'(2;2) \Rightarrow 2\vec{\Omega} = (0;0) - (2;2) = (-2;-2) \Rightarrow \Omega(-1;-1)$.
    \textbf{Réponse :} $S = H_{\Omega(-1;-1), 3}$.
    \item \textbf{Décomposition 2 : Homothétie + Translation.}
    On veut $S = T_{\vec{v}} \circ H_{O, 3}$ (choisissons centre $O$ pour $H$).
    $H_{O,3}(A) = (0;0) = A$. Il faut $T_{\vec{v}}(A) = A'(2;2) \Rightarrow \vec{v} = (2;2)$.
    Vérif : $H_{O,3}(B) = (3;0)$, $T_{(2;2)}(3;0) = (5;2) = B'$. $H_{O,3}(C) = (0;3)$, $T_{(2;2)}(0;3) = (2;5) = C'$. \textbf{Ça marche.}
    \textbf{Réponse :} $S = T_{(2;2)} \circ H_{O, 3}$.
    \item \textbf{Autre décomposition (Translation puis Homothétie) :} $S = H_{O,3} \circ T_{\vec{w}}$.
    $H_{O,3} \circ T_{\vec{w}} = T_{3\vec{w}} \circ H_{O,3}$. On veut $T_{3\vec{w}} \circ H_{O,3} = T_{(2;2)} \circ H_{O,3} \Rightarrow 3\vec{w} = (2;2) \Rightarrow \vec{w} = (2/3; 2/3)$.
    \textbf{Réponse :} $S = H_{O, 3} \circ T_{(2/3; 2/3)}$.
\end{enumerate}
\end{exemplebox}

\courssub{Complément pour la Terminale : l'écriture complexe unificatrice}

\begin{savaistu}[Le sais-tu ? — Vers l'écriture complexe $z \mapsto az + b$]
En classe de Terminale, vous découvrirez que le plan peut être muni d'une structure de corps : les \textbf{nombres complexes}. Un point $M(x;y)$ s'identifie au nombre complexe $z = x + iy$.

\textbf{Toutes les similitudes directes du plan} (rotations, homothéties, translations, et leurs composées) s'écrivent alors sous une \textbf{forme unique} :
\[
\boxed{z \longmapsto a z + b \qquad \text{avec } a \in \mathbb{C}^*,\ b \in \mathbb{C}}
\]
\begin{itemize}
    \item Le module $|a|$ donne le \textbf{rapport} de la similitude (rapport d'homothétie).
    \item L'argument $\arg(a)$ donne l'\textbf{angle} de la rotation (en radians).
    \item Le terme $b$ gère la \textbf{translation} (et le déplacement du centre).
\end{itemize}

\textbf{Exemples de traduction :}
\begin{itemize}
    \item Rotation centre $0$, angle $\theta$ : $a = e^{i\theta},\ b=0$.
    \item Homothétie centre $0$, rapport $k$ : $a = k,\ b=0$.
    \item Translation vecteur $\vec{v}(u;v)$ : $a = 1,\ b = u+iv$.
    \item Rotation centre $\omega$, angle $\theta$ : $z \mapsto e^{i\theta}(z-\omega)+\omega = e^{i\theta}z + \omega(1-e^{i\theta})$. Ici $a=e^{i\theta}$, $b=\omega(1-e^{i\theta})$.
\end{itemize}

\textbf{Pourquoi c'est puissant :} La composée de deux similitudes $z \mapsto a_1 z + b_1$ puis $z \mapsto a_2 z + b_2$ s'obtient par simple substitution :
\[
z \mapsto a_2(a_1 z + b_1) + b_2 = (a_2 a_1) z + (a_2 b_1 + b_2)
\]
\textbf{Le rapport se multiplie ($a_2 a_1$), l'angle s'additionne (argument du produit), la translation se combine.} Tous les théorèmes de cette leçon deviennent de simples calculs algébriques sur $a$ et $b$ ! C'est l'outil ultime pour la CAO, la robotique et le traitement d'images.
\end{savaistu}

\courssub{Bilan des savoir-faire et auto-évaluation}

\begin{aretenir}[Check-list des compétences pour le Probatoire]
\noindent\begin{tabularx}{\linewidth}{|X|c|}
\hline
\textbf{Compétence} & \textbf{Maîtrise} \\
\hline
Identifier la nature de deux transformations composées (rotation, homothétie, translation) & $\square$ \\
Relever correctement les paramètres : centre, angle, rapport, vecteur & $\square$ \\
Appliquer le bon théorème selon les cas (même centre / centres distincts / homothétie+translation) & $\square$ \\
Calculer les paramètres de la composée (produit des rapports, somme des angles, image du centre) & $\square$ \\
\textbf{Rédiger une justification complète :} « La transformation est une [nature] de centre/rapport/angle/vecteur [valeurs] car [théorème invoqué]. » & $\square$ \\
Résoudre un problème concret (plan, engrenage, logo) en modélisant par des transformations & $\square$ \\
Décomposer une similitude donnée en deux transformations simples (problème inverse) & $\square$ \\
\hline
\end{tabularx}
\end{aretenir}

\begin{exercice}[Entraînement mixte — Type Probatoire]
\begin{enumerate}
    \item Soit $R$ la rotation de centre $I(1;1)$ d'angle $+90^\circ$ et $H$ l'homothétie de centre $I$ de rapport $k=-2$. Déterminer la nature et les caractéristiques de $H \circ R$.
    \item Un technicien programme un bras robotisé. Il applique une rotation $R_1$ de centre $O$ angle $+60^\circ$, puis une rotation $R_2$ de centre $A(4;0)$ angle $-60^\circ$. Quelle est la transformation globale ? Préciser sa nature et son/ses élément(s) caractéristique(s).
    \item Une pièce mécanique subit une homothétie $H$ de centre $\Omega$ rapport $k=0,5$ (rétrécissement), puis une translation $T$ de vecteur $\vec{v} = \overrightarrow{\Omega A}$ où $A$ est un point tel que $\Omega A = 10$ cm. Déterminer la composée $T \circ H$ (nature, rapport, centre).
    \item \textbf{Problème inverse :} On donne la similitude $S$ définie par $S(M) = M'$ tel que pour tout point $M$, $\overrightarrow{OM'} = 2 \times \text{rot}_{+90^\circ}(\overrightarrow{OM}) + \begin{pmatrix} 1 \\ -1 \end{pmatrix}$. Décomposer $S$ en une rotation-homothétie de centre $O$ suivie d'une translation. Puis décomposer $S$ en une translation suivie d'une rotation-homothétie de centre $O$.
\end{enumerate}
\end{exercice}

\begin{corrige}[Exercice n°15 - Entraînement mixte]
\begin{enumerate}
    \item \textbf{Identification :} $R(I, +90^\circ)$, $H(I, -2)$. Même centre $I$.
    \textbf{Analyse :} $H_{I, -2}$ est une homothétie de rapport $2$ composée d'une rotation de $180^\circ$ (similitude directe d'angle $180^\circ$, rapport $2$).
    \textbf{Théorème :} Composée de deux similitudes directes de même centre = Similitude directe de centre $I$, angle somme des angles, rapport produit des rapports.
    \textbf{Calcul :} Angle $= 90^\circ + 180^\circ = 270^\circ$ (ou $-90^\circ$). Rapport $= 1 \times 2 = 2$.
    \textbf{Conclusion :} La transformation $H \circ R$ est une \textbf{similitude directe de centre $I(1;1)$, d'angle $270^\circ$ (ou $-90^\circ$) et de rapport $2$}.
    \item \textbf{Identification :} $R_1(O, +60^\circ)$, $R_2(A(4;0), -60^\circ)$. Centres distincts. Somme angles $= 0^\circ \equiv 0 [360^\circ]$.
    \textbf{Théorème :} Composée = Translation.
    \textbf{Vecteur :} $\vec{v} = \overrightarrow{OA} - R_{+60^\circ}(\overrightarrow{OA})$.
    $\overrightarrow{OA} = (4;0)$. $R_{+60^\circ}(4;0) = (4\cos 60^\circ; 4\sin 60^\circ) = (2; 2\sqrt{3})$.
    $\vec{v} = (4;0) - (2; 2\sqrt{3}) = (2; -2\sqrt{3})$.
    \textbf{Conclusion :} Translation de vecteur $\vec{v}(2; -2\sqrt{3})$.
    \item \textbf{Identification :} $H(\Omega, 0,5)$, $T(\vec{v}=\overrightarrow{\Omega A})$. $k=0,5 \neq 1$.
    \textbf{Théorème :} $T \circ H$ = Homothétie rapport $k=0,5$, centre $\Omega' = \Omega - \frac{\vec{v}}{k-1}$.
    \textbf{Calcul :} $k-1 = -0,5$. $\Omega' = \Omega - \frac{\overrightarrow{\Omega A}}{-0,5} = \Omega + 2\overrightarrow{\Omega A} = \Omega + 2(\vec{A} - \vec{\Omega}) = 2\vec{A} - \vec{\Omega}$.
    Le centre $\Omega'$ est le symétrique de $\Omega$ par rapport à $A$ (ou $A$ est le milieu de $[\Omega\Omega']$). Distance $\Omega\Omega' = 2 \times \Omega A = 20$ cm.
    \textbf{Conclusion :} Homothétie de centre $\Omega'$ (symétrique de $\Omega$ par rapport à $A$) et de rapport $0,5$.
    \item \textbf{Analyse (notation vectorielle) :} $S$ est la similitude directe de rapport $2$, angle $+90^\circ$, centre $O$, composée avec la translation de vecteur $\vec{b} = (1;-1)$.
    \textbf{Décomp 1 (Rotation-Homothétie puis Translation) :} $S = T_{\vec{b}} \circ (R_{O,90^\circ} \circ H_{O,2})$.
    \textbf{Décomp 2 (Translation puis Rotation-Homothétie) :} On cherche $\vec{w}$ tel que $S = (R_{O,90^\circ} \circ H_{O,2}) \circ T_{\vec{w}}$.
    $(R_{O,90^\circ} \circ H_{O,2}) \circ T_{\vec{w}} = T_{S_O(\vec{w})} \circ (R_{O,90^\circ} \circ H_{O,2})$ où $S_O$ est la similitude centre $O$ (rapport 2, angle $90^\circ$).
    Il faut $S_O(\vec{w}) = \vec{b} \Rightarrow \vec{w} = S_O^{-1}(\vec{b})$.
    $S_O^{-1}$ : rapport $1/2$, angle $-90^\circ$.
    $\vec{w} = \frac{1}{2} \text{rot}_{-90^\circ}(1;-1) = \frac{1}{2}(-1;-1) = (-0,5; -0,5)$.
    \textbf{Réponse :} $S = T_{(-0,5;-0,5)} \circ (R_{O,90^\circ} \circ H_{O,2})$.
\end{enumerate}
\end{corrige}

\begin{propriete}[Propriété récapitulative : Affinité des composées]
Toutes les transformations étudiées dans cette leçon (rotations, homothéties, translations, et leurs composées deux à deux) sont des \textbf{transformations affines}.
Elles satisfont les deux propriétés fondamentales :
\begin{enumerate}
    \item \textbf{Conservation de l'alignement :} Si $A, B, C$ sont alignés, alors leurs images $A', B', C'$ sont alignés.
    \item \textbf{Conservation des rapports de longueurs sur une droite :} Si $A, B, C$ alignés, alors $\frac{A'B'}{A'C'} = \frac{AB}{AC}$ (pour les isométries et homothéties) ou plus généralement conservation des barycentres.
\end{enumerate}
\textbf{Conséquence immédiate :} L'image d'un segment est un segment, l'image d'un parallélogramme est un parallélogramme, l'image d'une droite est une droite. C'est ce qui garantit qu'un plan technique reste « lisible » et géométriquement cohérent après agrandissement, rotation ou déplacement.
\end{propriete}

\courssub{En résumé}

Cette leçon sur les composées de transformations affines directes vous a permis de :
\begin{itemize}
    \item Maîtriser les \textbf{4 cas fondamentaux} du programme (Rotation+Rotation même centre, Rotation+Rotation centres distincts, Homothétie+Homothétie même centre, Homothétie+Translation).
    \item Acquérir une \textbf{méthode rigoureuse en 4 étapes} (Identifier, Relever, Appliquer, Conclure) exigée au Probatoire.
    \item Vous exercer sur des \textbf{problèmes métiers} (topographie, mécanique, PAO) où la modélisation par transformations est l'outil standard.
    \item Entrevoir la \textbf{puissance de l'écriture complexe} ($z \mapsto az+b$) qui unifiera tout cela en Terminale.
\end{itemize}
La clé de la réussite : \textbf{ne jamais oublier les éléments caractéristiques dans la conclusion}. Un centre sans coordonnées, un angle sans valeur, un rapport sans calcul = réponse incomplète = points perdus. Entraînez-vous à rédiger « comme au Probatoire » dès maintenant.

\newpage

\coursec{Activité d'intégration}

\courssub{Objectifs de l'activité}

À la fin de cette activité d'intégration, l'élève sera capable de :
\begin{itemize}
\item appliquer concrètement une \cle{homothétie} et une \cle{translation} à une figure donnée par les coordonnées de ses sommets dans un repère ;
\item déterminer la nature et les éléments caractéristiques (rapport, centre) de la \cle{composée d'une homothétie et d'une translation} ;
\item composer deux \cle{homothéties de même centre} et interpréter le résultat ;
\item rédiger et justifier une démarche complète sous la forme d'une note technique destinée à un non-spécialiste.
\end{itemize}

\courssub{Situation}

La coopérative des producteurs de cacao de la région du Centre (COPROCAO), installée à Mbalmayo, souhaite moderniser son image. Son logo est un triangle stylisé représentant une cabosse de cacao. Pour l'imprimer sur ses sacs de jute de différentes tailles, la coopérative fait appel à un graphiste de Yaoundé.

Le graphiste travaille dans un repère orthonormé $(O\,;\vec{\imath},\vec{\jmath})$ de l'écran de son ordinateur, l'unité étant le centimètre. Le logo original est le triangle $ABC$ avec :
\[ A(2\,;1), \qquad B(6\,;1), \qquad C(4\,;5). \]

Pour l'impression sur les grands sacs de \SI{50}{kg}, le graphiste procède en deux étapes :
\begin{itemize}
\item \textbf{Étape 1 :} il agrandit le logo par l'homothétie $h$ de centre $O$ (l'origine du repère) et de rapport $k=\dfrac{3}{2}$ ; le triangle obtenu est noté $A'B'C'$ ;
\item \textbf{Étape 2 :} il positionne le logo agrandi sur le sac en lui appliquant la translation $t$ de vecteur $\vec{u}\begin{pmatrix} 4 \\ 2 \end{pmatrix}$ ; le triangle final est noté $A''B''C''$.
\end{itemize}

Pour les petits emballages de \SI{5}{kg}, le graphiste applique au logo original une homothétie $h'$ de même centre $O$ et de rapport $k'=\dfrac{2}{3}$.

Le président de la coopérative, qui n'est pas mathématicien, demande une note technique claire expliquant la démarche, les dimensions du logo final et les propriétés des transformations utilisées.

\courssub{Consignes}

\begin{enumerate}
\item \textbf{Logo agrandi.} Déterminer les coordonnées des points $A'$, $B'$ et $C'$, images de $A$, $B$ et $C$ par l'homothétie $h$ de centre $O$ et de rapport $\dfrac{3}{2}$.
\item \textbf{Logo final sur le grand sac.} Déterminer les coordonnées des points $A''$, $B''$ et $C''$, images de $A'$, $B'$ et $C'$ par la translation $t$ de vecteur $\vec{u}\begin{pmatrix} 4 \\ 2 \end{pmatrix}$.
\item \textbf{Nature de la transformation globale.} On note $f = t \circ h$ la transformation qui envoie directement le logo original $ABC$ sur le logo final $A''B''C''$.
\begin{enumerate}
\item Pour un point $M(x\,;y)$ quelconque, exprimer les coordonnées $(x''\,;y'')$ de son image $M'' = f(M)$ en fonction de $x$ et $y$.
\item En déduire que $f$ est une homothétie dont on précisera le rapport, puis déterminer les coordonnées de son centre $\Omega$.
\item Vérifier que $\Omega$, $A$ et $A''$ sont alignés.
\end{enumerate}
\item \textbf{Petit emballage.} On applique au logo original l'homothétie $h'$ de centre $O$ et de rapport $\dfrac{2}{3}$.
\begin{enumerate}
\item Donner les coordonnées des sommets du logo réduit.
\item Que dire de la composée $h' \circ h$ des deux homothéties de même centre $O$ ? Justifier et interpréter concrètement le résultat pour la coopérative.
\end{enumerate}
\item \textbf{Note technique.} Rédiger une note technique d'une quinzaine de lignes à l'attention du président de la COPROCAO, expliquant sans formules compliquées la démarche du graphiste, les dimensions obtenues et l'intérêt pratique des résultats des questions 3 et 4.
\end{enumerate}

\courssub{Ressources à mobiliser}

\begin{aretenir}[Formules utiles]
\begin{itemize}
\item \textbf{Homothétie de centre $O$ (origine) et de rapport $k$ :} l'image de $M(x\,;y)$ est $M'(kx\,;ky)$.
\item \textbf{Translation de vecteur $\vec{u}\begin{pmatrix} a \\ b \end{pmatrix}$ :} l'image de $M(x\,;y)$ est $M'(x+a\,;y+b)$.
\item \textbf{Composée d'une homothétie de rapport $k \neq 1$ et d'une translation :} c'est une homothétie de même rapport $k$, dont le centre se détermine en cherchant le point invariant $\Omega$ tel que $f(\Omega)=\Omega$.
\item \textbf{Composée de deux homothéties de même centre $O$, de rapports $k$ et $k'$ :} c'est l'homothétie de centre $O$ et de rapport $k \times k'$. Si $k \times k' = 1$, cette composée est l'\cle{identité} (chaque point est sa propre image).
\item Une homothétie de rapport $k$ multiplie les longueurs par $|k|$ et les aires par $k^2$.
\end{itemize}
\end{aretenir}

\courssub{Démarche et corrigé commenté}

\begin{exoresolu}[Le logo de la coopérative : corrigé complet]
\textbf{Énoncé rappelé.} $A(2\,;1)$, $B(6\,;1)$, $C(4\,;5)$ ; $h$ homothétie de centre $O$ et de rapport $\dfrac{3}{2}$ ; $t$ translation de vecteur $\vec{u}\begin{pmatrix} 4 \\ 2 \end{pmatrix}$ ; $h'$ homothétie de centre $O$ et de rapport $\dfrac{2}{3}$.

\tcblower

\textbf{Question 1 : coordonnées de $A'$, $B'$, $C'$.}

Pour une homothétie de centre $O$ et de rapport $k$, on multiplie chaque coordonnée par $k$. Ici $k=\dfrac{3}{2}$ :
\begin{align*}
A' &= \left(\tfrac{3}{2}\times 2\,;\ \tfrac{3}{2}\times 1\right) = (3\,;\ 1{,}5),\\
B' &= \left(\tfrac{3}{2}\times 6\,;\ \tfrac{3}{2}\times 1\right) = (9\,;\ 1{,}5),\\
C' &= \left(\tfrac{3}{2}\times 4\,;\ \tfrac{3}{2}\times 5\right) = (6\,;\ 7{,}5).
\end{align*}

\textbf{Question 2 : coordonnées de $A''$, $B''$, $C''$.}

La translation de vecteur $\vec{u}\begin{pmatrix} 4 \\ 2 \end{pmatrix}$ ajoute $4$ à l'abscisse et $2$ à l'ordonnée :
\begin{align*}
A'' &= (3+4\,;\ 1{,}5+2) = (7\,;\ 3{,}5),\\
B'' &= (9+4\,;\ 1{,}5+2) = (13\,;\ 3{,}5),\\
C'' &= (6+4\,;\ 7{,}5+2) = (10\,;\ 9{,}5).
\end{align*}
Le logo final sur le grand sac est donc le triangle $A''(7\,;3{,}5)$, $B''(13\,;3{,}5)$, $C''(10\,;9{,}5)$. Sa base mesure $A''B'' = 13-7 = \SI{6}{cm}$ et sa hauteur $9{,}5-3{,}5 = \SI{6}{cm}$.

\textbf{Question 3 : nature de $f = t \circ h$.}

\textbf{a)} Soit $M(x\,;y)$. Par $h$, on obtient $M'\left(\dfrac{3}{2}x\,;\dfrac{3}{2}y\right)$, puis par $t$ :
\[ M''\left(\tfrac{3}{2}x + 4\,;\ \tfrac{3}{2}y + 2\right). \]

\textbf{b)} Cherchons un éventuel point invariant $\Omega(a\,;b)$, c'est-à-dire tel que $f(\Omega)=\Omega$ :
\[
\begin{cases}
\dfrac{3}{2}a + 4 = a \\[4pt]
\dfrac{3}{2}b + 2 = b
\end{cases}
\Longleftrightarrow
\begin{cases}
4 = a - \dfrac{3}{2}a = -\dfrac{1}{2}a \\[4pt]
2 = b - \dfrac{3}{2}b = -\dfrac{1}{2}b
\end{cases}
\Longleftrightarrow
\begin{cases}
a = -8 \\
b = -4
\end{cases}
\]
Il existe donc un unique point invariant $\Omega(-8\,;-4)$. Vérifions que $f$ est l'homothétie de centre $\Omega$ et de rapport $\dfrac{3}{2}$. Pour $M(x\,;y)$ et $M'' = f(M)$ :
\[
\overrightarrow{\Omega M''} = \begin{pmatrix} \tfrac{3}{2}x + 4 - (-8) \\[2pt] \tfrac{3}{2}y + 2 - (-4) \end{pmatrix}
= \begin{pmatrix} \tfrac{3}{2}x + 12 \\[2pt] \tfrac{3}{2}y + 6 \end{pmatrix}
= \dfrac{3}{2}\begin{pmatrix} x + 8 \\ y + 4 \end{pmatrix}
= \dfrac{3}{2}\,\overrightarrow{\Omega M}.
\]
Donc $f$ est bien l'\cle{homothétie} de centre $\Omega(-8\,;-4)$ et de rapport $k=\dfrac{3}{2}$.

\textbf{Conclusion conforme au cours :} la composée d'une homothétie de rapport $k \neq 1$ et d'une translation est une homothétie de même rapport $k$.

\textbf{c)} Vérification de l'alignement de $\Omega$, $A$ et $A''$ :
\[
\overrightarrow{\Omega A} = \begin{pmatrix} 2-(-8) \\ 1-(-4) \end{pmatrix} = \begin{pmatrix} 10 \\ 5 \end{pmatrix},
\qquad
\overrightarrow{\Omega A''} = \begin{pmatrix} 7-(-8) \\ 3{,}5-(-4) \end{pmatrix} = \begin{pmatrix} 15 \\ 7{,}5 \end{pmatrix}.
\]
On constate que $\overrightarrow{\Omega A''} = \dfrac{3}{2}\,\overrightarrow{\Omega A}$ : les vecteurs sont colinéaires, donc $\Omega$, $A$ et $A''$ sont alignés. C'est cohérent : dans une homothétie, un point, son image et le centre sont toujours alignés.

\textbf{Question 4 : petit emballage.}

\textbf{a)} En multipliant les coordonnées par $\dfrac{2}{3}$ :
\[ A_1\left(\tfrac{4}{3}\,;\tfrac{2}{3}\right), \qquad B_1\left(4\,;\tfrac{2}{3}\right), \qquad C_1\left(\tfrac{8}{3}\,;\tfrac{10}{3}\right). \]

\textbf{b)} La composée $h' \circ h$ de deux homothéties de même centre $O$ est l'homothétie de centre $O$ et de rapport :
\[ k \times k' = \dfrac{3}{2}\times\dfrac{2}{3} = 1. \]
Une homothétie de rapport $1$ est l'\cle{identité} : tout point est sa propre image. Concrètement, si le graphiste applique successivement l'agrandissement de rapport $\dfrac{3}{2}$ puis la réduction de rapport $\dfrac{2}{3}$ (ou l'inverse), il retrouve exactement le logo original : les deux transformations se compensent. On dit que $h$ et $h'$ sont des homothéties \emph{réciproques} l'une de l'autre.

\textbf{Question 5 : note technique (modèle de rédaction).}

\emph{Note technique à l'attention du Président de la COPROCAO — Objet : reproduction du logo sur les emballages.}

\emph{Monsieur le Président, le logo triangulaire de la coopérative, dont la base mesure \SI{4}{cm} et la hauteur \SI{4}{cm} dans le fichier original, a été agrandi d'une fois et demie (rapport $\frac{3}{2}$) puis déplacé pour être centré sur les grands sacs de \SI{50}{kg}. Le logo imprimé mesurera \SI{6}{cm} de base et \SI{6}{cm} de hauteur, soit une surface d'impression $2{,}25$ fois plus grande que l'original, ce qui garantit une bonne visibilité. Nous avons vérifié que l'ensemble de l'opération (agrandissement puis déplacement) revient à un seul agrandissement de même rapport effectué depuis un point fixe bien précis : le dessin n'est donc ni déformé ni pivoté, les proportions sont parfaitement respectées. Pour les petits emballages de \SI{5}{kg}, une réduction aux deux tiers (rapport $\frac{2}{3}$) est appliquée ; cette réduction compense exactement l'agrandissement précédent, ce qui permet de revenir au fichier original à tout moment sans aucune perte de qualité. Les coûts d'impression pourront être calculés à partir de ces dimensions. Le graphiste.}
\end{exoresolu}

\begin{attention}[Erreurs fréquentes à éviter]
\begin{itemize}
\item Ne pas additionner les rapports de deux homothéties composées : on les \textbf{multiplie} ($\tfrac{3}{2}\times\tfrac{2}{3}=1$, et non $\tfrac{3}{2}+\tfrac{2}{3}$).
\item Dans la composée $t \circ h$, on applique d'abord $h$, puis $t$ : l'ordre de rédaction des calculs doit respecter cet enchaînement.
\item Le centre de la composée d'une homothétie et d'une translation n'est \textbf{pas} le centre de l'homothétie de départ : il faut résoudre l'équation du point invariant.
\item Une homothétie de rapport $k$ multiplie les longueurs par $|k|$ mais les aires par $k^2$ : ne pas confondre les deux.
\end{itemize}
\end{attention}

\courssub{Bilan de l'activité}

Cette activité a permis de mettre en évidence, sur un problème authentique de reproduction graphique ancré dans le contexte agricole camerounais, les résultats essentiels de la leçon :
\begin{itemize}
\item la composée d'une homothétie de rapport $k \neq 1$ et d'une translation est une \cle{homothétie de même rapport} $k$, dont le centre s'obtient en résolvant une équation de point invariant ;
\item la composée de deux homothéties de même centre est une homothétie de même centre dont le rapport est le \textbf{produit} des rapports ; lorsque ce produit vaut $1$, on obtient l'identité, ce qui traduit concrètement la possibilité d'agrandir puis de réduire un visuel sans le déformer ;
\item les transformations affines conservent les formes (pas de déformation du logo), ce qui justifie leur usage en graphisme, en sérigraphie et en impression industrielle ;
\item la rédaction d'une note technique montre qu'un résultat mathématique doit pouvoir être communiqué clairement à un décideur non spécialiste : c'est une compétence professionnelle à part entière.
\end{itemize}

\newpage

\coursec{Méthodes et exercices résolus}

\coursec{Transformations affines du plan : composées}

\courssub{Généralités sur la composée de deux transformations}

\begin{definition}[Composée de deux transformations]
Soient $f$ et $g$ deux transformations du plan. La \cle{composée} de $f$ par $g$, notée $f \circ g$, est la transformation qui associe à tout point $M$ le point $f(g(M))$.
L'ordre est essentiel : on applique \textbf{d'abord} $g$, \textbf{ensuite} $f$.
\end{definition}

\begin{propriete}[Propriétés générales]
\begin{itemize}
\item \textbf{Associativité :} $(f \circ g) \circ h = f \circ (g \circ h)$. On peut écrire $f \circ g \circ h$ sans parenthèses.
\item \textbf{Non-commutativité :} En général, $f \circ g \neq g \circ f$. L'ordre des transformations change le résultat.
\item \textbf{Inverse :} Si $f$ et $g$ sont bijectives, $(f \circ g)^{-1} = g^{-1} \circ f^{-1}$.
\end{itemize}
\end{propriete}

\begin{aretenir}[Vocabulaire et notation au Probatoire]
\begin{itemize}
\item « La transformation $f$ suivie de la transformation $g $ » $\rightarrow$ $g \circ f$.
\item « La composée de $r_1$ puis $r_2$ » $\rightarrow$ $r_2 \circ r_1$.
\item On lit $r_2 \circ r_1$ : « $r_2$ rond $r_1$ » ou « $r_2$ composée à $r_1$ ».
\end{itemize}
\end{aretenir}

\begin{exemplebox}[Ordre des opérations]
Soit $h$ l'homothétie de centre $O$, rapport $2$ et $t$ la translation de vecteur $\vec{u}$.
\begin{itemize}
\item $t \circ h$ : on agrandit d'abord (centre $O$), puis on déplace.
\item $h \circ t$ : on déplace d'abord, puis on agrandit (centre $O$).
Les images d'un même point $M$ sont généralement différentes.
\end{itemize}
\end{exemplebox}

\courssub{Composée de deux rotations de même centre}

\begin{propriete}[Composée de deux rotations de même centre]
Soient $r_1$ et $r_2$ deux rotations de \textbf{même centre} $\Omega$ et d'angles orientés respectifs $\alpha$ et $\beta$.
La composée $r_2 \circ r_1$ est la \cle{rotation} de centre $\Omega$ et d'angle $\alpha + \beta$ (modulo $360°$).
\begin{itemize}
\item Si $\alpha + \beta \equiv 0 \pmod{360°}$, $r_2 \circ r_1 = \text{Id}$ (identité).
\item Si $\alpha + \beta \equiv 180° \pmod{360°}$, $r_2 \circ r_1$ est la \cle{symétrie centrale} de centre $\Omega$.
\end{itemize}
\textbf{Démonstration exigible :} Pour tout point $M$, $r_1(M)=M_1$ avec $(\overrightarrow{\Omega M}, \overrightarrow{\Omega M_1})=\alpha$ et $\Omega M_1=\Omega M$. Ensuite $r_2(M_1)=M_2$ avec $(\overrightarrow{\Omega M_1}, \overrightarrow{\Omega M_2})=\beta$ et $\Omega M_2=\Omega M_1$. Donc $(\overrightarrow{\Omega M}, \overrightarrow{\Omega M_2})=\alpha+\beta$ et $\Omega M_2=\Omega M$. C'est bien la rotation de centre $\Omega$, angle $\alpha+\beta$.
\end{propriete}

\begin{methode}[Composer deux rotations de même centre]
Pour déterminer la nature de $r_2 \circ r_1$ où $r_1$ et $r_2$ sont deux rotations de \textbf{même centre} $\Omega$ et d'angles respectifs $\alpha$ et $\beta$ :
\begin{enumerate}
\item Vérifier que les deux rotations ont le \cle{même centre} $\Omega$.
\item Additionner les angles orientés : la composée est la \cle{rotation} de centre $\Omega$ et d'angle $\alpha + \beta$.
\item Réduire l'angle somme modulo $360°$ :
      \begin{itemize}
      \item si $\alpha + \beta \equiv 0 \pmod{360°}$ : composée = \cle{identité} ;
      \item si $\alpha + \beta \equiv 180° \pmod{360°}$ : composée = \cle{symétrie centrale} de centre $\Omega$ ;
      \item sinon : composée = rotation de centre $\Omega$, angle $\alpha+\beta$.
      \end{itemize}
\item Conclure par une phrase complète : « $r_2 \circ r_1$ est la rotation de centre $\Omega$ et d'angle $\alpha+\beta$ » (ou identité / symétrie centrale).
\end{enumerate}
\textbf{Réflexe :} l'ordre n'a pas d'importance : $r_2 \circ r_1 = r_1 \circ r_2$ car l'addition des angles est commutative.
\end{methode}

\begin{exoresolu}[Composée de deux rotations de même centre -- Atelier mécanique]
Dans un atelier de mécanique à Douala, une pièce circulaire est fixée sur un plateau tournant de centre $O$. On effectue d'abord une rotation $r_1$ de centre $O$ et d'angle $70°$, puis une rotation $r_2$ de centre $O$ et d'angle $110°$.
\begin{enumerate}
\item Déterminer la nature et les éléments caractéristiques de $r_2 \circ r_1$.
\item Que vaut $r_1 \circ r_2$ ? Justifier.
\end{enumerate}
\tcblower
\textbf{Solution.}
\begin{enumerate}
\item Les rotations $r_1$ et $r_2$ ont le même centre $O$. D'après la propriété, la composée de deux rotations de même centre est une rotation de ce centre, dont l'angle est la somme des angles orientés :
\[ \alpha + \beta = 70° + 110° = 180°. \]
Une rotation d'angle $180°$ est une symétrie centrale.

\textbf{Conclusion :} $r_2 \circ r_1$ est la \textbf{symétrie centrale de centre $O$}.

\item La somme des angles est commutative : $110° + 70° = 180°$. Donc $r_1 \circ r_2$ est aussi la symétrie centrale de centre $O$. On a $r_1 \circ r_2 = r_2 \circ r_1$ : pour deux rotations de même centre, la composition est commutative.
\end{enumerate}
\end{exoresolu}

\begin{exercice}[Application directe]
Soient $r_1$ la rotation de centre $I$ d'angle $-120°$ et $r_2$ la rotation de centre $I$ d'angle $240°$. Déterminer $r_2 \circ r_1$ et $r_1 \circ r_2$.
\end{exercice}
\begin{corrige}[Exercice 1]
$r_2 \circ r_1$ : angle $-120° + 240° = 120°$. C'est la rotation de centre $I$, angle $120°$.
$r_1 \circ r_2$ : angle $240° + (-120°) = 120°$. Même rotation. Commutativité vérifiée.
\end{corrige}

\courssub{Composée de deux rotations de centres distincts}

\begin{propriete}[Composée de deux rotations de centres distincts]
Soient $r_1$ de centre $A$, angle $\alpha$ et $r_2$ de centre $B \neq A$, angle $\beta$. Soit $f = r_2 \circ r_1$.
\begin{itemize}
\item \textbf{Cas 1 :} Si $\alpha + \beta \not\equiv 0 \pmod{360°}$, alors $f$ est une \cle{rotation} d'angle $\alpha + \beta$. Son centre $\Omega$ est l'unique point fixe de $f$ ($f(\Omega)=\Omega$).
\item \textbf{Cas 2 :} Si $\alpha + \beta \equiv 0 \pmod{360°}$, alors $f$ est une \cle{translation}. Son vecteur est $\overrightarrow{A\, r_2(A)}$ (image de $A$ par $r_2$).
\end{itemize}
\textbf{Démonstration (idée) :} Une rotation se décompose en deux symétries axiales. La composée de quatre symétries axiales se réduit à deux symétries (rotation) ou à une translation si les axes finaux sont parallèles (somme des angles nulle).
\end{propriete}

\begin{methode}[Composer deux rotations de centres distincts]
Soit $r_1$ de centre $A$ et d'angle $\alpha$, et $r_2$ de centre $B$ ($B \neq A$) et d'angle $\beta$. Pour déterminer $f = r_2 \circ r_1$ :
\begin{enumerate}
\item Calculer la somme des angles orientés $S = \alpha + \beta$ (modulo $360°$).
\item \textbf{Si $S \not\equiv 0 \pmod{360°}$ :} $f$ est une \textbf{rotation} d'angle $S$.
      \begin{itemize}
      \item Pour trouver son centre $\Omega$ (point fixe), on résout $f(\Omega)=\Omega$.
      \item Méthode pratique : calculer $A' = f(A) = r_2(A)$. Le centre $\Omega$ est l'intersection de la médiatrice de $[AA']$ avec la droite passant par $A$ formant un angle $\frac{S}{2}$ avec $(AA')$ (côté choisi selon le sens de $S$).
      \end{itemize}
\item \textbf{Si $S \equiv 0 \pmod{360°}$ :} $f$ est une \textbf{translation}.
      \begin{itemize}
      \item Son vecteur est $\vec{v} = \overrightarrow{A\, f(A)} = \overrightarrow{A\, r_2(A)}$.
      \end{itemize}
\item Conclure en précisant la nature et tous les éléments caractéristiques (centre et angle, ou vecteur).
\end{enumerate}
\textbf{Réflexe :} toujours tester la somme des angles \textbf{avant} toute construction géométrique.
\end{methode}

\begin{popfigure}
\begin{tikzpicture}[scale=0.9, >=Stealth]
\coordinate (A) at (0,0);
\coordinate (B) at (4,0);
\coordinate (A1) at (1.5, 2.6); % r2(A) approx rotation 60 around B? No, r1 center A, r2 center B.
% Let's draw a generic case: r1 center A angle 60, r2 center B angle 60 -> sum 120 rotation.
% A' = r2(r1(A)) = r2(A).
\coordinate (Aprime) at (5.5, 1.5); % arbitrary image of A by r2
\coordinate (Omega) at (2.5, 1.5); % arbitrary center
\draw[thick, popblue] (A) -- (Aprime) node[midway, above] {$AA'$};
\draw[dashed, popmuted] (A) -- (Omega) node[midway, left] {$\Omega A$};
\draw[dashed, popmuted] (Aprime) -- (Omega) node[midway, right] {$\Omega A'$};
\draw[poporange, thick] (Omega) circle (0.1) node[below] {$\Omega$};
\draw[popblue, thick] (A) circle (0.1) node[below left] {$A$};
\draw[popblue, thick] (Aprime) circle (0.1) node[above right] {$A'=f(A)$};
\draw[popgreen, thick] (B) circle (0.1) node[below] {$B$};
% Angle mark
\draw[poppurple, ->] ($(A)!0.8!(Omega)$) arc (200:260:0.8) node[midway, below left] {$S/2$};
\end{tikzpicture}
\legende{Construction du centre $\Omega$ pour $S \neq 0$ : $\Omega$ sur la médiatrice de $[AA']$ et $(\Omega A, \Omega A') = S$.}
\end{popfigure}

\begin{exoresolu}[Composée de deux rotations de centres distincts -- Translation]
Sur un chantier, un gabarit triangulaire subit la rotation $r_1$ de centre $A$ et d'angle $60°$, puis la rotation $r_2$ de centre $B$ ($B \neq A$) et d'angle $-60°$.
\begin{enumerate}
\item Déterminer la nature de $f = r_2 \circ r_1$.
\item Exprimer le vecteur de cette transformation à l'aide de $r_2$ et du point $A$.
\end{enumerate}
\tcblower
\textbf{Solution.}
\begin{enumerate}
\item La somme des angles orientés est :
\[ \alpha + \beta = 60° + (-60°) = 0°. \]
Les centres $A$ et $B$ étant distincts et la somme des angles étant nulle (modulo $360°$), $f = r_2 \circ r_1$ est une \textbf{translation}.

\item Cherchons son vecteur en calculant l'image du point $A$ (centre de $r_1$). Comme $A$ est le centre de $r_1$, on a $r_1(A) = A$. Donc :
\[ f(A) = r_2(r_1(A)) = r_2(A). \]
Posons $A' = r_2(A)$. Alors $f(A) = A'$, donc le vecteur de la translation est $\overrightarrow{AA'}$ où $A'$ est l'image de $A$ par $r_2$.

\textbf{Conclusion :} $f$ est la translation de vecteur $\overrightarrow{AA'}$, avec $A' = r_2(A)$.
\end{enumerate}
\end{exoresolu}

\begin{exoresolu}[Composée de deux rotations de centres distincts -- Rotation]
Soit $r_1$ la rotation de centre $A$ et d'angle $90°$, et $r_2$ la rotation de centre $B$ ($B \neq A$) et d'angle $90°$. Déterminer la nature, l'angle et le centre de $f = r_2 \circ r_1$.
\tcblower
\textbf{Solution.}

La somme des angles orientés est :
\[ \alpha + \beta = 90° + 90° = 180° \not\equiv 0 \pmod{360°}. \]
La somme n'étant pas nulle, $f$ est une \textbf{rotation d'angle $180°$}, c'est-à-dire une \textbf{symétrie centrale}.

Son centre $\Omega$ est l'unique point fixe de $f$ : $f(\Omega) = \Omega$.
On construit $A' = f(A) = r_2(r_1(A)) = r_2(A)$ (car $r_1(A)=A$).
Comme $f$ est une symétrie centrale, elle échange $A$ et $A'$ ; son centre $\Omega$ est donc le \textbf{milieu du segment $[AA']$}.

\textbf{Conclusion :} $f$ est la symétrie centrale de centre $\Omega$, milieu de $[A\, r_2(A)]$.
\end{exoresolu}

\begin{exercice}[Rotation composée]
Soient $r_1$ rotation de centre $C_1$, angle $40°$ et $r_2$ rotation de centre $C_2 \neq C_1$, angle $140°$.
\begin{enumerate}
\item Nature de $f = r_2 \circ r_1$ ?
\item Nature de $g = r_1 \circ r_2$ ? Les angles sont-ils égaux ? Les centres sont-ils les mêmes ?
\end{enumerate}
\end{exercice}
\begin{corrige}[Exercice 2]
1. Somme $= 180° \neq 0$. $f$ est symétrie centrale. Centre $\Omega_f$ = milieu de $[C_1\, r_2(C_1)]$.
2. Somme $= 180° \neq 0$. $g$ est symétrie centrale. Angle identique ($180°$). Centre $\Omega_g$ = milieu de $[C_2\, r_1(C_2)]$. En général $\Omega_f \neq \Omega_g$ : non-commutativité des centres.
\end{corrige}

\courssub{Composée de deux homothéties de même centre}

\begin{propriete}[Composée de deux homothéties de même centre]
Soient $h_1$ et $h_2$ deux homothéties de \textbf{même centre} $\Omega$ et de rapports respectifs $k_1$ et $k_2$ ($k_1, k_2 \in \mathbb{R}^*$).
La composée $h_2 \circ h_1$ est :
\begin{itemize}
\item l'\cle{homothétie} de centre $\Omega$ et de rapport $k = k_1 k_2$ si $k_1 k_2 \neq 1$ ;
\item l'\cle{identité} du plan si $k_1 k_2 = 1$ (les rapports sont inverses).
\end{itemize}
\textbf{Démonstration exigible :} Pour tout point $M$, $M' = h_1(M) \Rightarrow \overrightarrow{\Omega M'} = k_1 \overrightarrow{\Omega M}$. $M'' = h_2(M') \Rightarrow \overrightarrow{\Omega M''} = k_2 \overrightarrow{\Omega M'} = k_2(k_1 \overrightarrow{\Omega M}) = (k_1 k_2) \overrightarrow{\Omega M}$.
Si $k_1 k_2 \neq 1$, c'est une homothétie de centre $\Omega$, rapport $k_1 k_2$.
Si $k_1 k_2 = 1$, $\overrightarrow{\Omega M''} = \overrightarrow{\Omega M} \Rightarrow M''=M$ pour tout $M$ : c'est l'identité.
\end{propriete}

\begin{methode}[Composer deux homothéties de même centre]
Soit $h_1$ et $h_2$ deux homothéties de \textbf{même centre} $\Omega$ et de rapports respectifs $k_1$ et $k_2$. Pour déterminer $h_2 \circ h_1$ :
\begin{enumerate}
\item Vérifier que les deux homothéties ont le même centre $\Omega$.
\item Multiplier les rapports : $k = k_1 \times k_2$.
\item \textbf{Cas 1 :} si $k \neq 1$, alors $h_2 \circ h_1$ est l'\cle{homothétie} de centre $\Omega$ et de rapport $k_1 k_2$.
\item \textbf{Cas 2 :} si $k_1 k_2 = 1$ (rapports inverses), alors $h_2 \circ h_1$ est l'\cle{identité} du plan.
\item Rédiger la conclusion avec le centre et le rapport (ou « identité »).
\end{enumerate}
\textbf{Formule à retenir :} $\forall M,\quad \overrightarrow{\Omega M''} = k_1 k_2 \, \overrightarrow{\Omega M}$.
\end{methode}

\begin{exoresolu}[Agrandissement puis réduction d'un plan d'atelier]
Un dessinateur industriel reproduit le plan d'une pièce en appliquant d'abord l'homothétie $h_1$ de centre $O$ et de rapport $3$, puis l'homothétie $h_2$ de même centre $O$ et de rapport $\dfrac{1}{2}$.
\begin{enumerate}
\item Déterminer la nature et les éléments caractéristiques de $h_2 \circ h_1$.
\item Un segment du plan initial mesure \SI{4}{cm}. Quelle est la longueur du segment image par $h_2 \circ h_1$ ?
\end{enumerate}
\tcblower
\textbf{Solution.}
\begin{enumerate}
\item Les deux homothéties ont le même centre $O$. Le rapport de la composée est le produit des rapports :
\[ k = k_1 \times k_2 = 3 \times \frac{1}{2} = \frac{3}{2}. \]
Comme $k = \dfrac{3}{2} \neq 1$, la composée est une homothétie.

\textbf{Conclusion :} $h_2 \circ h_1$ est l'homothétie de centre $O$ et de rapport $\dfrac{3}{2}$.

\item Une homothétie de rapport $k$ multiplie les longueurs par $|k|$. Ici $|k| = \dfrac{3}{2}$, donc :
\[ L' = \frac{3}{2} \times 4 = 6. \]
Le segment image mesure \SI{6}{cm}.
\end{enumerate}
\end{exoresolu}

\begin{exoresolu}[Cas où la composée est l'identité]
Soit $h_1$ l'homothétie de centre $\Omega$ et de rapport $-2$, et $h_2$ l'homothétie de même centre $\Omega$ et de rapport $-\dfrac{1}{2}$. Déterminer $h_2 \circ h_1$.
\tcblower
\textbf{Solution.}

Les deux homothéties ont le même centre $\Omega$. Le rapport de la composée est :
\[ k = k_1 \times k_2 = (-2) \times \left(-\frac{1}{2}\right) = 1. \]
Comme $k = 1$, pour tout point $M$ d'image $M''$ on a $\overrightarrow{\Omega M''} = 1 \times \overrightarrow{\Omega M}$, c'est-à-dire $M'' = M$.

\textbf{Conclusion :} $h_2 \circ h_1$ est l'\textbf{identité du plan} (les deux homothéties sont réciproques l'une de l'autre).
\end{exoresolu}

\begin{exercice}[Homothéties successives]
Une pièce est photographiée avec un agrandissement de rapport $k_1 = 1,5$ (centre $O$), puis le tirage est réduit par une homothétie de même centre $O$ et de rapport $k_2 = \frac{2}{3}$.
\begin{enumerate}
\item Nature de la transformation globale ?
\item Si un trou de diamètre \SI{10}{mm} sur la pièce paraît de \SI{15}{mm} sur la photo intermédiaire, quel est son diamètre sur le tirage final ?
\end{enumerate}
\end{exercice}
\begin{corrige}[Exercice 3]
1. $k = 1,5 \times \frac{2}{3} = 1$. Composée = Identité. Le tirage final est identique à la pièce originale (même taille, même orientation).
2. Diamètre final = \SI{10}{mm}.
\end{corrige}

\courssub{Composée d'une homothétie et d'une translation}

\begin{propriete}[Composée d'une homothétie et d'une translation]
Soit $h$ une homothétie de centre $\Omega$ et de rapport $k \neq 1$, et $t$ une translation de vecteur $\vec{u}$.
Les composées $t \circ h$ et $h \circ t$ sont des \cle{homothéties} de même rapport $k$, mais de centres \textbf{différents} de $\Omega$ (et différents l'un de l'autre en général).
Si $k = 1$, $h = \text{Id}$ et la composée est la translation $t$.
\end{propriete}

\begin{methode}[Composer une homothétie et une translation]
Soit $h$ une homothétie de centre $\Omega$ et de rapport $k$, et $t$ une translation de vecteur $\vec{u}$. Pour déterminer $f = t \circ h$ (ou $h \circ t$) :
\begin{enumerate}
\item Regarder la valeur du rapport $k$.
\item \textbf{Si $k \neq 1$ :} la composée est une \textbf{homothétie de rapport $k$}. Son centre $\Omega'$ est l'unique point fixe : on résout $f(\Omega') = \Omega'$.
      \begin{itemize}
      \item Pour $f = t \circ h$ : $\overrightarrow{\Omega \Omega'} = k \overrightarrow{\Omega \Omega'} + \vec{u} \;\Rightarrow\; (1-k)\overrightarrow{\Omega \Omega'} = \vec{u} \;\Rightarrow\; \overrightarrow{\Omega \Omega'} = \frac{1}{1-k}\vec{u}$.
      \item Pour $f = h \circ t$ : $\overrightarrow{\Omega \Omega'} = k(\overrightarrow{\Omega \Omega'} + \vec{u}) \;\Rightarrow\; (1-k)\overrightarrow{\Omega \Omega'} = k\vec{u} \;\Rightarrow\; \overrightarrow{\Omega \Omega'} = \frac{k}{1-k}\vec{u}$.
      \end{itemize}
\item \textbf{Si $k = 1$ :} $h = \text{Id}$, la composée est la translation $t$.
\item Conclure : nature, rapport, centre (ou vecteur).
\end{enumerate}
\textbf{Attention :} en général $t \circ h \neq h \circ t$ : la composition n'est pas commutative. Les centres obtenus sont différents.
\end{methode}

\begin{popfigure}
\begin{tikzpicture}[scale=0.8, >=Stealth]
\coordinate (O) at (0,0);
\coordinate (u) at (2,1);
\coordinate (OmegaPrime) at (-2,-1); % For k=2, t oh : OOmega' = -u
\draw[->, thick, popblue] (O) -- (u) node[midway, above right] {$\vec{u}$};
\draw[->, thick, poporange] (O) -- (OmegaPrime) node[midway, below left] {$\overrightarrow{O\Omega'}=-\vec{u}$};
\draw[popblue, thick] (O) circle (0.1) node[above right] {$O$};
\draw[poporange, thick] (OmegaPrime) circle (0.1) node[below left] {$\Omega'$};
\draw[dashed, popmuted] (O) -- ++(3,1.5) node[above right] {Direction $\vec{u}$};
\draw[dashed, popmuted] (OmegaPrime) -- ++(3,1.5);
\end{tikzpicture}
\legende{Cas $f = t \circ h$, $k=2$ : centre $\Omega'$ tel que $\overrightarrow{O\Omega'} = -\vec{u}$.}
\end{popfigure}

\begin{exoresolu}[Homothétie suivie d'une translation -- Atelier de soudure]
Dans un atelier de soudure, un motif est reproduit par l'homothétie $h$ de centre $O$ et de rapport $2$, puis déplacé par la translation $t$ de vecteur $\vec{u}$. On note $f = t \circ h$.
\begin{enumerate}
\item Quelle est la nature de $f$ ? Préciser son rapport.
\item Soit $A$ un point, $A' = h(A)$ et $A'' = f(A)$. Exprimer $\overrightarrow{OA''}$ en fonction de $\overrightarrow{OA}$ et $\vec{u}$.
\item Montrer que $f$ admet un unique point fixe $\Omega'$ et exprimer $\overrightarrow{O\Omega'}$ en fonction de $\vec{u}$.
\end{enumerate}
\tcblower
\textbf{Solution.}
\begin{enumerate}
\item Le rapport de $h$ est $k = 2 \neq 1$. La composée d'une homothétie de rapport $k \neq 1$ et d'une translation est une \textbf{homothétie de même rapport $2$} (mais de centre différent de $O$ en général).

\item Par définition de $h$ : $\overrightarrow{OA'} = 2\,\overrightarrow{OA}$. Par définition de $t$ : $\overrightarrow{A'A''} = \vec{u}$. D'après la relation de Chasles :
\[ \overrightarrow{OA''} = \overrightarrow{OA'} + \overrightarrow{A'A''} = 2\,\overrightarrow{OA} + \vec{u}. \]

\item $\Omega'$ est un point fixe de $f$ si $f(\Omega') = \Omega'$, c'est-à-dire, d'après la question précédente :
\[ \overrightarrow{O\Omega'} = 2\,\overrightarrow{O\Omega'} + \vec{u}. \]
En passant $2\,\overrightarrow{O\Omega'}$ au premier membre :
\[ \overrightarrow{O\Omega'} - 2\,\overrightarrow{O\Omega'} = \vec{u} \quad\Longleftrightarrow\quad -\overrightarrow{O\Omega'} = \vec{u} \quad\Longleftrightarrow\quad \overrightarrow{O\Omega'} = -\vec{u}. \]
Ce point $\Omega'$ existe et est unique car l'équation vectorielle a une unique solution.

\textbf{Conclusion :} $f$ est l'homothétie de rapport $2$ et de centre $\Omega'$ défini par $\overrightarrow{O\Omega'} = -\vec{u}$.
\end{enumerate}
\end{exoresolu}

\begin{exercice}[Ordre inverse : translation puis homothétie]
Mêmes $h$ (centre $O$, rapport $2$) et $t$ (vecteur $\vec{u}$). On considère $g = h \circ t$.
\begin{enumerate}
\item Nature et rapport de $g$ ?
\item Déterminer le centre $\Omega''$ de $g$ (exprimer $\overrightarrow{O\Omega''}$ en fonction de $\vec{u}$).
\item Comparer $\Omega'$ (de $t \circ h$) et $\Omega''$.
\end{enumerate}
\end{exercice}
\begin{corrige}[Exercice 4]
1. $g$ est homothétie de rapport $2$.
2. Point fixe : $\overrightarrow{O\Omega''} = 2(\overrightarrow{O\Omega''} + \vec{u}) \Rightarrow -\overrightarrow{O\Omega''} = 2\vec{u} \Rightarrow \overrightarrow{O\Omega''} = -2\vec{u}$.
3. $\Omega'$ tel que $\overrightarrow{O\Omega'} = -\vec{u}$. $\Omega''$ tel que $\overrightarrow{O\Omega''} = -2\vec{u}$. $\Omega' \neq \Omega''$ (sauf si $\vec{u}=\vec{0}$). Non-commutativité confirmée.
\end{corrige}

\courssub{Synthèse, méthodes et applications concrètes}

\begin{aretenir}[Tableau récapitulatif des composées (Programme 1re Tech)]
\begin{center}
\begin{tabularx}{\linewidth}{|l|X|X|}\hline
\rowcolor{popblueL}
\textbf{Composée} & \textbf{Condition} & \textbf{Résultat} \\ \hline
$r_2 \circ r_1$ (même centre $\Omega$) & Toujours & Rotation centre $\Omega$, angle $\alpha+\beta$ \\ \hline
$r_2 \circ r_1$ (centres $A \neq B$) & $\alpha+\beta \not\equiv 0 \pmod{360°}$ & Rotation angle $\alpha+\beta$, centre $\Omega$ (point fixe) \\ \hline
$r_2 \circ r_1$ (centres $A \neq B$) & $\alpha+\beta \equiv 0 \pmod{360°}$ & Translation vecteur $\overrightarrow{A\, r_2(A)}$ \\ \hline
$h_2 \circ h_1$ (même centre $\Omega$) & $k_1 k_2 \neq 1$ & Homothétie centre $\Omega$, rapport $k_1 k_2$ \\ \hline
$h_2 \circ h_1$ (même centre $\Omega$) & $k_1 k_2 = 1$ & Identité \\ \hline
$t \circ h$ ou $h \circ t$ ($h$ centre $\Omega$, rapport $k$) & $k \neq 1$ & Homothétie rapport $k$, centre $\Omega' \neq \Omega$ \\ \hline
$t \circ h$ ou $h \circ t$ & $k = 1$ & Translation $t$ \\ \hline
\end{tabularx}
\end{center}
\end{aretenir}

\begin{methode}[Résoudre un problème concret avec les transformations]
Face à une situation concrète (atelier, chantier, plan, pignon, engrenage) :
\begin{enumerate}
\item \textbf{Traduire} l'énoncé en langage mathématique : identifier chaque transformation (centre, angle orienté, rapport, vecteur) et l'ordre de composition ($f \circ g$ = $g$ puis $f$).
\item \textbf{Identifier} la nature de la composée grâce aux résultats du cours (somme des angles, produit des rapports, test $k=1$).
\item \textbf{Calculer} les éléments demandés :
      \begin{itemize}
      \item Longueurs : multipliées par $|k|$ (homothétie) ou conservées (rotation, translation).
      \item Angles : conservés par toutes ces transformations.
      \item Aires : multipliées par $k^2$ (homothétie de rapport $k$) ou conservées (rotation, translation).
      \end{itemize}
\item \textbf{Rédiger} la réponse par une phrase complète qui répond à la question posée, avec les unités.
\end{enumerate}
\textbf{Réflexes :} Une rotation ou une homothétie conserve les angles ; une homothétie de rapport $k$ multiplie les longueurs par $|k|$ et les aires par $k^2$.
\end{methode}

\begin{exoresolu}[Problème de chantier : double agrandissement d'un plan]
Pour la construction d'un hangar à Bafoussam, un technicien agrandit le plan au sol en appliquant l'homothétie $h_1$ de centre $O$ et de rapport $\dfrac{3}{2}$, puis son chef applique au résultat l'homothétie $h_2$ de même centre $O$ et de rapport $2$.
\begin{enumerate}
\item Déterminer la nature et le rapport de $h_2 \circ h_1$.
\item Sur le plan initial, une poutre est représentée par un segment de \SI{5}{cm} et l'aire du hangar est de \SI{12}{cm^2}. Calculer la longueur de la poutre et l'aire du hangar sur le plan final.
\end{enumerate}
\tcblower
\textbf{Solution.}
\begin{enumerate}
\item Les deux homothéties ont le même centre $O$. Le rapport de la composée est :
\[ k = \frac{3}{2} \times 2 = 3. \]
Comme $k = 3 \neq 1$, $h_2 \circ h_1$ est l'\textbf{homothétie de centre $O$ et de rapport $3$}.

\item Une homothétie de rapport $3$ multiplie les longueurs par $|3| = 3$ et les aires par $3^2 = 9$.

Longueur de la poutre sur le plan final :
\[ L' = 3 \times 5 = 15 \ \text{cm}. \]

Aire du hangar sur le plan final :
\[ A' = 9 \times 12 = 108 \ \text{cm}^2. \]

\textbf{Conclusion :} sur le plan final, la poutre mesure \SI{15}{cm} et l'aire du hangar est de \SI{108}{cm^2}.
\end{enumerate}
\end{exoresolu}

\begin{exoresolu}[Engrenages : rotations successives]
Dans une boîte de vitesses, un pignon $P_1$ (centre $O_1$) entraîne un pignon $P_2$ (centre $O_2$) qui entraîne un pignon $P_3$ (centre $O_3$). Le mouvement de $P_1$ vers $P_2$ correspond à une rotation $r_1$ de centre $O_1$, angle $60°$. Le mouvement de $P_2$ vers $P_3$ correspond à une rotation $r_2$ de centre $O_2$, angle $-60°$. On suppose $O_1 \neq O_2$.
Déterminer la transformation globale $f = r_2 \circ r_1$ qui relie directement la position initiale de $P_1$ à la position finale de $P_3$.
\tcblower
\textbf{Solution.}
Somme des angles : $60° + (-60°) = 0°$.
Centres distincts ($O_1 \neq O_2$).
D'après le cours (Cas 2), $f$ est une \textbf{translation}.
Son vecteur est $\overrightarrow{O_1\, r_2(O_1)}$.
\textbf{Interprétation :} L'axe du pignon $P_3$ est l'image de l'axe du pignon $P_1$ par une translation (pas de rotation nette, juste un déplacement).
\end{exoresolu}

\begin{savaistu}[Histoire et applications : Les transformations affines]
\begin{itemize}
\item Le terme « \emph{affine} » vient du latin \emph{affinis} (parent, voisin). Les transformations affines (rotations, translations, homothéties, symétries, et leurs composées) conservent le parallélisme et les rapports de longueurs sur une même droite.
\item En \textbf{infographie} et \textbf{CAO/DAO} (Conception Assistée par Ordinateur), tout placement d'objet (déplacement, rotation, mise à l'échelle) est géré par des matrices de transformation affine ($3 \times 3$ en coordonnées homogènes). La composée correspond au produit matriciel.
\item En \textbf{robotique} (bras articulés), la position de l'effecteur (pince, outil) est obtenue par la composée des rotations successives des articulations. C'est le \emph{modèle géométrique direct}.
\item Au \textbf{Cameroun}, les bureaux d'études (BTP, mécanique, topographie) utilisent ces notions pour le calage de plans, l'implantation de bâtiments (rotations + translations) et la mise à l'échelle de maquettes (homothéties).
\end{itemize}
\end{savaistu}

\begin{attention}[Les 5 erreurs qui coûtent des points au Probatoire]
\begin{enumerate}
\item \textbf{Additionner les rapports d'homothéties au lieu de les multiplier.} La composée de deux homothéties de même centre a pour rapport le \emph{produit} $k_1 k_2$, jamais la somme $k_1 + k_2$. De même, ce sont les \emph{angles} des rotations qui s'additionnent, pas les rapports.
\item \textbf{Oublier de tester la somme des angles avant de conclure.} Pour deux rotations de centres distincts, si $\alpha + \beta = 0°$ (mod $360°$), la composée est une \emph{translation}, pas une rotation. Affirmer « rotation d'angle nul » est une erreur de nature qui fait perdre tous les points de la question.
\item \textbf{Confondre le cas $k_1 k_2 = 1$.} Si le produit des rapports vaut $1$, la composée de deux homothéties de même centre est l'\emph{identité} (tout point est fixe), et non une homothétie de rapport $1$ de centre $\Omega$ : il faut le signaler explicitement (« c'est l'identité »).
\item \textbf{Multiplier les aires par $|k|$ au lieu de $k^2$.} Une homothétie de rapport $k$ multiplie les longueurs par $|k|$ mais les aires par $k^2$. Erreur classique dans les problèmes de plans et de surfaces.
\item \textbf{Négliger la rédaction et les conditions d'application.} Il faut toujours justifier : « les deux rotations ont le même centre, donc... » ou « la somme des angles est non nulle, donc... ». Une conclusion sans justification, ou sans préciser le centre et l'angle (ou le rapport), est sanctionnée même si le résultat est correct.
\end{enumerate}
\end{attention}

\newpage

\coursec{Exercices}

\courssub{Je vérifie mes connaissances}

\begin{exercice}[QCM : nature d'une composée]
Pour chaque question, une seule réponse est exacte. Recopier la lettre de la bonne réponse.
\begin{enumerate}
\item La composée de deux rotations de même centre $O$, d'angles $\alpha$ et $\beta$, est :
\begin{enumerate}
\item une translation ;
\item une rotation de centre $O$ et d'angle $\alpha+\beta$ ;
\item une homothétie de centre $O$.
\end{enumerate}
\item La composée de deux homothéties de même centre $O$, de rapports $k$ et $k'$, est :
\begin{enumerate}
\item une homothétie de centre $O$ et de rapport $k\times k'$ ;
\item une homothétie de centre $O$ et de rapport $k+k'$ ;
\item toujours une translation.
\end{enumerate}
\item Si $k\times k'=1$, la composée $h(O,k')\circ h(O,k)$ est :
\begin{enumerate}
\item la symétrie centrale de centre $O$ ;
\item l'identité du plan ;
\item une translation.
\end{enumerate}
\item La composée d'une homothétie de rapport $k\neq 1$ et d'une translation de vecteur non nul est :
\begin{enumerate}
\item une translation ;
\item une homothétie de rapport $k$ ;
\item une rotation.
\end{enumerate}
\end{enumerate}
\end{exercice}

\begin{exercice}[Vrai ou faux justifié]
Répondre par VRAI ou FAUX en justifiant chaque réponse.
\begin{enumerate}
\item La composée de deux rotations d'angles $30^{\circ}$ et $-30^{\circ}$, de centres quelconques, est toujours une translation.
\item La composée de deux rotations de centres distincts est toujours une rotation.
\item La composée de deux homothéties de même centre $O$ et de rapports $2$ et $-\dfrac{3}{4}$ est une homothétie de rapport $-\dfrac{3}{2}$.
\item La composée d'une homothétie de rapport $1$ et d'une translation est une homothétie.
\end{enumerate}
\end{exercice}

\begin{exercice}[Composée de deux rotations de même centre]
Soient $r_1$ et $r_2$ deux rotations de même centre $O$, d'angles respectifs $50^{\circ}$ et $70^{\circ}$.
\begin{enumerate}
\item Déterminer la nature et les éléments caractéristiques de $r_2\circ r_1$.
\item On note $r_3$ la rotation de centre $O$ et d'angle $-50^{\circ}$. Déterminer la nature et les éléments caractéristiques de $r_3\circ r_1$. Que remarque-t-on ?
\end{enumerate}
\end{exercice}

\begin{exercice}[Composée de deux homothéties de même centre]
Dans un repère $(O\,;\vec{\imath},\vec{\jmath})$, on considère les homothéties $h_1$ de centre $O$ et de rapport $2$, et $h_2$ de centre $O$ et de rapport $-\dfrac{3}{4}$.
\begin{enumerate}
\item Déterminer la nature et les éléments caractéristiques de $h_2\circ h_1$.
\item Calculer les coordonnées de l'image du point $A(4\,;-2)$ par $h_2\circ h_1$.
\end{enumerate}
\end{exercice}

\courssub{Je m'entraîne}

\begin{exercice}[Rotations de centres distincts]
Soient $r_1$ la rotation de centre $O_1$ et d'angle $60^{\circ}$, et $r_2$ la rotation de centre $O_2$ et d'angle $-60^{\circ}$, avec $O_1\neq O_2$.
\begin{enumerate}
\item Quelle est la somme des angles de $r_1$ et $r_2$ ? En déduire la nature de $r_2\circ r_1$.
\item Que vaut $r_1(O_1)$ ? Justifier.
\item Déterminer l'image de $O_1$ par $r_2\circ r_1$ et en déduire le vecteur de la translation $r_2\circ r_1$.
\end{enumerate}
\end{exercice}

\begin{exercice}[Homothétie et translation dans un repère]
Dans un repère $(O\,;\vec{\imath},\vec{\jmath})$, on considère l'homothétie $h$ de centre $O$ et de rapport $3$, et la translation $t$ de vecteur $\vec{u}(-6\,;3)$.
\begin{enumerate}
\item Soit $M(x\,;y)$ un point du plan. Calculer les coordonnées de $M'=h(M)$, puis celles de $M''=t(M')$.
\item Montrer que $f=t\circ h$ est une homothétie de rapport $3$.
\item Déterminer les coordonnées du centre $\Omega$ de $f$ (on cherchera le point invariant par $f$).
\end{enumerate}
\end{exercice}

\begin{exercice}[Problème inverse]
La composée de deux homothéties $h_1$ et $h_2$ de même centre $O$ est l'homothétie de centre $O$ et de rapport $-6$. On sait que $h_1$ a pour rapport $2$.
\begin{enumerate}
\item Déterminer le rapport de $h_2$.
\item Interpréter géométriquement l'effet de $h_2$ sur une figure (agrandissement ou réduction, sens).
\end{enumerate}
\end{exercice}

\begin{exercice}[Plan d'une concession]
Un topographe réalise le plan d'une concession. Le plan est agrandi par une homothétie de rapport $\dfrac{5}{4}$, puis le résultat est agrandi par une homothétie de même centre et de rapport $\dfrac{6}{5}$.
\begin{enumerate}
\item Quel est le rapport de l'agrandissement global ? Justifier.
\item L'aire de la concession sur le plan initial est de $800\ \mathrm{cm}^2$. Quelle est son aire sur le plan final ?
\end{enumerate}
\end{exercice}

\courssub{Je me prépare au Probatoire}

\begin{exercice}[Image d'un carré — type Probatoire]
$ABCD$ est un carré de centre $O$, de côté $4\ \mathrm{cm}$, tel que $\left(\overrightarrow{OA},\overrightarrow{OB}\right)=90^{\circ}$. On considère la rotation $r$ de centre $O$ et d'angle $90^{\circ}$, et l'homothétie $h$ de centre $A$ et de rapport $2$.
\begin{enumerate}
\item Faire une figure que l'on complétera au fur et à mesure.
\item Déterminer $r(A)$, $r(B)$, $r(C)$ et $r(D)$. Justifier.
\item Construire les points $A'$, $B'$, $C'$, $D'$, images respectives de $A$, $B$, $C$, $D$ par $h\circ r$.
\item Quelle est la nature du quadrilatère $A'B'C'D'$ ? Justifier et calculer son aire.
\item La transformation $h\circ r$ est-elle une homothétie, une rotation, une translation ? Justifier la réponse.
\end{enumerate}
\end{exercice}

\begin{exercice}[Atelier de mécanique — type Probatoire]
Dans un atelier de mécanique de Douala, une pièce métallique est modélisée dans un repère orthonormé $(O\,;\vec{\imath},\vec{\jmath})$ (unité : le centimètre). La pièce subit deux opérations successives : d'abord l'homothétie $h_1$ de centre $O$ et de rapport $\dfrac{1}{2}$ (réduction pour le plan), puis l'homothétie $h_2$ de centre $O$ et de rapport $-4$ (mise à l'échelle de fabrication avec retournement).
\begin{enumerate}
\item Déterminer la nature et les éléments caractéristiques de la transformation $f=h_2\circ h_1$.
\item Un trou de perçage est représenté par le point $P(6\,;-2)$. Calculer les coordonnées de $P'=f(P)$.
\item Une arête de la pièce a pour longueur $5\ \mathrm{cm}$. Quelle est la longueur de son image par $f$ ? Justifier.
\item L'aire de la face de la pièce est $12\ \mathrm{cm}^2$. Quelle est l'aire de son image par $f$ ? Justifier.
\item Le chef d'atelier affirme : « $f$ est une symétrie centrale suivie d'un agrandissement ». A-t-il raison ? Justifier.
\end{enumerate}
\end{exercice}

\begin{exercice}[Chantier et déplacement d'une structure — type Probatoire]
Sur un chantier, une structure triangulaire $ABC$ est modélisée dans un repère orthonormé $(O\,;\vec{\imath},\vec{\jmath})$ par les points $A(1\,;1)$, $B(4\,;1)$ et $C(1\,;3)$. On considère l'homothétie $h$ de centre $O$ et de rapport $2$, et la translation $t$ de vecteur $\vec{u}(-4\,;2)$. On pose $f=t\circ h$.
\begin{enumerate}
\item Calculer les coordonnées des points $A'$, $B'$, $C'$, images de $A$, $B$, $C$ par $f$.
\item Montrer que pour tout point $M(x\,;y)$, son image $M''$ par $f$ a pour coordonnées $(2x-4\,;\,2y+2)$.
\item En déduire que $f$ est une homothétie dont on précisera le rapport, et déterminer les coordonnées de son centre $\Omega$.
\item Calculer l'aire du triangle $ABC$, puis celle du triangle $A'B'C'$. Vérifier le résultat à l'aide du rapport de $f$.
\item Sur le chantier, on déplace ensuite la structure $A'B'C'$ par la translation $t'$ de vecteur $\vec{v}(4\,;-2)$. Quelle est la nature de $t'\circ f$ ? Déterminer ses éléments caractéristiques.
\end{enumerate}
\end{exercice}

\begin{attention}[Erreurs fréquentes]
\begin{itemize}
\item Ne pas \emph{additionner} les rapports de deux homothéties : les rapports se \textbf{multiplient} ($k\times k'$), ce sont les \emph{angles} des rotations qui s'additionnent.
\item Pour l'effet sur les longueurs, c'est $|k|$ qui intervient, pas $k$ : une homothétie de rapport $-2$ double les longueurs.
\item Pour les aires, on multiplie par $k^2$ (et non par $k$).
\item La composée d'une homothétie de rapport $k\neq 1$ et d'une translation est une homothétie de \textbf{même rapport} $k$, mais son centre n'est \emph{pas} celui de l'homothétie initiale : il faut chercher le point invariant.
\item Dans $t\circ h$, on applique d'abord $h$, ensuite $t$ : l'ordre de lecture est inverse de l'ordre d'application.
\end{itemize}
\end{attention}

\newpage

\coursec{Corrigés des exercices}

\coursec{Corrigés détaillés — Transformations affines du plan}

\begin{corrige}[Exercice 1 — QCM : nature d'une composée]
\begin{enumerate}
\item Réponse \textbf{b}. La composée de deux rotations de même centre $O$, d'angles $\alpha$ et $\beta$, est la rotation de centre $O$ et d'angle $\alpha+\beta$ : le centre est conservé et les angles s'additionnent.
\item Réponse \textbf{a}. La composée de deux homothéties de même centre $O$, de rapports $k$ et $k'$, est l'homothétie de centre $O$ et de rapport $k\times k'$ : les rapports se multiplient.
\item Réponse \textbf{b}. Si $k\times k'=1$, alors $h(O,k')\circ h(O,k)$ est l'homothétie de centre $O$ et de rapport $1$, c'est-à-dire l'\cle{identité} du plan : chaque point est ramené à sa position initiale.
\item Réponse \textbf{b}. La composée d'une homothétie de rapport $k\neq 1$ et d'une translation est une homothétie de \textbf{même rapport} $k$ (son centre est en général différent de celui de l'homothétie initiale).
\end{enumerate}
\end{corrige}

\begin{corrige}[Exercice 2 — Vrai ou faux justifié]
\begin{enumerate}
\item \textbf{VRAI.} La somme des angles est $30^{\circ}+(-30^{\circ})=0^{\circ}$. Or la composée de deux rotations dont la somme des angles est nulle est une \cle{translation}, quels que soient les centres.
\item \textbf{FAUX.} La composée de deux rotations de centres distincts est une rotation \emph{uniquement} si la somme des angles est non nulle ; si cette somme vaut $0^{\circ}$, la composée est une translation. Contre-exemple : deux rotations d'angles $60^{\circ}$ et $-60^{\circ}$ de centres distincts donnent une translation.
\item \textbf{VRAI.} Les rapports se multiplient :
\[ 2\times\left(-\dfrac{3}{4}\right)=-\dfrac{6}{4}=-\dfrac{3}{2}. \]
La composée est donc l'homothétie de centre $O$ et de rapport $-\dfrac{3}{2}$.
\item \textbf{FAUX.} Une homothétie de rapport $1$ est l'identité du plan. Sa composée avec une translation est donc cette \cle{translation} elle-même, et non une homothétie (une translation non triviale n'a aucun point invariant, contrairement à une homothétie).
\end{enumerate}
\end{corrige}

\begin{corrige}[Exercice 3 — Composée de deux rotations de même centre]
\begin{enumerate}
\item La composée de deux rotations de même centre $O$ est la rotation de centre $O$ dont l'angle est la somme des angles :
\[ r_2\circ r_1 = r\big(O\,;\,50^{\circ}+70^{\circ}\big)=r\big(O\,;\,120^{\circ}\big). \]
C'est donc la \cle{rotation} de centre $O$ et d'angle $120^{\circ}$.
\item De même :
\[ r_3\circ r_1 = r\big(O\,;\,50^{\circ}+(-50^{\circ})\big)=r\big(O\,;\,0^{\circ}\big). \]
Une rotation d'angle nul laisse tout point invariant : c'est l'\cle{identité} du plan. On remarque que $r_3$ est la \emph{rotation réciproque} de $r_1$ : composer une rotation avec sa réciproque « annule » le déplacement, comme avancer puis reculer du même pas.
\end{enumerate}
\end{corrige}

\begin{corrige}[Exercice 4 — Composée de deux homothéties de même centre]
\begin{enumerate}
\item La composée de deux homothéties de même centre $O$ est l'homothétie de centre $O$ dont le rapport est le produit des rapports :
\[ h_2\circ h_1 = h\left(O\,;\,2\times\left(-\dfrac{3}{4}\right)\right)=h\left(O\,;\,-\dfrac{3}{2}\right). \]
\item Pour une homothétie de centre $O$ (origine du repère) et de rapport $k$, l'image de $M(x\,;y)$ est $M'(kx\,;ky)$. Ici $k=-\dfrac{3}{2}$ et $A(4\,;-2)$ :
\[ A'\left(-\dfrac{3}{2}\times 4\ ;\ -\dfrac{3}{2}\times(-2)\right)=\left(-\dfrac{12}{2}\ ;\ \dfrac{6}{2}\right)=(-6\,;\,3). \]
Vérification : $\overrightarrow{OA'}=-\dfrac{3}{2}\,\overrightarrow{OA}$, ce qui correspond bien aux coordonnées trouvées.
\end{enumerate}
\end{corrige}

\begin{corrige}[Exercice 5 — Rotations de centres distincts]
\begin{enumerate}
\item La somme des angles est $60^{\circ}+(-60^{\circ})=0^{\circ}$. La somme étant nulle, $r_2\circ r_1$ est une \cle{translation}.
\item $r_1(O_1)=O_1$ car le centre d'une rotation est \emph{invariant} par cette rotation (c'est le seul point fixe d'une rotation d'angle non nul).
\item On calcule l'image de $O_1$ par la composée :
\[ r_2\circ r_1(O_1)=r_2\big(r_1(O_1)\big)=r_2(O_1). \]
Notons $O_1'=r_2(O_1)$, image de $O_1$ par la rotation de centre $O_2$ et d'angle $-60^{\circ}$. Comme $O_1\neq O_2$, le point $O_1$ n'est pas le centre de $r_2$, donc $O_1'\neq O_1$. Une translation est entièrement déterminée par l'image d'un seul point : le vecteur de la translation $r_2\circ r_1$ est donc
\[ \vec{u}=\overrightarrow{O_1O_1'}. \]
\end{enumerate}
\textbf{Point de vigilance :} pour identifier le vecteur d'une translation, il suffit de connaître un couple (point, image) ; choisir le centre $O_1$ est astucieux car son image par $r_1$ est immédiate.
\end{corrige}

\begin{corrige}[Exercice 6 — Homothétie et translation dans un repère]
\begin{enumerate}
\item Par l'homothétie $h$ de centre $O$ et de rapport $3$ : $M'(3x\,;3y)$. La translation $t$ de vecteur $\vec{u}(-6\,;3)$ ajoute ensuite les coordonnées de $\vec{u}$ :
\[ M''(3x-6\,;\,3y+3). \]
\item Posons $\Omega(a\,;b)$ et cherchons $a$ et $b$ tels que $\overrightarrow{\Omega M''}=3\,\overrightarrow{\Omega M}$ pour tout point $M$, c'est-à-dire :
\[ \begin{cases} x''-a=3(x-a) \\ y''-b=3(y-b) \end{cases} \iff \begin{cases} x''=3x-2a \\ y''=3y-2b \end{cases} \]
Par identification avec $x''=3x-6$ et $y''=3y+3$ :
\[ -2a=-6 \iff a=3 \qquad\text{et}\qquad -2b=3 \iff b=-\dfrac{3}{2}. \]
On a alors bien $\overrightarrow{\Omega M''}=3\,\overrightarrow{\Omega M}$ pour tout $M$ : $f=t\circ h$ est l'\cle{homothétie} de rapport $3$ et de centre $\Omega\left(3\,;-\dfrac{3}{2}\right)$.
\item Le centre d'une homothétie de rapport différent de $1$ est son unique point invariant. On résout $f(\Omega)=\Omega$ :
\[ \begin{cases} 3a-6=a \\ 3b+3=b \end{cases} \iff \begin{cases} 2a=6 \\ 2b=-3 \end{cases} \iff \begin{cases} a=3 \\ b=-\dfrac{3}{2} \end{cases} \]
Donc $\Omega\left(3\,;-\dfrac{3}{2}\right)$. Vérification : $f(\Omega)=\left(3\times 3-6\,;\,3\times\left(-\tfrac{3}{2}\right)+3\right)=\left(3\,;\,-\tfrac{3}{2}\right)=\Omega$. \checkmark
\end{enumerate}
\textbf{Point de vigilance :} le centre de $f$ n'est \emph{pas} $O$ : la translation a « déplacé » le centre de l'homothétie.
\end{corrige}

\begin{corrige}[Exercice 7 — Problème inverse]
\begin{enumerate}
\item Les rapports de deux homothéties de même centre se multiplient. En notant $k_2$ le rapport de $h_2$ :
\[ 2\times k_2=-6 \iff k_2=\dfrac{-6}{2}=-3. \]
\item $h_2$ est une homothétie de rapport $-3$ :
\begin{itemize}
\item les distances sont multipliées par $|-3|=3$ : c'est un \cle{agrandissement} de facteur $3$ ;
\item le rapport étant \emph{négatif}, l'image de chaque point se trouve de l'\emph{autre côté} du centre $O$ : la figure est retournée. L'effet est équivalent à un agrandissement de rapport $3$ suivi d'une symétrie centrale de centre $O$.
\end{itemize}
\end{enumerate}
\end{corrige}

\begin{corrige}[Exercice 8 — Plan d'une concession]
\begin{enumerate}
\item La composée de deux homothéties de même centre est une homothétie dont le rapport est le produit des rapports :
\[ k=\dfrac{5}{4}\times\dfrac{6}{5}=\dfrac{5\times 6}{4\times 5}=\dfrac{6}{4}=\dfrac{3}{2}. \]
L'agrandissement global est donc l'homothétie de rapport $\dfrac{3}{2}$.
\item Par une homothétie de rapport $k$, les aires sont multipliées par $k^2$. Ici :
\[ k^2=\left(\dfrac{3}{2}\right)^2=\dfrac{9}{4}. \]
L'aire sur le plan final est donc :
\[ \mathcal{A}'=800\times\dfrac{9}{4}=\dfrac{7200}{4}=1800\ \mathrm{cm}^2. \]
\end{enumerate}
\textbf{Point de vigilance :} ne pas multiplier l'aire par $k=\dfrac{3}{2}$ : c'est bien $k^2$ qui intervient pour les aires.
\end{corrige}

\begin{corrige}[Exercice 9 — Image d'un carré — type Probatoire]
\begin{enumerate}
\item Figure : tracer le carré $ABCD$ de centre $O$ et de côté $4\ \mathrm{cm}$, ses diagonales $[AC]$ et $[BD]$, puis placer les images demandées aux questions suivantes.
\item Les diagonales d'un carré sont perpendiculaires, de même longueur, et se coupent en leur milieu $O$. On a donc $OA=OB=OC=OD$ et
\[ \left(\overrightarrow{OA},\overrightarrow{OB}\right)=\left(\overrightarrow{OB},\overrightarrow{OC}\right)=\left(\overrightarrow{OC},\overrightarrow{OD}\right)=\left(\overrightarrow{OD},\overrightarrow{OA}\right)=90^{\circ}. \]
La rotation $r$ de centre $O$ et d'angle $90^{\circ}$ envoie donc chaque sommet sur le suivant :
\[ r(A)=B,\qquad r(B)=C,\qquad r(C)=D,\qquad r(D)=A. \]
\item On applique ensuite $h$, homothétie de centre $A$ et de rapport $2$, aux images par $r$ :
\[ A'=h(B),\qquad B'=h(C),\qquad C'=h(D),\qquad D'=h(A). \]
Construction : pour $A'$, on trace la demi-droite $[AB)$ et on reporte $\overrightarrow{AA'}=2\overrightarrow{AB}$ ; de même pour $B'$ et $C'$. Pour $D'$ : $\overrightarrow{AD'}=2\overrightarrow{AA}=\vec{0}$, donc $D'=A$ (le centre d'une homothétie est invariant).
\item L'image d'un carré par une homothétie est un \cle{carré} : l'homothétie conserve les angles et multiplie les longueurs par $|k|$. Le côté de $A'B'C'D'$ est donc $2\times 4=8\ \mathrm{cm}$, et son aire :
\[ \mathcal{A}'=8^2=64\ \mathrm{cm}^2. \]
Vérification par le rapport : $\mathcal{A}'=k^2\times\mathcal{A}=2^2\times 16=64\ \mathrm{cm}^2$. \checkmark
\item $h\circ r$ n'est ni une translation, ni une rotation : ces deux transformations conservent les longueurs, or ici les longueurs sont doublées. Ce n'est pas non plus une homothétie : une homothétie conserve les directions (angle nul), or $h\circ r$ fait tourner les figures de $90^{\circ}$. C'est une \cle{similitude directe} de rapport $2$ et d'angle $90^{\circ}$ (composée d'une rotation et d'une homothétie de centres distincts).
\end{enumerate}
\textbf{Barème indicatif (sur 10) :} figure soignée : 2 pts ; question 2 (justification par les diagonales) : 2 pts ; construction des images : 2 pts ; nature et aire de $A'B'C'D'$ : 2 pts ; question 5 argumentée : 2 pts.

\textbf{Points de vigilance :} à la question 2, il faut \emph{justifier} par les propriétés des diagonales du carré (conservation des distances et angle de $90^{\circ}$), pas seulement lire sur la figure ; à la question 5, citer la non-conservation des longueurs pour éliminer translation/rotation, et la rotation non nulle pour éliminer l'homothétie.
\end{corrige}

\begin{corrige}[Exercice 10 — Atelier de mécanique — type Probatoire]
\begin{enumerate}
\item $f=h_2\circ h_1$ est la composée de deux homothéties de même centre $O$ : c'est l'homothétie de centre $O$ et de rapport
\[ k=\dfrac{1}{2}\times(-4)=-\dfrac{4}{2}=-2. \]
\item Pour une homothétie de centre $O$ et de rapport $k$, l'image de $M(x\,;y)$ est $(kx\,;ky)$. Avec $k=-2$ et $P(6\,;-2)$ :
\[ P'\big(-2\times 6\,;\,-2\times(-2)\big)=(-12\,;\,4). \]
\item Une homothétie de rapport $k$ multiplie les longueurs par $|k|$. Ici $|k|=|-2|=2$, donc la longueur de l'image de l'arête est :
\[ 5\times 2=10\ \mathrm{cm}. \]
\item Les aires sont multipliées par $k^2=(-2)^2=4$ :
\[ \mathcal{A}'=12\times 4=48\ \mathrm{cm}^2. \]
\item Oui, le chef d'atelier a raison. En effet, les rapports se multiplient :
\[ h(O,-2)=h(O,2)\circ h(O,-1) \quad\text{car}\quad 2\times(-1)=-2. \]
Or $h(O,-1)$ est la \cle{symétrie centrale} de centre $O$, et $h(O,2)$ est un agrandissement de rapport $2$. Effectuer $f$ revient donc à effectuer une symétrie centrale suivie d'un agrandissement de rapport $2$ (l'ordre n'importe pas car même centre).
\end{enumerate}
\textbf{Barème indicatif (sur 10) :} question 1 : 2 pts ; question 2 : 2 pts ; question 3 : 2 pts ; question 4 : 2 pts ; question 5 (décomposition justifiée) : 2 pts.

\textbf{Points de vigilance :} pour les longueurs, citer $|k|$ et non $k$ (le rapport $-2$ est négatif) ; pour les aires, utiliser $k^2$ ; à la question 5, la justification repose sur la multiplication des rapports et sur l'identification de $h(O,-1)$ comme symétrie centrale.
\end{corrige}

\begin{corrige}[Exercice 11 — Chantier et déplacement d'une structure — type Probatoire]
\begin{enumerate}
\item On applique d'abord $h$ (on multiplie les coordonnées par $2$), puis $t$ (on ajoute $(-4\,;2)$) :
\begin{align*}
A(1\,;1) &\xrightarrow{h} (2\,;2) \xrightarrow{t} A'(-2\,;\,4)\\
B(4\,;1) &\xrightarrow{h} (8\,;2) \xrightarrow{t} B'(4\,;\,4)\\
C(1\,;3) &\xrightarrow{h} (2\,;6) \xrightarrow{t} C'(-2\,;\,8)
\end{align*}
\item Pour $M(x\,;y)$ : $h(M)=(2x\,;2y)$, puis
\[ f(M)=t\big(h(M)\big)=(2x-4\,;\,2y+2). \]
Donc $M''(2x-4\,;\,2y+2)$.
\item Le centre d'une homothétie de rapport différent de $1$ est son unique point invariant. On résout $f(\Omega)=\Omega$ avec $\Omega(a\,;b)$ :
\[ \begin{cases} 2a-4=a \\ 2b+2=b \end{cases} \iff \begin{cases} a=4 \\ b=-2 \end{cases} \]
Donc $\Omega(4\,;-2)$. Vérifions la relation vectorielle pour tout $M(x\,;y)$ :
\[ \overrightarrow{\Omega M''}\big(2x-4-4\,;\,2y+2-(-2)\big)=(2x-8\,;\,2y+4)=2(x-4\,;\,y+2)=2\,\overrightarrow{\Omega M}. \]
Donc $f$ est l'\cle{homothétie} de centre $\Omega(4\,;-2)$ et de rapport $2$.
\item On a $\overrightarrow{AB}(3\,;0)$ et $\overrightarrow{AC}(0\,;2)$ : ces vecteurs sont portés par les axes du repère orthonormé, donc le triangle $ABC$ est rectangle en $A$. Son aire est :
\[ \mathcal{A}=\dfrac{AB\times AC}{2}=\dfrac{3\times 2}{2}=3\ \text{unités d'aire}. \]
Par l'homothétie $f$ de rapport $2$, les aires sont multipliées par $2^2=4$ :
\[ \mathcal{A}'=3\times 4=12\ \text{unités d'aire}. \]
Vérification directe : $\overrightarrow{A'B'}(6\,;0)$ et $\overrightarrow{A'C'}(0\,;4)$, donc $A'B'C'$ est rectangle en $A'$ et
\[ \mathcal{A}'=\dfrac{6\times 4}{2}=12\ \text{unités d'aire}. \checkmark \]
\item Pour $M(x\,;y)$, on a $f(M)=(2x-4\,;\,2y+2)$, puis $t'$ ajoute $(4\,;-2)$ :
\[ t'\circ f(M)=\big(2x-4+4\,;\,2y+2-2\big)=(2x\,;\,2y). \]
On reconnaît l'expression analytique de l'homothétie de centre $O$ et de rapport $2$ :
\[ t'\circ f=h(O,2). \]
Interprétation : le vecteur $\vec{v}(4\,;-2)$ est l'opposé de $\vec{u}(-4\,;2)$, donc la translation $t'$ a exactement « compensé » la translation $t$, et il ne reste que l'homothétie initiale $h$.
\end{enumerate}
\textbf{Barème indicatif (sur 12) :} question 1 : 2 pts ; question 2 : 2 pts ; question 3 (point invariant + relation vectorielle) : 3 pts ; question 4 (aire + vérification) : 3 pts ; question 5 : 2 pts.

\textbf{Points de vigilance :} respecter l'ordre d'application ($f=t\circ h$ : d'abord $h$, ensuite $t$) ; à la question 3, le centre de $f$ n'est pas $O$ — il faut résoudre $f(\Omega)=\Omega$ ; à la question 4, justifier que le triangle est rectangle avant d'appliquer la formule de l'aire.
\end{corrige}

\newpage

\coursec{Fiche bilan}

\begin{popfigure}
\begin{tikzpicture}[mindmap, grow cyclic, every node/.style={concept, text=white, font=\small}, root concept/.append style={concept color=popink, font=\bfseries\small}, level 1/.append style={level distance=3.6cm, sibling angle=72}, level 1 concept/.append style={font=\scriptsize}]
\node[root concept]{Compos\'ees de transformations affines}
  child[concept color=popblue]{ node {G\'en\'eralit\'es : $g\circ f$, image de $M$} }
  child[concept color=poporange]{ node {Rotations m\^eme centre : $r(O,\alpha)\circ r(O,\beta)=r(O,\alpha+\beta)$} }
  child[concept color=popgreen]{ node {Rotations centres distincts : rotation ou translation} }
  child[concept color=poppurple]{ node {Homoth\'eties m\^eme centre : rapport $k\times k'$} }
  child[concept color=popgold]{ node {Homoth\'etie et translation : homoth\'etie ou translation} };
\end{tikzpicture}
\legende{Carte mentale : les compos\'ees de transformations affines du plan}
\end{popfigure}

\begin{aretenir}[L'essentiel de la le\c{c}on]
\textbf{D\'efinitions cl\'es}
\begin{itemize}
\item La \cle{compos\'ee} de $f$ suivie de $g$ est not\'ee $g\circ f$ : $M \xrightarrow{f} M' \xrightarrow{g} M''$, donc $(g\circ f)(M)=g\bigl(f(M)\bigr)$.
\item Attention \`a l'\cle{ordre} : dans $g\circ f$, c'est $f$ qui agit en premier.
\item Une transformation est une application bijective du plan dans lui-m\^eme ; la compos\'ee de deux transformations est une transformation.
\end{itemize}

\textbf{Th\'eor\`emes \`a conna\^itre par c\oe ur}
\begin{center}
\begin{tabularx}{\linewidth}{|l|X|}
\hline
\rowcolor{popblueL} \textbf{Compos\'ee} & \textbf{Nature et \'el\'ements caract\'eristiques} \\
\hline
$r(O,\alpha)\circ r(O,\beta)$ & Rotation de centre $O$, d'angle $\alpha+\beta$ (m\^eme sens). \\
\hline
$r(A,\alpha)\circ r(B,\beta)$, $A\neq B$, $\alpha+\beta\neq 0$ & Rotation d'angle $\alpha+\beta$ (centre \`a construire). \\
\hline
$r(A,\alpha)\circ r(B,-\alpha)$, $A\neq B$ & Translation (la somme des angles est nulle). \\
\hline
$h(O,k)\circ h(O,k')$ & Homoth\'etie de centre $O$, de rapport $k\,k'$. \\
\hline
$h(O,k)\circ h(O,k')$, $kk'=1$ & Identit\'e du plan (les deux homoth\'eties sont r\'eciproques l'une de l'autre). \\
\hline
$h(O,k)\circ h(O,k')$, $kk'=-1$ & Sym\'etrie centrale de centre $O$ (homoth\'etie de rapport $-1$). \\
\hline
$h(O,k)\circ t_{\vec{u}}$, $k\neq 1$ & Homoth\'etie de rapport $k$ (centre \`a d\'eterminer). \\
\hline
$h(O,k)\circ t_{\vec{u}}$, $k=1$ & Translation de vecteur $\vec{u}$ (car $h(O,1)$ est l'identit\'e). \\
\hline
\end{tabularx}
\end{center}

\textbf{M\'ethodes express}
\begin{itemize}
\item Pour identifier $g\circ f$ : chercher l'image de deux ou trois points bien choisis, puis comparer avec les transformations de r\'ef\'erence.
\item Pour deux rotations de m\^eme centre : additionner les angles, garder le centre.
\item Pour deux homoth\'eties de m\^eme centre : multiplier les rapports, garder le centre ; si $kk'=1$, c'est l'identit\'e ; si $kk'=-1$, c'est la sym\'etrie centrale de centre $O$.
\item Pour homoth\'etie + translation avec $k\neq 1$ : le r\'esultat est une homoth\'etie de rapport $k$ ; trouver le centre en cherchant le point invariant ($\Omega$ tel que $(g\circ f)(\Omega)=\Omega$).
\item Toujours conclure en r\'edigeant : \og la compos\'ee est une \ldots de \ldots \fg{} (nature + \'el\'ements caract\'eristiques).
\end{itemize}
\end{aretenir}

\begin{attention}[Erreurs fr\'equentes \`a \'eviter]
\begin{itemize}
\item Confondre l'ordre de composition : $g\circ f$ signifie \og $f$ puis $g$ \fg. En g\'en\'eral $g\circ f \neq f\circ g$.
\item Oublier la condition $\alpha+\beta \neq 0$ pour deux rotations de centres distincts : si $\alpha+\beta=0$, la compos\'ee n'est pas une rotation mais une translation.
\item Croire que $h(O,k)\circ h(O,k')$ avec $kk'=1$ donne une sym\'etrie centrale : c'est faux, on obtient l'\cle{identit\'e}. La sym\'etrie centrale correspond au cas $kk'=-1$.
\item Affirmer que la compos\'ee d'une homoth\'etie et d'une translation est toujours une homoth\'etie : si $k=1$, l'homoth\'etie est l'identit\'e et la compos\'ee est une translation.
\item Donner la nature de la transformation sans ses \'el\'ements caract\'eristiques (centre, angle, rapport, vecteur) : la r\'eponse est alors incompl\`ete au Probatoire.
\end{itemize}
\end{attention}

\begin{exoresolu}[Compos\'ee de deux homoth\'eties de m\^eme centre]
Soit $h_1=h(O,2)$ et $h_2=h(O,-3)$. D\'eterminer la nature et les \'el\'ements caract\'eristiques de $h_2\circ h_1$.
\tcblower
\textbf{Solution.} Les deux homoth\'eties ont le m\^eme centre $O$. La compos\'ee de deux homoth\'eties de m\^eme centre est l'homoth\'etie de m\^eme centre dont le rapport est le produit des rapports :
\[ k = 2\times(-3) = -6. \]
Comme $kk'=-6\neq 1$ et $-6\neq -1$, la compos\'ee $h_2\circ h_1$ est l'\cle{homoth\'etie} de centre $O$ et de rapport $-6$ : $h_2\circ h_1 = h(O,-6)$.
\end{exoresolu}

\begin{exoresolu}[Compos\'ee d'une homoth\'etie et d'une translation]
Soit $h=h(O,2)$ et $t=t_{\vec{u}}$ avec $\vec{u}\neq\vec{0}$. Quelle est la nature de $h\circ t_{\vec{u}}$ ?
\tcblower
\textbf{Solution.} Le rapport de l'homoth\'etie est $k=2\neq 1$. D'apr\`es le th\'eor\`eme, la compos\'ee d'une homoth\'etie de rapport $k\neq 1$ et d'une translation est une \cle{homoth\'etie} de rapport $k$. Donc $h\circ t_{\vec{u}}$ est une homoth\'etie de rapport $2$, dont le centre $\Omega$ est l'unique point invariant, d\'efini par $(h\circ t_{\vec{u}})(\Omega)=\Omega$.
\end{exoresolu}

\begin{aretenir}[Checklist avant le Probatoire]
\begin{itemize}
\item[$\square$] Je sais que dans $g\circ f$, la transformation $f$ agit en premier.
\item[$\square$] Je sais composer deux rotations de m\^eme centre : j'additionne les angles.
\item[$\square$] Je sais que deux rotations de centres distincts donnent une rotation d'angle $\alpha+\beta$, ou une translation si $\alpha+\beta=0$.
\item[$\square$] Je sais composer deux homoth\'eties de m\^eme centre : je multiplie les rapports.
\item[$\square$] Je reconna\^is le cas $kk'=1$ : la compos\'ee est l'identit\'e ; et le cas $kk'=-1$ : c'est la sym\'etrie centrale de centre $O$.
\item[$\square$] Je sais que la compos\'ee d'une homoth\'etie de rapport $k\neq 1$ et d'une translation est une homoth\'etie de rapport $k$.
\item[$\square$] Je sais d\'eterminer le centre de l'homoth\'etie r\'esultante en cherchant un point invariant.
\item[$\square$] Je v\'erifie toujours les conditions d'application (somme des angles nulle ? rapport \'egal \`a $1$ ?).
\item[$\square$] Je r\'edige ma conclusion en donnant la nature ET les \'el\'ements caract\'eristiques de la transformation.
\item[$\square$] Je sais illustrer la r\'eponse par une figure soign\'ee (centre, angle, rapport, vecteur).
\end{itemize}
\end{aretenir}

\findecours

\end{document}
